\documentclass[10pt]{ctexart}
\usepackage[a4paper,top=2.3cm,bottom=2.3cm,left=1.8cm,right=1.8cm,headheight=15pt,headsep=0.45cm,footskip=1.0cm]{geometry}
\usepackage{fvextra}
\usepackage[table]{xcolor}
\usepackage{colortbl}
\usepackage{fontspec}
\usepackage[most]{tcolorbox}
\usepackage{booktabs}
\usepackage{enumitem}
\usepackage{titlesec}
\usepackage{array}
\usepackage{tikz}
\usetikzlibrary{arrows.meta,positioning,calc}
\usepackage{fancyhdr}
\usepackage{lastpage}
\usepackage{hyperref}
\setmonofont{Consolas}
\setCJKmonofont{SimHei}
\setlength{\parindent}{0pt}
\setlist{nosep,leftmargin=1.4em}
\linespread{1.06}
% ---------- 配色 ----------
\definecolor{deep}{RGB}{24,43,73}
\definecolor{accent}{RGB}{31,95,168}
\definecolor{rowtint}{RGB}{232,239,247}
\definecolor{codebg}{RGB}{247,249,252}
\definecolor{thinline}{RGB}{178,193,214}
\definecolor{codeframe}{RGB}{203,215,230}
\definecolor{codenum}{RGB}{140,155,175}
\definecolor{moteal}{RGB}{23,124,122}
\definecolor{ggray}{RGB}{118,126,138}
% ---------- 页眉页脚 ----------
\pagestyle{fancy}
\fancyhf{}
\fancyhead[L]{\small\color{deep}\bfseries fpga-ptz-tracker\ \ 代码与设计说明}
\fancyhead[R]{\small\color{accent}\nouppercase{\rightmark}}
\fancyfoot[C]{\small\color{black!60}\thepage\ / \pageref{LastPage}}
\renewcommand{\headrulewidth}{0.7pt}
\renewcommand{\headrule}{{\color{accent}\hrule width\headwidth height\headrulewidth}}
\renewcommand{\footrulewidth}{0pt}
\fancypagestyle{plain}{%
  \fancyhf{}%
  \fancyfoot[C]{\small\color{black!60}\thepage\ / \pageref{LastPage}}%
  \renewcommand{\headrulewidth}{0pt}}
% ---------- 标题样式 ----------
\titleformat{\section}{\large\bfseries\color{deep}}%
  {\colorbox{deep}{\textcolor{white}{\,\thesection\,}}}{0.55em}{}%
  [{\vspace{2pt}\color{accent}\titlerule[1pt]}]
\titlespacing{\section}{0pt}{16pt}{9pt}
\titleformat{\subsection}{\normalsize\bfseries\color{deep}}%
  {\textcolor{accent}{\thesubsection}}{0.5em}{}
\titlespacing{\subsection}{0pt}{12pt}{5pt}
% ---------- 代码块 ----------
\fvset{framerule=0.5pt,framesep=2.5mm,rulecolor=\color{codeframe}}
\renewcommand{\theFancyVerbLine}{\textcolor{codenum}{\scriptsize\arabic{FancyVerbLine}}}
\newcommand{\code}[1]{\texttt{#1}}
\newcommand{\fileintro}[1]{{\small\color{accent}\sloppy\emergencystretch=2.5em #1\par}\vspace{2pt}}
\newtcolorbox{notebox}[1]{enhanced, breakable, colback=codebg, colframe=thinline,
  boxrule=0.5pt, arc=1.5mm, left=3mm, right=3mm, top=2mm, bottom=2mm,
  title={\bfseries #1}, fonttitle=\small\bfseries, coltitle=white,
  colbacktitle=deep}
% ---------- 超链接 ----------
\hypersetup{colorlinks=true, linkcolor=deep, urlcolor=accent, citecolor=accent,
  linktocpage=true,
  pdftitle={fpga-ptz-tracker 代码与设计说明（外部评审版）},
  pdfauthor={fpga-ptz-tracker 项目组}}
\setcounter{tocdepth}{1}
\begin{document}

\begin{titlepage}
\begin{tcolorbox}[enhanced, colback=deep, colframe=deep, sharp corners,
  boxrule=0pt, left=9mm, right=9mm, top=13mm, bottom=13mm,
  borderline west={1.4mm}{0pt}{accent}]
\raggedright
{\color{white}\Huge\bfseries fpga-ptz-tracker}\\[10pt]
{\color{rowtint}\Large\bfseries 代码与设计说明（外部评审版）}\\[12pt]
{\color{thinline}\large 2026 FPGA 创新设计赛道（安路）赛题二 · 视觉跟踪系统}
\end{tcolorbox}
\vspace{1.4cm}
\begin{center}
{\large PH1P35MDG324 / TD 6.2.1 / ModelSim SE}\\[10pt]
{\color{accent}\rule{0.32\textwidth}{0.9pt}}
\end{center}
\vspace{1.1cm}
\begin{notebox}{本文档用途}
\centering
{\small 汇总全部自研代码（RTL / 测试平台 / 仿真设施 / Python 工具，\textbf{约 7600 行}）\\
与架构说明、验证状态、开放问题清单，供外部 AI / 同行独立评审。\\
\textbf{版本 v8}（2026-10-03）：状态总览 / 已闭环项与证据指针 / 开放 Gate / 复现速查前置（见 \S1 首页）。\\
代码均通过 18 项位精确仿真回归，TD gate 综合 0 违规；开放项与上板核实清单见第 8 节。}
\end{notebox}
\vfill
{\small\color{black!55} 生成脚本：tools/make\_review\_pdf.py（可复现） · 2026-10}
\end{titlepage}
\tableofcontents
\newpage


\section{项目总览与阅读指南}

\textbf{【本版状态总览（v8）】}本版把“当前状态 / 已闭环项 / 开放 Gate / 证据索引 / 复现方式”
前置到本节开头，便于评审快速锁定基线。\textbf{已闭环项请当作待复核命题而非免检项}——
能给出可复现失效路径的反驳最有价值（本项目历史上两轮真 bug 均由此类反驳钉出）。

\textbf{验证与综合现状}
\begin{center}
{\small
\begin{tabular}{p{4.6cm}p{5.6cm}p{6.6cm}}
\toprule
\rowcolor{rowtint}\textbf{指标} & \textbf{数值} & \textbf{证据} \\
\midrule
位精确回归 & \textbf{18/18 PASS} & \code{测试运行\_20261003\_202418.log} \\
行突发全链验收 & 2764800 px / sof=3 / eol=2160 / drop=0 & \code{tb\_isp\_burst}（行节奏馈送）\\
FIFO 最大占用 & 1601/2048 px（78\%）& \code{tb\_isp\_burst} 观测 \\
TD 综合（像素域 74.25MHz）& \textbf{WNS +2.441ns}，0 违规端点 & \code{02\_td\_synth\_gate/reports/} \\
TD 综合（控制域 50MHz）& WNS +8.747ns & 同上 \\
资源（视觉链）& 3179 lut / 1963 reg / dram 256 / \textbf{eram 62/108} / dsp 8 & 同上 \\
\bottomrule
\end{tabular}}
\end{center}

\textbf{已闭环项（处置 + 证据指针）}
\begin{center}
{\small
\begin{tabular}{p{4.2cm}p{5.4cm}p{7.2cm}}
\toprule
\rowcolor{rowtint}\textbf{问题} & \textbf{处置} & \textbf{证据} \\
\midrule
算法 mux 提交沿丢 SOF & 提交沿组合选择 + 超时回退 & \code{tb\_algo\_switch}（SOF$\to$SOF 严格 scoreboard）\\
128bit 重排 tuser 错位 & 抽头改 AWB 96bit 直抽 & \code{tb\_isp\_burst} + system\_top \\
system\_top 悬空网 & 握手修正（I\_tvalid 直用 RAW / O\_tready 显式）& 代码 + 迭代记录 \#N+6 \\
串化器容量不足（128px）& v4 弹性缓冲 2048px + 寄存器读 & \code{tb\_isp\_burst} 验收 \\
resync 抱死（tuser 缺失）& v4.1 看门狗 + tlast 门控恢复 & \code{tb\_ser96to24} 阶段 2 故障注入 \\
跑马灯 ROM mux 链超时 & 前瞻 4 像素 + 4 级移位流水 & TD +2.441ns + \code{tb\_osd\_track} 位精确 \\
LUTRAM 映射未证 & TD 实证（bin\_smooth 256 dram lut）& \code{02\_td\_synth\_gate/reports\_prefix/} \\
motion ERAM 同址依赖 & 写地址延迟一拍结构解耦 & TD（eram 57）+ \code{tb\_motion\_detect} \\
OSD 影子换装撕裂 & 结构成立（SOF 换装）& \code{tb\_osd\_track} / \code{tb\_vision\_top} \\
闭环控制滞后 & T1/T4 时序修正（12 帧锁定）& \code{tb\_closed\_loop} \\
\bottomrule
\end{tabular}}
\end{center}

\textbf{开放 Gate（评审请聚焦）}
\begin{center}
{\small
\begin{tabular}{p{1.1cm}p{7.4cm}p{8.2cm}}
\toprule
\rowcolor{rowtint}\textbf{Gate} & \textbf{内容} & \textbf{验证途径} \\
\midrule
A & 真机 ISP 最大连续 burst / 多行合并行数上限 & 上板：O\_drop\_cnt / O\_occ\_max / sync\_err + 帧计数 \\
B & P\&R 时序与最终资源（gate 级为估算，净延迟占 \~{}68\%）& \code{opt\_place/opt\_route} 后复核 \\
C & 算法切换短帧对真实视频消费者的兼容性 & 先测两分支延迟差是否固定，再决定是否对齐 \\
D & ERAM 最终布局（视觉 62 + 官方链 39 $\Rightarrow$ \~{}94\%）& P\&R 布局报告 \\
E & 边缘支路未在合成激励下保留；cdc busy 断言待补 & 正式 P\&R 激励覆盖 / 断言 TODO \\
\bottomrule
\end{tabular}}
\end{center}

\textbf{证据索引（可向维护者索取对应工件）}
\begin{itemize}
  \item 回归日志：\code{03\_vision/step4\_osd\_integration/04\_迭代文件/测试运行\_*.log}
  \item TD 综合证据：\code{08\_board\_bringup/02\_td\_synth\_gate/}（reports/ 终态；reports\_prefix/ 修复前）
  \item 板级突发测量/验收 TB：\code{08\_board\_bringup/01\_video\_passthrough/sim/tb\_isp\_burst.v}
  \item 全量迭代记录（\#N+6 \~{} \#N+9 为关键节点）：\code{04\_embedded/04\_迭代文件/迭代记录\_嵌入式对接.md}
\end{itemize}

\textbf{复现速查}
\begin{itemize}
  \item 全量回归：\code{cd tools \&\& python run\_all.py}（约 4 分钟，18 项）
  \item 行突发验收：\code{vsim -c -gFEED\_MODE=1 tb\_isp\_burst}（板级 sim 目录，用法见其文件头注释）
  \item TD gate 综合：\code{td\_commands\_prompt.exe <TD>/doc/scripts/DefaultFlow.tcl}
        （run 目录放 settings.cfg 与 vision\_gate.sdc；完整说明见 \code{02\_td\_synth\_gate/README.md}）
\end{itemize}

\textbf{卷宗版本历史}
\begin{center}
{\small
\begin{tabular}{p{1.8cm}p{3.0cm}p{11.6cm}}
\toprule
\rowcolor{rowtint}\textbf{版本} & \textbf{口径} & \textbf{关键变化} \\
\midrule
v1 & 旧排版 105 页 & 初版（全代码 + 架构说明）\\
v2\~{}v4 & --- & 第 1\~{}6 轮外部评审落地（18 项回归成型）\\
v5 & 旧排版 117 页 & 架构级复审实测：ISP 行突发/全灭证据 + TD 两 Gate 解除\\
v6 & 旧排版 119 页 & 严格帧 scoreboard（SOF$\to$SOF）+ FIFO 占用实测 \\
v7 & 旧排版 122 页 & v4.1 异常恢复 + 跑马灯结构性流水（WNS +2.441ns）\\
\textbf{v8（本版）} & \textbf{精排版 + 内容同步} & 页眉/页脚/彩色标题/架构图；状态总览、证据索引、复现速查前置 \\
\bottomrule
\end{tabular}}
\end{center}
{\small\color{black!55} 注：v8 起排版升级（每页更密），页数与旧版不可直接比较；内容以本版为准。}

\subsection{一句话与架构}

\textbf{2026 FPGA 创新设计赛道（安路）赛题二}：基于 PH1P35 的实时视觉跟踪系统。
摄像头 SC520CS（1280x720@60）经官方 ISP 链路输出 96bit AXI4-Stream（4 像素/拍），
自研检测链在纯旁路窃流点接入，完成 颜色识别 → 阈值分割 → 3x3 去噪 → 质心/包围框，
结果经跨时钟域送至云台控制（双路 50Hz 舵机 PWM）与 OSD 叠加/遥测。

\begin{center}
\begin{tikzpicture}[x=1cm,y=1cm,>={Stealth[length=2.4mm]},
  main/.style={draw=deep,fill=rowtint,rounded corners=2pt,align=center,font=\small,inner sep=3pt,minimum height=8mm},
  mod/.style={draw=accent,fill=codebg,rounded corners=2pt,align=center,font=\footnotesize\ttfamily,inner sep=2.5pt,minimum height=7mm},
  io/.style={draw=deep,fill=white,rounded corners=2pt,align=center,font=\footnotesize,inner sep=3pt,minimum height=7mm},
  chip/.style={draw=thinline,fill=codebg,rounded corners=6pt,align=center,font=\scriptsize,inner sep=3pt,text width=3.5cm,minimum height=9mm},
  lab/.style={font=\tiny\itshape,text=black!60,align=center},
  arr/.style={->,draw=accent!75!black,semithick},
  elbow/.style={arr,rounded corners=3pt}]
% ---- 主通路（显示链路，零扰动）----
\node[main] (sensor) at (1.1,0) {SC520CS\\[-2pt]{\scriptsize 1280x720@60}};
\node[main] (isp) at (4.7,0) {官方 ISP\\[-2pt]{\scriptsize AWB 后 96bit}};
\node[io] (thru) at (11.1,0) {直通出 · 显示链路（零扰动）};
\coordinate (tap) at (8.2,0);
\draw[arr] (sensor.east) -- (isp.west);
\draw[arr] (isp.east) -- (tap);
\draw[arr] (tap) -- (thru.west);
\fill[accent] (tap) circle (1.5mm);
\node[lab,above=2pt] at (tap) {旁路窃流点};
% ---- 检测分支（旁路窃流）----
\node[mod] (ser) at (1.6,-1.8) {ser96to24};
\node[mod] (rgb) at (4.1,-1.8) {rgb2hsv};
\node[mod] (hsv) at (7.0,-1.8) {hsv\_threshold};
\node[mod] (sm3) at (10.1,-1.8) {bin\_smooth3x3};
\node[mod] (cen) at (13.4,-1.8) {centroid\_accum};
\node[lab] at (3.9,-1.05) {检测分支（旁路窃流）};
\draw[elbow] (tap) -- (8.2,-0.55) -- (1.6,-0.55) -- (ser.north);
\draw[arr] (ser.east) -- (rgb.west);
\draw[arr] (rgb.east) -- (hsv.west);
\draw[arr] (hsv.east) -- (sm3.west);
\draw[arr] (sm3.east) -- (cen.west);
% ---- 输出：OSD 与云台 ----
\node[io] (osd) at (13.4,-3.5) {OSD 叠加};
\node[mod] (gim) at (10.0,-3.5) {gimbal\_ctrl};
\node[mod] (pwm) at (6.7,-3.5) {pwm\_servo};
\draw[elbow] (cen.south) -- (13.4,-2.7) -- (10.0,-2.7) -- (gim.north);
\draw[arr] (13.4,-2.7) -- (osd.north);
\node[lab,anchor=west] at (13.75,-2.95) {135bit 目标包 · cdc};
\node[lab] at (11.7,-2.45) {23bit 控制包 · cdc};
\draw[arr] (gim.west) -- (pwm.east);
\node[lab] at (8.35,-3.25) {50Hz};
\node[lab] at (6.7,-4.2) {双路 50Hz 舵机};
% ---- 边缘分支（常开）----
\node[mod] (edge) at (2.4,-5.1) {edge\_detect\_96};
\node[lab] at (2.4,-5.65) {边缘分支（常开，Lab2 复用）};
\node[io] (mux) at (6.4,-5.1) {帧首 mux 切换显示};
\draw[arr] (edge.east) -- (mux.west);
\draw[elbow,dashed] (mux.north) -- (6.4,-4.35) -- (16.7,-4.35) -- (16.7,-0.02) -- (thru.east);
% ---- 辅助模块 ----
\node[chip] at (2.1,-6.35) {video\_adj\\亮度/对比度（画质，显示路径）};
\node[chip] at (6.2,-6.35) {motion\_detect\\帧差（运动，115KB ERAM）};
\node[chip] at (10.3,-6.35) {led\_buzzer\\声光提示};
\node[chip] at (14.4,-6.35) {uart\_telemetry\\115200，每帧 9 字节（遥测）};
\node[lab] at (8.5,-7.15) {实线：数据流；虚线：显示切换；圆点：窃流接入点};
\end{tikzpicture}
\end{center}

\subsection{平台、工具链与关键约定}

\begin{itemize}
  \item 器件/工具：安路 PH1P35MDG324（-3）；TD 6.2.1（build 168116）；仿真 ModelSim SE 10.4；Verilog-2001。
  \item 时钟域：像素域 74.25MHz（1 像素/拍，96bit 流为 1 beat/4 clk）；PWM/控制域 50MHz；跨域仅经 pulse\_sync（单 bit）与 cdc\_4phase（多 bit）。
  \item 96bit 契约：4 像素 RGB888/拍；tuser=帧首、tlast=行尾、tready 恒 1（不反压）。
  \item 舵机：50Hz，0.5～2.5ms 对应 0～180 度；共周期计数器；上电中位 1500us；双道限幅。
  \item UART 遥测：115200-8N1；帧格式 A5|valid|cx(2B)|cy(2B)|preset|累加和|5A。
\end{itemize}

\subsection{目标包位图（IF-VISION v1.1.0，135bit，位偏移唯一权威定义）}

\begin{center}
\begin{tikzpicture}[x=1cm,y=1cm,
  seg/.style={rounded corners=1.2pt},
  nam/.style={font=\tiny\bfseries,text=white},
  rng/.style={font=\tiny,text=white!75!deep},
  nar/.style={font=\tiny\bfseries,text=black!70},
  rngn/.style={font=\tiny,text=black!55}]
% ---- 第一行：bit 134 .. 66 ----
\path[fill=deep,seg] (0,-0.375) rectangle (0.38,0.375);
\path[fill=deep,seg] (0.46,-0.375) rectangle (0.84,0.375);
\path[fill=deep,seg] (0.92,-0.375) rectangle (1.58,0.375);
\path[fill=accent,seg] (1.66,-0.375) rectangle (4.08,0.375);
\path[fill=accent,seg] (4.16,-0.375) rectangle (6.58,0.375);
\path[fill=accent,seg] (6.66,-0.375) rectangle (9.08,0.375);
\path[fill=accent,seg] (9.16,-0.375) rectangle (11.58,0.375);
\path[fill=accent,seg] (11.66,-0.375) rectangle (13.86,0.375);
\node[nar,rotate=90,anchor=west] at (0.19,-0.45) {valid};
\node[rngn,rotate=-90,anchor=west] at (0.19,0.45) {[134]};
\node[nar,rotate=90,anchor=west] at (0.65,-0.45) {algo};
\node[rngn,rotate=-90,anchor=west] at (0.65,0.45) {[133]};
\node[nar,rotate=90,anchor=west] at (1.25,-0.45) {preset};
\node[rngn,rotate=-90,anchor=west] at (1.25,0.45) {[132:130]};
\node[nam] at (2.87,0.10) {cx};
\node[rng] at (2.87,-0.17) {[129:119]};
\node[nam] at (5.37,0.10) {cy};
\node[rng] at (5.37,-0.17) {[118:108]};
\node[nam] at (7.87,0.10) {minx};
\node[rng] at (7.87,-0.17) {[107:97]};
\node[nam] at (10.37,0.10) {maxx};
\node[rng] at (10.37,-0.17) {[96:86]};
\node[nam] at (12.76,0.10) {miny};
\node[rng] at (12.76,-0.17) {[85:76]};
\node[rngn,anchor=east] at (-0.12,0) {bit 134};
\node[rngn,anchor=west] at (13.98,0) {bit 66};
% ---- 第二行：bit 65 .. 0 ----
\path[fill=accent,seg] (0,-2.675) rectangle (1.76,-1.925);
\path[fill=moteal,seg] (1.84,-2.675) rectangle (2.22,-1.925);
\path[fill=moteal,seg] (2.30,-2.675) rectangle (5.38,-1.925);
\path[fill=moteal,seg] (5.46,-2.675) rectangle (7.88,-1.925);
\path[fill=moteal,seg] (7.96,-2.675) rectangle (10.38,-1.925);
\path[fill=moteal,seg] (10.46,-2.675) rectangle (12.66,-1.925);
\path[fill=moteal,seg] (12.74,-2.675) rectangle (14.94,-1.925);
\path[fill=ggray,seg] (15.02,-2.675) rectangle (15.40,-1.925);
\node[nam] at (0.88,-2.20) {fid};
\node[rng] at (0.88,-2.47) {[65:58]};
\node[nar,rotate=90,anchor=west] at (2.03,-2.75) {mot\_valid};
\node[rngn,rotate=-90,anchor=west] at (2.03,-1.85) {[57]};
\node[nam] at (3.84,-2.20) {mot\_cnt};
\node[rng] at (3.84,-2.47) {[56:43]};
\node[nam] at (6.67,-2.20) {mot\_minx};
\node[rng] at (6.67,-2.47) {[42:32]};
\node[nam] at (9.17,-2.20) {mot\_maxx};
\node[rng] at (9.17,-2.47) {[31:21]};
\node[nam] at (11.56,-2.20) {mot\_miny};
\node[rng] at (11.56,-2.47) {[20:11]};
\node[nam] at (13.84,-2.20) {mot\_maxy};
\node[rng] at (13.84,-2.47) {[10:1]};
\node[nar,rotate=90,anchor=west] at (15.21,-2.75) {reserved（填 0）};
\node[rngn,rotate=-90,anchor=west] at (15.21,-1.85) {[0]};
\node[rngn,anchor=east] at (-0.12,-2.3) {bit 65};
\node[rngn,anchor=west] at (15.52,-2.3) {bit 0};
% ---- 图例 ----
\node[rngn,anchor=west] at (0.3,-3.35)
  {\textcolor{deep}{\rule[2.2pt]{7pt}{7pt}}\ 头部（valid/algo/preset）\quad
   \textcolor{accent}{\rule[2.2pt]{7pt}{7pt}}\ 目标/坐标\quad
   \textcolor{moteal}{\rule[2.2pt]{7pt}{7pt}}\ 运动统计\quad
   \textcolor{ggray}{\rule[2.2pt]{7pt}{7pt}}\ 保留};
\end{tikzpicture}

{\tiny\color{black!55} 段宽示意（非等比例）；位段范围以各段标注为准，为唯一权威定义。}
\end{center}

\subsection{验证体系（18 项回归，2026-10-03 全绿）}

方法：\textbf{金标准先行}。所有算术先在 Python 金标准 \code{golden\_model.py} 中实现，
RTL 逐位复刻；激励/期望由 \code{gen\_vectors.py} 生成；\code{run\_all.py} 一键回归。

\begin{center}
\renewcommand{\arraystretch}{1.35}
\arrayrulecolor{thinline}
\rowcolors{2}{white}{rowtint}
{\small
\begin{tabular}{@{}>{\ttfamily}p{5.2cm}p{4.9cm}p{4.6cm}@{}}
\toprule
\rowcolor{deep}
\textcolor{white}{\textbf{TB}} & \textcolor{white}{\textbf{验证对象}} & \textcolor{white}{\textbf{粒度}} \\
\midrule
tb\_ser96to24 / tb\_tpg & 串化器 / 测试图案源 & 逐位 \\
tb\_pulse\_sync / tb\_cdc4phase & 两种跨域桥 & 行为/时序 \\
tb\_rgb2hsv / tb\_hsv\_threshold & 颜色变换 / 阈值 & 逐位（10010/4000 向量）\\
tb\_bin\_smooth / tb\_centroid & 去噪 / 质心 & 逐位（12288 像素/3 帧）\\
tb\_detect\_chain & 检测链整链 & 3 帧真实场景逐字段 \\
tb\_gimbal / tb\_pwm & 控制律 / PWM 波形 & 10 步进表 / 逐拍 \\
tb\_osd\_track / tb\_vision\_top & OSD / 整机 & 逐像素 921600 / 5 断言 \\
tb\_closed\_loop & 相机-舵机在环 & 12 帧锁定 \\
tb\_motion\_detect & 帧差运动 & 6 帧逐字段（含末像素边界）\\
tb\_uart\_telemetry & 遥测（含 cx/cy 全坐标字节） & 逐字节 \\
tb\_cdc\_reset\_skew & CDC 复位偏斜（第 17 项） & 4 轮错开释放 \\
tb\_algo\_switch & 算法切换帧对齐（第 18 项） & SOF→SOF 严格 scoreboard \\
\bottomrule
\end{tabular}}
\end{center}

\subsection{已完成的独立评审（供本轮评审参考，不必重复）}

自研代码已经过\textbf{六轮外部 AI 评审}（三版独立模型、累计 100+ 条意见，逐条实证甄别），
已修复并回归的问题包括：

\begin{itemize}
  \item \textbf{TD 综合阻断类}：多驱动网络/寄存器、层次化引用取跨域信号；
  \item \textbf{综合模式依赖类}：motion\_detect 同址读写依赖 ERAM 写模式
        （已改"写地址延迟一拍"结构性解耦）；组合读行缓存/缓冲的 LUTRAM 映射
        （列入第 8 节硬 Gate）；
  \item \textbf{系统级接线类}：vision\_top 显示链曾断接（video\_adj 输出无下游，
        edge/亮度对比度为死分支）——已重接 mux->video\_adj->OSD；
  \item \textbf{边界/时序类}：motion 帧末漏最后像素（锁存再延一拍，并新增
        末像素专项向量）；PWM 占空采样与 gimbal 更新同沿致多一帧滞后（已修）；
        ser96to24 满缓冲拒收 SOF 后帧同步丢失窗口（已改"仅 SOF 真正入队才解除"）；
  \item \textbf{恢复/复位类}：CDC 两侧复位无同步释放（vision\_top 已加 pix/pwm
        双域两拍同步释放网络）；
  \item \textbf{功能类}：跑马灯索引截断、亮度负值符号、D 字段滞后一包、
        蜂鸣时长与文档不符；
  \item \textbf{控制性能}：控制节拍改为"新包到达即节拍"，配合上述修复，
        相机-舵机在环锁定帧数 36 -> 12。
\end{itemize}

已确认"不成立"的被驳意见（供参考）：组合读被静默强转 BRAM、OSD 影子寄存器
同沿撕裂、大 LUT-ROM 必然时序违例（均给出过实证反驳，详见仓库迭代记录）。

第四轮落地批次（三版独立模型评审后已完成，供本轮参考）：（1）算法 mux 改为
\textbf{输出侧 SOF 提交}（每个输出帧严格单算法，消除切换帧内新旧拼接）；
（2）CDC \textbf{复位同源契约}写入 cdc\_4phase 头部，并新增
\code{tb\_cdc\_reset\_skew}（第 17 项回归，4 轮错开释放全部 PASS）；
（3）system\_top 的 OSD 流补全 \code{O\_osd\_sof/O\_osd\_eol} 正式端口，
异步输入（algo/bright/contrast）加 2FF 同步；（4）运动框与 M-字段的
\textbf{像素级覆盖}测试上线后立刻抓出并修复"OSD M 字段读影子旧值"缺陷
（与 D 字段同类，此前只修了 D）。另：本轮评审中出现的若干"接线 bug"指控
（UART cy、运动框坐标、video\_adj 复位极性）经逐行实证均为引用幻觉，
代码本身正确——文档与代码冲突时以代码为准。

第五轮落地批次（ChatGPT 第三版复审）：外部评审指出算法 mux 的"输出侧提交"存在
NBA 同沿首拍丢失——mux 选择块与提交寄存器分处两个 always，提交沿扇出的首个
帧首脉冲落到旧选择分支，新算法帧的 SOF 被吃掉。此项指控成立：新增
\code{tb\_algo\_switch}（第 18 项回归）首跑即复现（MISM win=5 sof=0，切换后
首窗零 SOF）。修复为\textbf{提交沿组合选择}（\code{algo\_sel} 提交当拍即生效）
并加入\textbf{请求超时回退}（输入帧首仍未提交则强制提交，兑现上一轮评审的
"绝不卡死"建议）；修复后专项 PASS，18 项回归全绿。

第六轮（架构级复审，2026-10-03 深夜）：以"不依赖文档自证"为前提，改用\textbf{实测与工具}
替代论证——新增 \code{tb\_isp\_burst}（三档馈送节拍）坐实 ISP 行级突发与 ser96to24
实链功能性全灭，并连带抓出 128bit 重排 tuser 结构性错位、system\_top isp\_tready
悬空网两个独立致命；用\textbf{安路 TD 无头综合}（read\_design→opt\_rtl→opt\_gate）解除
两个 Gate：LUTRAM 映射实证（ser96to24 56 / bin\_smooth3x3 256 dram lut）、
50MHz 域 WNS +8.747ns（gimbal 乘加无虞），并指认真正的时序风险在 OSD 跑马灯路径
（74.25MHz 域 WNS −0.113ns）。详见迭代记录 \#N+6 与第 8 节。

\textbf{修复轮（\#N+7，用户授权后实施）}：三处致命全部修复并复验——ser96to24 v4
（2048px 寄存器读弹性缓冲）、抽头改 AWB 96bit 直抽（isp\_top 新增 tap\_* 端口）、
system\_top 握手修复、跑马灯数组 ROM 化。验收：tb\_isp\_burst 行突发模式
\code{pixels=2764800/sof=3/eol=2160/drop=0} PASS（修复前 188 像素/0 SOF）；
18 项回归全绿；TD 像素域 \textbf{WNS −0.113ns → +0.541ns（0 违规端点）}。

\textbf{加固轮（\#N+9，第七轮复审响应）}：采纳其两条核心意见——\textbf{ser96to24 v4.1}
（帧同步看门狗 + tlast 门控受控恢复，杜绝 tuser 缺失死锁；故障注入 TB 验证；O\_occ\_max/
sync\_err 观测量经 status[6] 引出）；\textbf{跑马灯结构性流水化}（并诚实更正：\#N+7 的
"数组 ROM 化"并未真正消除 mux 树，WNS 曾回退 −1.025ns；前瞻 4 像素+移位对齐后
像素域 \textbf{WNS +2.441ns、0 违规端点}）。多行合并突发与短帧兼容性维持"真机确认"Gate。

第二波待办与上板硬 Gate 见第 8 节；\textbf{如你不同意以上任何结论，请明确指出并给出依据}。

\subsection{请重点评审的问题（回答时请注明严重度与依据）}

已闭环项（详见 §1.5 与 §8，请勿重复报告，除非有新证据）：算法 mux 首拍 SOF、128bit 重排
tuser 错位、system\_top 悬空网、ser96to24 突发（v4 已修）、LUTRAM 映射（TD 已证）、
50MHz 时序（+8.747ns）、OSD 跑马灯与文本影子、motion ERAM 解耦。当前开放边界如下：

\begin{enumerate}
  \item \textbf{ser96to24 v4.1 的对抗性审计}：容量 2048px 依据=实测行突发 320 beat、
        最大占用 1601/2048（78\%）、drop=0；寄存器读（+1 拍）；溢出帧级重同步 +
        \textbf{帧同步看门狗（SYNC\_TIMEOUT=2 帧预算）+ tlast 门控恢复}（tuser 永久缺失
        不再死锁，故障注入 TB 已验）。请继续找：多行合并突发（3 行=2880px > 2048）
        的分布边界是否可在板级可靠界定；恢复后的帧相位偏移对下游有无隐性影响；
        以及 DEPTH 是否需要更大的量化判断（代价：4 块 ERAM/实例，整机预算 ≈94\%）。
  \item \textbf{抽头点选择}：视觉从 AWB 96bit 直抽（isp\_top 只增 4 个 tap\_* 端口，
        未改 ISP 逻辑），弃用官方 128bit 重排路径（其 tuser 与 valid 结构性错位已实测）。
        请评估该改动对"官方显示链零扰动"承诺是否有任何破坏，以及 tap 端口加法
        在 TD 综合/P\&R 上可能引入的负担或副作用。
  \item \textbf{无反压直抽 ISP 的完整性边界}：窃流语义下 O\_tready 恒 1。
        实测：行突发节奏下零丢行；但满速馈送（非真实节奏）下曾观察到 AWB 仅输出
        \~{}668/720 行。请评估"无背压"的残余风险，并给出上板必须确认的观测项清单
        （O\_drop\_cnt / O\_occ\_max / sync\_err 已可观测）。
  \item \textbf{P\&R 前的时序风险预判}：gate 级 0 违规（pwm +8.747ns / pix \textbf{+2.441ns}，
        跑马灯已结构性流水化）。但 68\% 为估算净延迟。请指出 P\&R 后最可能恶化的路径簇
        与兜底改法（成本排序），特别复核：OSD 文本影子装载、cdc 目的侧数据到文本数组、
        AWB tap fanout。
  \item \textbf{严格帧 scoreboard 的语义边界}：SOF→SOF 判据（行宽严格 1280 / 非切换跨度
        恰 720 行 / 切换跨度允许旧帧截断 ≤2 次）已 PASS；实测两次切换截断形态为
        \code{0 行+1px} 与 \code{719 行+1279px}（源于两分支流水延迟差）。请评估该语义
        对下游（显示/推流/录像）的可接受性，以及"对齐两分支延迟实现零截断"是否值得做。
  \item \textbf{ERAM 预算与资源分配}：视觉链 61/108（motion 57 + ser 8），
        叠加官方 DDR/HDMI 链 39 块 ⇒ 真机 ≈93\%。请给出资源分配建议
        （降 ser DEPTH / 共享串化器 / 裁剪 motion 分辨率，按性价比排序）。
  \item \textbf{竞赛层面}：以"基础+扩展+进阶"评分口径看，此实现还有哪些短板可补？
  \item \textbf{接线语义审计（继续）}：上一轮同类审计抓到真实缺陷（system\_top 的
        isp\_tready 悬空网）。请继续沿 vision\_top / system\_top 每条数据通路核对
        "文档声明意图"与"实际连线"的闭环性，重点：tap\_* 四端口、cdc\_4phase 两侧
        握手、OSD 的 sof/eol/valid 对齐链。
\end{enumerate}

\section{视觉链路 RTL（数据流顺序）}
\subsection{ser96to24.v}
\fileintro{检测链入口适配层 v4.1（2026-10-03 弹性缓冲重写 + 异常恢复加固）：96bit(4px/拍) AXI-S 拆成1 像素/拍 vpix。官方 ISP AWB 输出实测为行级突发（320 beat=1280px/行），v3 的 128px 缓冲远不够（实链 188/2764800 像素）；v4 容量 2048 像素（512 beat）+ 寄存器读（输出 +1 拍，兼容 LUTRAM/ERAM）。无反压；溢出丢弃 + 帧级重同步；v4.1 增加帧同步看门狗（SYNC\_TIMEOUT 后置 O\_sync\_err）与tlast 门控受控恢复（tuser 永久缺失不再死锁；故障注入 TB 验证）+ O\_occ\_max 占用水位观测。验收：tb\_isp\_burst 行突发模式像素/帧计数/drop 全对；18 项回归全绿。}
\begin{Verbatim}[numbers=left,numbersep=5pt,frame=single,fontsize=\footnotesize,breaklines=true,breakanywhere=true]
// =============================================================================
// ser96to24.v —— 96bit(4像素/拍) AXI4-S -> 1 像素/拍 vpix 串化器（v4，弹性缓冲重写）
// -----------------------------------------------------------------------------
// 视觉检测链入口适配层。挂在官方 ISP 的 AWB 96bit 输出上（4 像素 RGB888 打包，
// tuser=帧首、tlast=行尾，两者均在有效拍上；见 system_top 的抽头与迭代记录 #N+6）。
//
// v4 重写缘由（2026-10-03 架构级复审后实测坐实）：
//   官方 ISP 的 AWB 输出为"行级突发"：每行 320 连续 beat（=1280px @4px/拍），
//   行间有消隐间隔（实测 2160 段=3 帧x720 行，长度全部精确 320，无合并、无丢行）。
//   v3 的 128 像素缓冲（合 32 beat）净缺 3px/拍、连续 42 拍即溢出；且满缓冲后
//   SOF 无法入队导致 resync 不可恢复——与 system_top 同接法的实测输出仅
//   188/2764800 像素、零 SOF/EOL（功能性全灭）。
// v4 选型（数据驱动）：
//   缓冲深度 = 2048 像素（512 beat）：覆盖实测最坏 320-beat 行突发的峰值缺口
//     （320x(4-1)=960px）两倍余量；平均输入率 <1px/拍 时长期不增长。
//   读出改"寄存器读"（输出相对输入 +1 拍）：同时兼容 LUTRAM / ERAM 两种映射，
//     消除"组合读被综合成同步 RAM 会整帧错行"的实现依赖风险（复审 Gate 2 根治）。
//
// 语义：
//   O_sof=帧首像素；O_eol=行尾像素；O_valid=数据有效（滞后 I_tvalid 1 拍）。
//   无反压（O_tready 恒 1，窃流语义）；溢出丢弃 + 帧级重同步（SOF 能入队才解除，
//   损害限定当前帧）；O_drop_cnt 计数（饱和）。
//
// v4.1 加固（2026-10-03 第 7 轮外部复审，[中] 异常恢复缺口）：
//   原语义下 resync 只由"真实 tuser 入队"解除——若 tuser 永久缺失（帧同步契约被破坏），
//   模块将进入永久丢弃态（O_valid 恒 0、drop 饱和）。v4.1 增加：
//   1) 帧同步看门狗：resync 期间计 valid beat，超过 SYNC_TIMEOUT（默认 2 帧预算）
//      仍未等到真实 tuser -> 置 O_sync_err（粘滞至下一次真实 tuser 入队）；
//   2) tlast 门控受控恢复：超时后，下一个 tlast 拍被当作"合成 SOF"入队并解除 resync
//      （tlast 是真实行边界标记，恢复后帧相位可能偏移若干行——对 SOF 复位的下游
//      无害；若 tlast 也缺失则保持丢弃等复位，不猜 SOF）；
//   3) 观测量：O_occ_max（缓冲最大占用水位，ChipWatcher/维修诊断用）。
// =============================================================================
module ser96to24 #(
    parameter [23:0] SYNC_TIMEOUT = 24'd460800   // resync 看门狗阈值（beat；默认 2 帧）
)(
    input  wire        I_clk,
    input  wire        I_rst_n,
    // 96bit AXI4-Stream 视频输入（仅有效像素拍）
    input  wire [95:0] I_tdata,
    input  wire        I_tvalid,
    input  wire        I_tuser,
    input  wire        I_tlast,
    output wire        O_tready,
    // 1 像素/拍 vpix 输出（相对输入滞后 1 拍）
    output reg         O_valid,
    output reg         O_sof,
    output reg         O_eol,
    output reg  [23:0] O_data,
    // 调试：因缓冲满被拒的 beat 计数（饱和）
    output reg  [15:0] O_drop_cnt,
    // 调试：帧同步错误（看门狗触发；下一次真实 tuser 入队时清除）
    output reg         O_sync_err,
    // 调试：缓冲最大占用水位（像素）
    output reg  [11:0] O_occ_max
);

    localparam DEPTH = 2048;   // 像素缓冲深度（=512 beat；实测行突发 320 beat 的 2 倍余量）
    localparam PW    = 11;     // 指针位宽（DEPTH = 2^PW）
    localparam CW    = 12;     // 计数位宽（0..DEPTH）

    assign O_tready = 1'b1;    // 窃流分支，绝不反压上游

    // 像素缓冲（数据+标志同址）
    (* ram_style = "distributed" *) reg [23:0] mem   [0:DEPTH-1];
    (* ram_style = "distributed" *) reg        sof_q [0:DEPTH-1];
    (* ram_style = "distributed" *) reg        eol_q [0:DEPTH-1];
    reg [PW-1:0] wp;           // 写指针
    reg [PW-1:0] rp;           // 读指针
    reg [CW-1:0] cnt;          // 缓冲像素数（0..DEPTH）

    wire [PW-1:0] wp_n  = wp + 11'd1;
    wire [PW-1:0] wp_n2 = wp + 11'd2;
    wire [PW-1:0] wp_n3 = wp + 11'd3;
    wire [PW-1:0] rp_n  = rp + 11'd1;

    wire [23:0] p0 = I_tdata[95:72];
    wire [23:0] p1 = I_tdata[71:48];
    wire [23:0] p2 = I_tdata[47:24];
    wire [23:0] p3 = I_tdata[23:0];

    wire can_wr = (cnt <= (DEPTH - 4));

    // 满载丢弃后帧级重同步（v3 语义）+ v4.1 看门狗/受控恢复
    reg  resync;
    reg  [23:0] sync_cnt;      // resync 期间的 valid beat 计数（看门狗计时）
    wire sync_timeout = (sync_cnt >= SYNC_TIMEOUT);
    // 可入队的"帧首"：真实 tuser，或（超时后）tlast 门控的合成帧首
    wire wr_as_sof = I_tuser | (resync & sync_timeout & I_tlast);

    always @(posedge I_clk or negedge I_rst_n) begin
        if (!I_rst_n)
            resync <= 1'b0;
        else if (I_tvalid && wr_as_sof && can_wr)
            resync <= 1'b0;                       // 仅当帧首拍真正能入队才解除
        else if (I_tvalid && !can_wr)
            resync <= 1'b1;
    end

    always @(posedge I_clk or negedge I_rst_n) begin
        if (!I_rst_n)
            sync_cnt <= 24'd0;
        else if (!resync)
            sync_cnt <= 24'd0;
        else if (I_tvalid && (sync_cnt < SYNC_TIMEOUT))
            sync_cnt <= sync_cnt + 24'd1;
    end

    always @(posedge I_clk or negedge I_rst_n) begin
        if (!I_rst_n)
            O_sync_err <= 1'b0;
        else if (I_tvalid && I_tuser && can_wr)
            O_sync_err <= 1'b0;                   // 真实帧首恢复 -> 清错误
        else if (resync && sync_timeout)
            O_sync_err <= 1'b1;                   // 看门狗触发（粘滞）
    end

    wire do_wr = I_tvalid & O_tready & can_wr & (~resync | wr_as_sof);
    wire do_rd = (cnt != {CW{1'b0}});

    always @(posedge I_clk or negedge I_rst_n) begin
        if (!I_rst_n) begin
            wp  <= {PW{1'b0}};
            rp  <= {PW{1'b0}};
            cnt <= {CW{1'b0}};
            O_valid <= 1'b0;
            O_sof   <= 1'b0;
            O_eol   <= 1'b0;
            O_data  <= 24'd0;
            O_drop_cnt <= 16'd0;
            O_occ_max  <= 12'd0;
        end else begin
            // 写口：4 像素入队（分辨率宽为 4 的倍数，末拍必满 4 像素）
            if (do_wr) begin
                mem[wp]    <= p0;
                mem[wp_n]  <= p1;
                mem[wp_n2] <= p2;
                mem[wp_n3] <= p3;
                sof_q[wp]    <= wr_as_sof;          // 真实或（超时后）合成帧首
                sof_q[wp_n]  <= 1'b0;
                sof_q[wp_n2] <= 1'b0;
                sof_q[wp_n3] <= 1'b0;
                eol_q[wp]    <= 1'b0;
                eol_q[wp_n]  <= 1'b0;
                eol_q[wp_n2] <= 1'b0;
                eol_q[wp_n3] <= I_tlast;
                wp <= wp + 11'd4;
            end
            // 读口：寄存器读（rp 当拍指向的像素在下一拍出现在输出寄存器上）
            O_data  <= mem[rp];
            O_sof   <= sof_q[rp];
            O_eol   <= eol_q[rp];
            O_valid <= do_rd;
            if (do_rd)
                rp <= rp_n;
            // 缓冲像素数 + 最大占用水位（诊断）
            case ({do_wr, do_rd})
                2'b10:   cnt <= cnt + 12'd4;
                2'b01:   cnt <= cnt - 12'd1;
                2'b11:   cnt <= cnt + 12'd3;
                default: cnt <= cnt;
            endcase
            if (cnt > O_occ_max)
                O_occ_max <= cnt;                   // 占用水位单调记录（复位清零）
            if (I_tvalid && !do_wr && (O_drop_cnt != 16'hFFFF))
                O_drop_cnt <= O_drop_cnt + 16'd1;   // 缓冲满导致的丢弃（饱和计数）
        end
    end

    // synthesis translate_off
    reg drop_dbg_seen;
    always @(posedge I_clk) begin
        if (!I_rst_n) drop_dbg_seen <= 1'b0;
        else if ((cnt > (DEPTH - 4)) && I_tvalid && !drop_dbg_seen) begin
            drop_dbg_seen <= 1'b1;
            $display("SER96TO24_DROP_FIRST: buffer full cnt=%0d t=%0t", cnt, $time);
        end
    end
    // synthesis translate_on

endmodule
\end{Verbatim}
\subsection{pulse\_sync.v}
\fileintro{单 bit 脉冲跨域（toggle 翻转 + 目标域 3 级同步 + 边沿还原），输出固定 1 拍脉冲。要求源脉冲间隔 > 目标域 2 拍（帧级脉冲天然满足）。当前设计中 T1 修正后已无实例化，保留备用。}
\begin{Verbatim}[numbers=left,numbersep=5pt,frame=single,fontsize=\footnotesize,breaklines=true,breakanywhere=true]
// =============================================================================
// pulse_sync.v —— 单 bit 脉冲跨时钟域（toggle 同步器）
// -----------------------------------------------------------------------------
// 源域脉冲 -> 目标域固定 1 拍脉冲（外脑评审措辞修正：与源脉冲宽度无关）。用于 frame_tick（检测域 -> 云台/PWM 域）。
// 约束：源脉冲间隔 > 目标域 3 个时钟周期（帧级脉冲天然满足）。
// =============================================================================
module pulse_sync (
    input  wire I_clk_src,
    input  wire I_rst_src_n,
    input  wire I_pulse_src,
    input  wire I_clk_dst,
    input  wire I_rst_dst_n,
    output wire O_pulse_dst
);

    reg toggle_src;
    always @(posedge I_clk_src or negedge I_rst_src_n) begin
        if (!I_rst_src_n)
            toggle_src <= 1'b0;
        else if (I_pulse_src)
            toggle_src <= ~toggle_src;
    end

    reg [2:0] sync_dst;
    always @(posedge I_clk_dst or negedge I_rst_dst_n) begin
        if (!I_rst_dst_n)
            sync_dst <= 3'b000;
        else
            sync_dst <= {sync_dst[1:0], toggle_src};
    end

    assign O_pulse_dst = sync_dst[1] ^ sync_dst[2];

endmodule
\end{Verbatim}
\subsection{cdc\_4phase.v}
\fileintro{N bit 数据跨域（4 相电平握手）：装载拍寄存数据、握手中恒定，目标域 2FF 同步 req、应答后输出 1 拍 valid。用于 135bit 目标包与 23bit 控制包跨 74.25M<->50M。复位契约：两侧复位必须同源（允许释放偏斜，已由 tb\_cdc\_reset\_skew 固化验证）。}
\begin{Verbatim}[numbers=left,numbersep=5pt,frame=single,fontsize=\footnotesize,breaklines=true,breakanywhere=true]
// =============================================================================
// cdc_4phase.v —— N bit 数据跨时钟域（4 相握手，电平比较式）
// -----------------------------------------------------------------------------
// 用途：检测域 -> OSD/控制域 的目标坐标/状态传递（每帧一次，低速）。
// 协议：req/ack 为电平（翻转式）握手。源域 load 时锁存数据并翻转 req；
// 目标域发现 req != ack（2FF 同步后）即锁存数据、翻转 ack、给出 1 拍
// O_dst_valid；源域发现 ack_2ff == req（对端已应答）即解除 busy。
// 数据在 req 翻转前即稳定，且整个握手期间保持不变——跨域采样安全。
//
// 【复位契约（三版外部评审共识，务必遵守）】
//   两侧复位必须**同源**（同一板级复位网络）：允许存在释放时刻的域间偏斜
//   （vision_top 内已用各域同步释放 + tb_cdc_reset_skew 固化验证），
//   但**禁止运行时单独复位其中一侧**——单侧复位会使 req/ack 失配，
//   产生一次伪 O_dst_valid（重复消费上一包）。若未来确需独立复位，
//   必须在此模块内增加 epoch/re-arm 协议。
// =============================================================================
module cdc_4phase #(
    parameter WIDTH = 67
)(
    // 源域
    input  wire             I_clk_src,
    input  wire             I_rst_src_n,
    input  wire [WIDTH-1:0] I_src_data,
    input  wire             I_src_load,
    output wire             O_src_busy,
    // 目标域
    input  wire             I_clk_dst,
    input  wire             I_rst_dst_n,
    output reg  [WIDTH-1:0] O_dst_data,
    output reg              O_dst_valid
);

    reg ack_dst;    // 目标域应答电平（跨域线，前向声明）

    // ---------------- 源域 ----------------
    reg [WIDTH-1:0] data_hold;
    reg             req;
    reg             busy;
    reg             ack_meta, ack_2ff;

    always @(posedge I_clk_src or negedge I_rst_src_n) begin
        if (!I_rst_src_n) begin
            data_hold <= {WIDTH{1'b0}};
            req       <= 1'b0;
            busy      <= 1'b0;
            ack_meta  <= 1'b0;
            ack_2ff   <= 1'b0;
        end else begin
            ack_meta <= ack_dst;
            ack_2ff  <= ack_meta;
            if (!busy) begin
                if (I_src_load) begin
                    data_hold <= I_src_data;
                    req       <= ~req;
                    busy      <= 1'b1;
                end
            end else if (ack_2ff == req) begin
                busy <= 1'b0;            // 对端已确认（ack 已翻到 req 新值）
            end
        end
    end

    assign O_src_busy = busy;

    // ---------------- 目标域 ----------------
    reg req_meta, req_2ff;

    always @(posedge I_clk_dst or negedge I_rst_dst_n) begin
        if (!I_rst_dst_n) begin
            req_meta    <= 1'b0;
            req_2ff     <= 1'b0;
            ack_dst     <= 1'b0;
            O_dst_data  <= {WIDTH{1'b0}};
            O_dst_valid <= 1'b0;
        end else begin
            req_meta <= req;
            req_2ff  <= req_meta;
            O_dst_valid <= 1'b0;
            if (req_2ff != ack_dst) begin
                O_dst_data  <= data_hold;   // req 变化时数据早已稳定
                ack_dst     <= ~ack_dst;    // 翻到与 req 一致
                O_dst_valid <= 1'b1;
            end
        end
    end

endmodule
\end{Verbatim}
\subsection{rgb2hsv.v}
\fileintro{RGB888->HSV 2 级流水。定点方案：V=max、S=delta（未归一化代理）、H 三扇区（R>G>B 优先级）+ 倒数表 inv[d] = floor(3840/d)、有符号乘、算术右移。已知简化：H 只落三段区间（0..60 / 120..180 / 240..300），品红归红扇区（方案对比已记录）。与 Python 金标准位精确（10010 向量回归）。}
\begin{Verbatim}[numbers=left,numbersep=5pt,frame=single,fontsize=\footnotesize,breaklines=true,breakanywhere=true]
// =============================================================================
// rgb2hsv.v —— RGB888 -> HSV 色彩空间转换（2 级流水，1 像素/拍）
// -----------------------------------------------------------------------------
// 定点算术与 tools/golden_model.py::rgb_to_hsv 逐位一致（见文件头注释）：
//   V = max(R,G,B); S = max-min（未归一化饱和度代理）
//   inv[d] = floor(3840/d) 查表；h = (差值*inv)>>>6 后截断到 [0,60]
//   3 扇区优先级 R>G>B；delta==0 时 H=0。
// 已知特性：H 仅落在 {0..60, 120..180, 240..300}（详见 01_方案对比）。
// 资源：256x12 ROM + 1 个 9x13 有符号乘法（可综合到 DSP/逻辑）。
// =============================================================================
module rgb2hsv (
    input  wire        I_clk,
    input  wire        I_rst_n,
    input  wire        I_valid,
    input  wire        I_sof,
    input  wire        I_eol,
    input  wire [23:0] I_data,      // {R[7:0], G[7:0], B[7:0]}
    output reg         O_valid,
    output reg         O_sof,
    output reg         O_eol,
    output reg  [8:0]  O_h,         // 0..359
    output reg  [7:0]  O_s,         // = delta
    output reg  [7:0]  O_v          // = max
);

    // ---------------- 第 0 拍：max/min/差值 ----------------
    wire [7:0] r = I_data[23:16];
    wire [7:0] g = I_data[15:8];
    wire [7:0] b = I_data[7:0];

    wire [7:0] rg_max = (r > g) ? r : g;
    wire [7:0] mx     = (rg_max > b) ? rg_max : b;
    wire [7:0] rg_min = (r < g) ? r : g;
    wire [7:0] mn     = (rg_min < b) ? rg_min : b;
    wire [7:0] d      = mx - mn;

    wire is_r = (r >= g) && (r >= b);
    wire is_g = ~is_r && (g >= r) && (g >= b);

    reg signed [8:0] diff_r;     // 扇区所需差值
    always @(*) begin
        if (is_r)      diff_r = $signed({1'b0, g}) - $signed({1'b0, b}); // G-B
        else if (is_g) diff_r = $signed({1'b0, b}) - $signed({1'b0, r}); // B-R
        else           diff_r = $signed({1'b0, r}) - $signed({1'b0, g}); // R-G
    end

    reg [7:0] d_r;
    reg       isg_r, isb_r;
    reg signed [8:0] diff_r1;
    reg [7:0] v_r, s_r;
    reg       valid_r, sof_r, eol_r;

    always @(posedge I_clk or negedge I_rst_n) begin
        if (!I_rst_n) begin
            d_r <= 8'd0; isg_r <= 1'b0; isb_r <= 1'b0; diff_r1 <= 9'sd0;
            v_r <= 8'd0; s_r <= 8'd0;
            valid_r <= 1'b0; sof_r <= 1'b0; eol_r <= 1'b0;
        end else begin
            d_r     <= d;
            isg_r   <= is_g;
            isb_r   <= ~is_r && ~is_g;
            diff_r1 <= diff_r;
            v_r     <= mx;
            s_r     <= d;
            valid_r <= I_valid;
            sof_r   <= I_sof;
            eol_r   <= I_eol;
        end
    end

    // ---------------- 第 1 拍：倒数查表 + 乘法 + 扇区 ----------------
    wire [11:0] inv;
    hsv_inv_lut u_inv (.I_d(d_r), .O_inv(inv));

    wire signed [12:0] inv_s  = {1'b0, inv};                    // 0..3840
    wire signed [21:0] prod   = diff_r1 * inv_s;                // 有符号乘
    wire signed [21:0] h_raw  = prod >>> 6;                     // floor
    wire signed [6:0]  h_clip = (h_raw < 0)   ? 7'sd0 :
                                (h_raw > 60)  ? 7'sd60 :
                                                h_raw[6:0];
    wire [8:0] h_sec = isg_r ? (9'd120 + {2'b0, h_clip}) :
                       isb_r ? (9'd240 + {2'b0, h_clip}) :
                               ({3'b0, h_clip});
    wire [8:0] h_fin = (d_r == 8'd0) ? 9'd0 : h_sec;

    always @(posedge I_clk or negedge I_rst_n) begin
        if (!I_rst_n) begin
            O_valid <= 1'b0; O_sof <= 1'b0; O_eol <= 1'b0;
            O_h <= 9'd0; O_s <= 8'd0; O_v <= 8'd0;
        end else begin
            O_valid <= valid_r;
            O_sof   <= sof_r;
            O_eol   <= eol_r;
            O_h     <= h_fin;
            O_s     <= s_r;
            O_v     <= v_r;
        end
    end

endmodule
\end{Verbatim}
\subsection{hsv\_threshold.v}
\fileintro{HSV 阈值分割（1 拍）：8 组预置 ROM（H 窗口支持环绕：红色 335..25 用"或"判断判），mask = H in 窗 \&\& S>=s\_lo \&\& v\_lo<=V<=v\_hi。预置组在 detect\_top 帧首锁存。}
\begin{Verbatim}[numbers=left,numbersep=5pt,frame=single,fontsize=\footnotesize,breaklines=true,breakanywhere=true]
// =============================================================================
// hsv_threshold.v —— HSV 阈值分割（8 组预置阈值，1 拍延迟）
// -----------------------------------------------------------------------------
// mask = H in [h_lo..h_hi]（支持环绕窗口，用于红色 335..25）且 S>=s_lo
//        且 v_lo<=V<=v_hi
// 阈值组与 golden_model.py::HSV_PRESETS 一致；I_preset_idx 由按键/拨码
// 切换（上板时接嵌入式方向的键控模块）。
// =============================================================================
module hsv_threshold (
    input  wire       I_clk,
    input  wire       I_rst_n,
    input  wire       I_valid,
    input  wire       I_sof,
    input  wire       I_eol,
    input  wire [8:0] I_h,
    input  wire [7:0] I_s,
    input  wire [7:0] I_v,
    input  wire [2:0] I_preset_idx,
    output reg        O_valid,
    output reg        O_sof,
    output reg        O_eol,
    output reg        O_mask
);

    reg [8:0] h_lo, h_hi;
    reg [7:0] s_lo, v_lo, v_hi;

    // 8 组预置阈值（位精确对齐 golden_model.HSV_PRESETS）
    always @(*) begin
        case (I_preset_idx)
            3'd0: begin h_lo = 9'd335; h_hi = 9'd25;  s_lo = 8'd80; v_lo = 8'd60;  v_hi = 8'd255; end // 红
            3'd1: begin h_lo = 9'd110; h_hi = 9'd145; s_lo = 8'd60; v_lo = 8'd40;  v_hi = 8'd255; end // 绿
            3'd2: begin h_lo = 9'd200; h_hi = 9'd250; s_lo = 8'd60; v_lo = 8'd40;  v_hi = 8'd255; end // 蓝
            3'd3: begin h_lo = 9'd20;  h_hi = 9'd70;  s_lo = 8'd80; v_lo = 8'd80;  v_hi = 8'd255; end // 橙黄
            3'd4: begin h_lo = 9'd170; h_hi = 9'd195; s_lo = 8'd60; v_lo = 8'd60;  v_hi = 8'd255; end // 青
            3'd5: begin h_lo = 9'd300; h_hi = 9'd330; s_lo = 8'd60; v_lo = 8'd40;  v_hi = 8'd255; end // 品红
            3'd6: begin h_lo = 9'd0;   h_hi = 9'd359; s_lo = 8'd10; v_lo = 8'd220; v_hi = 8'd255; end // 白
            default: begin h_lo = 9'd0; h_hi = 9'd359; s_lo = 8'd0; v_lo = 8'd0; v_hi = 8'd60;    end // 黑
        endcase
    end

    wire in_h;
    assign in_h = (h_lo <= h_hi) ? (I_h >= h_lo) && (I_h <= h_hi)
                                 : (I_h >= h_lo) || (I_h <= h_hi);
    wire mask_c = in_h && (I_s >= s_lo) && (I_v >= v_lo) && (I_v <= v_hi);

    always @(posedge I_clk or negedge I_rst_n) begin
        if (!I_rst_n) begin
            O_valid <= 1'b0; O_sof <= 1'b0; O_eol <= 1'b0; O_mask <= 1'b0;
        end else begin
            O_valid <= I_valid;
            O_sof   <= I_sof;
            O_eol   <= I_eol;
            O_mask  <= mask_c;
        end
    end

endmodule
\end{Verbatim}
\subsection{bin\_smooth3x3.v}
\fileintro{二值掩码 3x3 多数表决去噪（v7 索引行缓存自对齐结构）。读写同指针行缓存 + 2 拍水平抽头；边界复制与行尾窗口用抽头寄存器补出，输出流连续（每行恰 W 像素）。行缓存需 LUTRAM（组合读），TD 综合后须核对映射（已在上板检查单）。v3->v7 四轮相位缺陷的重构复盘见迭代记录。}
\begin{Verbatim}[numbers=left,numbersep=5pt,frame=single,fontsize=\footnotesize,breaklines=true,breakanywhere=true]
// =============================================================================
// bin_smooth3x3.v —— 二值掩码 3x3 多数表决平滑（v7：索引行缓存，自对齐结构）
// -----------------------------------------------------------------------------
// 规则（与 golden_model.py::smooth3x3_majority 逐位一致）：
//   窗口 9 bit 求和 >= 5 -> 输出 1；边界行/列复制（replicate）。
//
// 架构（v7，采用官方例程 edge_detect_96 同款"索引读写行缓存"）：
//   line1_arr[x] 存 row y-2、line0_arr[x] 存 row y-1，读写同用列指针 x：
//   先读旧值（d1/d0 = 上两行的当前列）后写新值（line0_arr[x]<=mask、
//   line1_arr[x]<=d0），读指针==写指针 -> 垂直数据天然对齐，不存在
//   移位延迟线的 +/-1 相位问题（v3~v6 的教训，见 04_迭代文件）。
//   水平方向每行用 2 拍 FF 抽头（x-2, x-1）+ 当拍 x 组成 3 列窗口。
// 输出（每令牌恰一输出，流连续有序）：
//     x>=1          -> 中心 (x-1, y)   列 {x-2, x-1, x}，行 {y-2, y-1, y}
//     x==0 且已启动 -> 中心 (W-1, y-1) 列 {W-2, W-1, W-1复制}，行 {y-3, y-2, y-1}
//   两个窗口全部取自抽头寄存器（行尾窗口不读行缓存），与输出同拍组合。
// =============================================================================
module bin_smooth3x3 #(
    parameter IMG_WIDTH  = 1280,
    parameter IMG_HEIGHT = 720
)(
    input  wire I_clk,
    input  wire I_rst_n,
    input  wire I_valid,
    input  wire I_sof,
    input  wire I_eol,
    input  wire I_mask,
    output reg  O_valid,
    output reg  O_sof,
    output reg  O_eol,
    output reg  O_mask
);

    localparam W = IMG_WIDTH;

    // ---------------- 坐标计数 ----------------
    reg [10:0] x;
    reg [9:0]  y;
    reg        started;

    always @(posedge I_clk or negedge I_rst_n) begin
        if (!I_rst_n) begin
            x <= 11'd0; y <= 10'd0; started <= 1'b0;
        end else if (I_valid) begin
            if (I_sof) begin
                x <= 11'd1; y <= 10'd0; started <= 1'b1;
            end else if (I_eol) begin
                x <= 11'd0;
                y <= (y == IMG_HEIGHT - 1) ? 10'd0 : y + 10'd1;
            end else begin
                x <= x + 11'd1;
            end
        end
    end

    // ---------------- 索引行缓存（分布式 RAM） ----------------
    (* ram_style = "distributed" *) reg line0_arr [0:W-1];   // row y-1
    (* ram_style = "distributed" *) reg line1_arr [0:W-1];   // row y-2

    // 读旧值（写前读，读指针==写指针==x）
    wire d0 = line0_arr[x];          // row y-1 列 x
    wire d1 = line1_arr[x];          // row y-2 列 x

    // 写新值：当前行进 line0，被挤出的 d0 进 line1（行滚动）
    always @(posedge I_clk) begin
        if (I_valid) begin
            line1_arr[x] <= d0;
            line0_arr[x] <= I_mask;
        end
    end

    // ---------------- 行选择（垂直复制规则） ----------------
    // t* 管道载运 row y-2、m* 载运 row y-1、b* 载运 row y（当前行）
    wire t_sel_r = (y == 10'd0) ? I_mask :
                   (y == 10'd1) ? d0    : d1;
    wire m_sel_r = (y == 10'd0) ? I_mask : d0;
    wire b_sel_r = I_mask;

    // ---------------- 水平 FF 抽头 ----------------
    reg t0, t1, m0, m1, b0, b1;
    always @(posedge I_clk or negedge I_rst_n) begin
        if (!I_rst_n) begin
            {t0,t1} <= 2'b00; {m0,m1} <= 2'b00; {b0,b1} <= 2'b00;
        end else if (I_valid) begin
            t0 <= t_sel_r; m0 <= m_sel_r; b0 <= b_sel_r;
            t1 <= t0;      m1 <= m0;      b1 <= b0;
            if (x == 11'd0) begin
                // 左边界复制：列 -1 复用列 0 的值（中心 0 的窗口需要）
                t1 <= t_sel_r; m1 <= m_sel_r; b1 <= b_sel_r;
            end
        end
    end

    // ---------------- 窗口组合（与输出同拍） ----------------
    // 主窗口：中心 (x-1, y)（x>=1），列 {x-2, x-1, x}
    wire [3:0] win_main = {3'd0, t1} + {3'd0, t0} + {3'd0, t_sel_r} +
                          {3'd0, m1} + {3'd0, m0} + {3'd0, m_sel_r} +
                          {3'd0, b1} + {3'd0, b0} + {3'd0, b_sel_r};
    // 行尾窗口：中心 (W-1, y-1)（x==0），列 {W-2, W-1, W-1复制}，全用抽头
    wire [3:0] win_last = {3'd0, t1} + {3'd0, t0} + {3'd0, t0} +
                          {3'd0, m1} + {3'd0, m0} + {3'd0, m0} +
                          {3'd0, b1} + {3'd0, b0} + {3'd0, b0};

    wire main_en = (x >= 11'd1);
    wire last_en = (x == 11'd0) && started;

    // ---------------- 输出 ----------------
    always @(posedge I_clk or negedge I_rst_n) begin
        if (!I_rst_n) begin
            O_valid <= 1'b0; O_sof <= 1'b0; O_eol <= 1'b0; O_mask <= 1'b0;
        end else begin
            O_valid <= I_valid && (main_en || last_en);
            O_sof   <= I_valid && main_en && (y == 10'd0) && (x == 11'd1);
            O_eol   <= I_valid && last_en;      // 行尾窗口即每行最后一个输出
            O_mask  <= I_valid && (main_en ? (win_main >= 4'd5) :
                                   (last_en ? (win_last >= 4'd5) : 1'b0));
        end
    end

endmodule
\end{Verbatim}
\subsection{centroid\_accum.v}
\fileintro{帧内累加质心/包围框 + 帧末恢复法除法（32 位商，X/Y 两轴串行）。关键设计：帧末 tick 拍全量快照（snap\_*），除法与结果锁存只读快照——避免下一帧 sof 清零竞争（迭代 8 抓出的真缺陷）。商末位与最后一次移位同拍预测锁存。}
\begin{Verbatim}[numbers=left,numbersep=5pt,frame=single,fontsize=\footnotesize,breaklines=true,breakanywhere=true]
// =============================================================================
// centroid_accum.v —— 帧内累加求质心/包围框 + 帧末除法（除零保护）
// -----------------------------------------------------------------------------
// 与 golden_model.py::centroid_bbox 逐位一致：
//   count = summask；sum_x = sum x·mask；sum_y = sum y·mask（对平滑后掩码）
//   count >= MIN_PIXELS 才有效；cx = floor(sum_x/count)，cy 同理。
// 帧定界：内部 x/y 计数器（与 bin_smooth3x3 输出约定一致：sof/eol 自带）。
// 帧末（最后一行 eol 后一拍）启动恢复法除法状态机，两轴共 ~2*DIVW 拍，
// 在场消隐期内完成；结果锁存到下一帧。lost_cnt 记录连续无效帧数。
// =============================================================================
module centroid_accum #(
    parameter IMG_WIDTH  = 1280,
    parameter IMG_HEIGHT = 720,
    parameter MIN_PIXELS = 30
)(
    input  wire        I_clk,
    input  wire        I_rst_n,
    input  wire        I_valid,
    input  wire        I_sof,
    input  wire        I_eol,
    input  wire        I_mask,
    output reg         O_target_valid,
    output reg  [10:0] O_cx,
    output reg  [10:0] O_cy,
    output reg  [20:0] O_count,
    output reg  [10:0] O_min_x,
    output reg  [10:0] O_max_x,
    output reg  [9:0]  O_min_y,
    output reg  [9:0]  O_max_y,
    output reg  [15:0] O_lost_cnt,
    output reg         O_results_upd,   // 结果刷新脉冲（除法完成拍）
    output wire        O_frame_tick     // 帧末脉冲（帧最后一个像素后一拍）
);

    // ---------------- 像素坐标计数 ----------------
    reg [10:0] x;
    reg [9:0]  y;
    wire last_line = (y == IMG_HEIGHT - 1);

    always @(posedge I_clk or negedge I_rst_n) begin
        if (!I_rst_n) begin
            x <= 11'd0; y <= 10'd0;
        end else if (I_valid) begin
            if (I_eol) begin
                x <= 11'd0;
                y <= last_line ? 10'd0 : y + 10'd1;
            end else begin
                x <= x + 11'd1;
            end
        end
    end

    // ---------------- 帧累加器 ----------------
    reg [31:0] sum_x, sum_y;
    reg [20:0] cnt;
    reg        seen;
    reg [10:0] bmin_x, bmax_x;
    reg [9:0]  bmin_y, bmax_y;

    wire hit = I_valid && I_mask;

    always @(posedge I_clk or negedge I_rst_n) begin
        if (!I_rst_n) begin
            sum_x <= 32'd0; sum_y <= 32'd0; cnt <= 21'd0; seen <= 1'b0;
            bmin_x <= 11'd0; bmax_x <= 11'd0; bmin_y <= 10'd0; bmax_y <= 10'd0;
        end else if (I_valid && I_sof) begin
            sum_x <= 32'd0; sum_y <= 32'd0; cnt <= 21'd0; seen <= 1'b0; // 帧首清零
        end else if (hit) begin
            sum_x <= sum_x + {21'd0, x};
            sum_y <= sum_y + {22'd0, y};
            cnt   <= cnt + 21'd1;
            if (!seen) begin
                seen   <= 1'b1;
                bmin_x <= x; bmax_x <= x; bmin_y <= y; bmax_y <= y;
            end else begin
                if (x < bmin_x) bmin_x <= x;
                if (x > bmax_x) bmax_x <= x;
                if (y < bmin_y) bmin_y <= y;
                if (y > bmax_y) bmax_y <= y;
            end
        end
    end

    // ---------------- 帧末脉冲 ----------------
    reg tick_r;
    always @(posedge I_clk or negedge I_rst_n) begin
        if (!I_rst_n)
            tick_r <= 1'b0;
        else
            tick_r <= I_valid && I_eol && last_line;
    end
    assign O_frame_tick = tick_r;

    // ---------------- 恢复法除法状态机（sum/cnt，两轴串行） ----------------
    localparam S_IDLE = 2'd0, S_DIV_X = 2'd1, S_DIV_Y = 2'd2, S_LATCH = 2'd3;
    reg [1:0]  st;
    reg [31:0] dividend;
    reg [20:0] divisor;
    reg [31:0] quotient;
    reg [21:0] remainder;
    reg [5:0]  bit_cnt;      // 32 位商逐位恢复
    // 除法快照：tick 时刻定格，下一帧 sof 清零累加器不影响本次结果
    reg [31:0] snap_sumx, snap_sumy;
    reg [20:0] snap_cnt;
    reg [10:0] snap_minx, snap_maxx;
    reg [9:0]  snap_miny, snap_maxy;

    wire [53:0] rem_sh   = {remainder, dividend[31]};
    wire [21:0] rem_sub  = rem_sh[21:0] - divisor;
    wire        rem_ge   = (rem_sh[21:0] >= divisor);

    always @(posedge I_clk or negedge I_rst_n) begin
        if (!I_rst_n) begin
            st <= S_IDLE; dividend <= 32'd0; divisor <= 21'd0;
            quotient <= 32'd0; remainder <= 22'd0; bit_cnt <= 6'd0;
        end else begin
            case (st)
                S_IDLE: if (tick_r) begin
                    snap_sumx <= sum_x; snap_sumy <= sum_y;
                    snap_cnt  <= cnt;
                    snap_minx <= bmin_x; snap_maxx <= bmax_x;
                    snap_miny <= bmin_y; snap_maxy <= bmax_y;
                    if (cnt < MIN_PIXELS) begin
                        st <= S_LATCH;                    // 无目标：直接锁存无效
                    end else begin
                        dividend <= sum_x; divisor <= cnt;
                        quotient <= 32'd0; remainder <= 22'd0;
                        bit_cnt  <= 6'd32;
                        st       <= S_DIV_X;
                    end
                end
                S_DIV_X, S_DIV_Y: begin
                    if (rem_ge) begin
                        remainder <= rem_sub;
                        quotient  <= {quotient[30:0], 1'b1};
                    end else begin
                        remainder <= rem_sh[21:0];
                        quotient  <= {quotient[30:0], 1'b0};
                    end
                    dividend <= {dividend[30:0], 1'b0};
                    bit_cnt  <= bit_cnt - 6'd1;
                    if (bit_cnt == 6'd1) begin
                        remainder <= 22'd0;
                        if (st == S_DIV_X) begin
                            dividend <= snap_sumy;          // 切换 Y 轴（用快照）
                            quotient <= 32'd0;
                            bit_cnt  <= 6'd32;
                            st       <= S_DIV_Y;
                        end else begin
                            st <= S_LATCH;
                        end
                    end
                end
                S_LATCH: st <= S_IDLE;
                default: st <= S_IDLE;
            endcase
        end
    end

    // ---------------- 结果锁存 ----------------
    // 商末位与移位同步预测（与 quotient 最后一次移位同拍）
    reg [31:0] q_x, q_y;
    always @(posedge I_clk) begin
        if (st == S_DIV_X && bit_cnt == 6'd1)
            q_x <= {quotient[30:0], rem_ge};
        if (st == S_DIV_Y && bit_cnt == 6'd1)
            q_y <= {quotient[30:0], rem_ge};
    end

    always @(posedge I_clk or negedge I_rst_n) begin
        if (!I_rst_n) begin
            O_target_valid <= 1'b0; O_cx <= 11'd0; O_cy <= 11'd0;
            O_count <= 21'd0; O_lost_cnt <= 16'd0;
            O_min_x <= 11'd0; O_max_x <= 11'd0; O_min_y <= 10'd0; O_max_y <= 10'd0;
            O_results_upd <= 1'b0;
        end else begin
            O_results_upd <= 1'b0;
            if (st == S_LATCH) begin
                O_results_upd <= 1'b1;
                if (snap_cnt >= MIN_PIXELS) begin
                    O_target_valid <= 1'b1;
                    O_cx <= q_x[10:0];
                    O_cy <= q_y[10:0];
                    O_count <= snap_cnt;
                    O_min_x <= snap_minx; O_max_x <= snap_maxx;
                    O_min_y <= snap_miny; O_max_y <= snap_maxy;
                    O_lost_cnt <= 16'd0;
                end else begin
                    O_target_valid <= 1'b0;
                    O_cx <= 11'd0; O_cy <= 11'd0;
                    O_count <= snap_cnt;
                    O_min_x <= 11'd0; O_max_x <= 11'd0;
                    O_min_y <= 10'd0; O_max_y <= 10'd0;
                    if (O_lost_cnt != 16'hFFFF)
                        O_lost_cnt <= O_lost_cnt + 16'd1;
                end
            end
        end
    end

    // synthesis translate_off
    always @(posedge I_clk) begin
        if (I_rst_n && I_valid && I_sof && (x != 11'd0) )
            $display("CENTROID_ALIGN_WARN: sof at x=%0d y=%0d t=%0t", x, y, $time);
    end
    // synthesis translate_on

endmodule
\end{Verbatim}
\subsection{detect\_top.v}
\fileintro{检测分支顶层：96bit 窃流纯旁路（O\_*=I\_* 逐拍一致，不反压）+ 检测链串联。含进阶项：fid 帧号计数器、motion\_detect 分支挂载、串化后 1ppc 流引出（OSD 底图用）。}
\begin{Verbatim}[numbers=left,numbersep=5pt,frame=single,fontsize=\footnotesize,breaklines=true,breakanywhere=true]
// =============================================================================
// detect_top.v —— 目标检测分支顶层（96bit 窃流 -> 1ppc 检测链 -> 坐标输出）
// -----------------------------------------------------------------------------
// 挂载点：官方链路 isp_top 输出（AWB 之后、data128_96 之前）的 96bit
// AXI4-Stream。本模块为纯旁路：O_* 直通输入（时序零改动），检测链
// ser96to24 -> rgb2hsv -> hsv_threshold -> bin_smooth3x3 -> centroid_accum
// 并行工作，不反压、不修改视频流（OSD 在显示侧叠加）。
// 输出：每帧目标坐标/包围框/计数 + 帧脉冲（供云台控制与 OSD 用）。
// =============================================================================
module detect_top #(
    parameter IMG_WIDTH   = 1280,
    parameter IMG_HEIGHT  = 720,
    parameter MIN_PIXELS  = 30
)(
    input  wire        I_clk,
    input  wire        I_rst_n,
    // 96bit AXI4-Stream 视频输入（窃流，不影响主线）
    input  wire [95:0] I_tdata,
    input  wire        I_tvalid,
    input  wire        I_tuser,
    input  wire        I_tlast,
    output wire        O_tready,
    // 直通输出（与输入逐拍一致，供继续传输/比对）
    output wire [95:0] O_tdata,
    output wire        O_tvalid,
    output wire        O_tuser,
    output wire        O_tlast,
    // 阈值组选择（接按键/拨码）
    input  wire [2:0]  I_preset_idx,
    // 检测结果（帧率更新，帧末锁存）
    output wire         O_target_valid,
    output wire [10:0]  O_cx,
    output wire [10:0]  O_cy,
    output wire [20:0]  O_count,
    output wire [10:0]  O_min_x,
    output wire [10:0]  O_max_x,
    output wire [9:0]   O_min_y,
    output wire [9:0]   O_max_y,
    output wire [15:0]  O_lost_cnt,
    output wire         O_results_upd,
    output wire         O_frame_tick,
    // 调试：1ppc 掩码流（可选观察）
    output wire         O_dbg_mask_valid,
    output wire         O_dbg_mask,
    // 进阶④：运动检测输出（帧级锁存）
    output wire         O_mot_valid,
    output wire [13:0]  O_mot_cnt,
    output wire [10:0]  O_mot_minx,
    output wire [10:0]  O_mot_maxx,
    output wire [9:0]   O_mot_miny,
    output wire [9:0]   O_mot_maxy,
    // 串化后 1ppc 原始流（供 OSD 叠加层取底图）
    output wire         O_s_valid,
    output wire         O_s_sof,
    output wire         O_s_eol,
    output wire [23:0]  O_s_data,
    // 进阶③：采集帧号（每输入帧 sof 自增，随目标包携带用于延迟测量）
    output wire [7:0]   O_fid,
    // 串化器帧同步错误（v4.1 看门狗：帧契约破坏后的持久丢弃态指示；0=正常）
    output wire         O_sync_err
);

    assign O_tready = 1'b1;
    assign O_tdata  = I_tdata;
    assign O_tvalid = I_tvalid;
    assign O_tuser  = I_tuser;
    assign O_tlast  = I_tlast;

    // 96 -> 1 像素/拍
    wire        s_valid, s_sof, s_eol, s_sync_err;
    wire [23:0] s_data;
    ser96to24 u_ser (
        .I_clk(I_clk), .I_rst_n(I_rst_n),
        .I_tdata(I_tdata), .I_tvalid(I_tvalid),
        .I_tuser(I_tuser), .I_tlast(I_tlast),
        .O_tready(),
        .O_valid(s_valid), .O_sof(s_sof), .O_eol(s_eol), .O_data(s_data),
        .O_sync_err(s_sync_err)
    );
    assign O_sync_err = s_sync_err;

    assign O_s_valid = s_valid;
    assign O_s_sof   = s_sof;
    assign O_s_eol   = s_eol;
    assign O_s_data  = s_data;

    // 进阶③：采集帧号计数器
    reg [7:0] fid;
    always @(posedge I_clk or negedge I_rst_n) begin
        if (!I_rst_n)
            fid <= 8'd0;
        else if (I_tvalid && I_tuser)
            fid <= fid + 8'd1;
    end
    assign O_fid = fid;

    // RGB -> HSV（2 拍延迟）
    wire        h_valid, h_sof, h_eol;
    wire [8:0]  h_h;
    wire [7:0]  h_s, h_v;
    rgb2hsv u_hsv (
        .I_clk(I_clk), .I_rst_n(I_rst_n),
        .I_valid(s_valid), .I_sof(s_sof), .I_eol(s_eol), .I_data(s_data),
        .O_valid(h_valid), .O_sof(h_sof), .O_eol(h_eol),
        .O_h(h_h), .O_s(h_s), .O_v(h_v)
    );

    // 阈值组帧首锁存（外脑评审: 按键帧中切换会造成半帧旧半帧新掩码）
    reg [2:0] preset_q;
    always @(posedge I_clk or negedge I_rst_n) begin
        if (!I_rst_n)
            preset_q <= 3'd0;
        else if (I_tvalid && I_tuser)
            preset_q <= I_preset_idx;
    end

    // 阈值分割（1 拍延迟）
    wire        t_valid, t_sof, t_eol, t_mask;
    hsv_threshold u_thr (
        .I_clk(I_clk), .I_rst_n(I_rst_n),
        .I_valid(h_valid), .I_sof(h_sof), .I_eol(h_eol),
        .I_h(h_h), .I_s(h_s), .I_v(h_v),
        .I_preset_idx(preset_q),
        .O_valid(t_valid), .O_sof(t_sof), .O_eol(t_eol), .O_mask(t_mask)
    );

    // 3x3 多数表决平滑
    wire        m_valid, m_sof, m_eol, m_mask;
    bin_smooth3x3 #(
        .IMG_WIDTH(IMG_WIDTH), .IMG_HEIGHT(IMG_HEIGHT)
    ) u_smooth (
        .I_clk(I_clk), .I_rst_n(I_rst_n),
        .I_valid(t_valid), .I_sof(t_sof), .I_eol(t_eol), .I_mask(t_mask),
        .O_valid(m_valid), .O_sof(m_sof), .O_eol(m_eol), .O_mask(m_mask)
    );

    // 质心/包围框累加
    centroid_accum #(
        .IMG_WIDTH(IMG_WIDTH), .IMG_HEIGHT(IMG_HEIGHT),
        .MIN_PIXELS(MIN_PIXELS)
    ) u_cent (
        .I_clk(I_clk), .I_rst_n(I_rst_n),
        .I_valid(m_valid), .I_sof(m_sof), .I_eol(m_eol), .I_mask(m_mask),
        .O_target_valid(O_target_valid), .O_cx(O_cx), .O_cy(O_cy),
        .O_count(O_count),
        .O_min_x(O_min_x), .O_max_x(O_max_x),
        .O_min_y(O_min_y), .O_max_y(O_max_y),
        .O_lost_cnt(O_lost_cnt),
        .O_results_upd(O_results_upd),
        .O_frame_tick(O_frame_tick)
    );

    assign O_dbg_mask_valid = m_valid;
    assign O_dbg_mask       = m_mask;

    // 进阶④：帧差运动检测分支（与 centroid 共享平滑掩码流）
    wire         md_hit;
    wire [13:0]  md_cnt;
    wire [10:0]  md_minx, md_maxx;
    wire [9:0]   md_miny, md_maxy;
    motion_detect #(
        .IMG_WIDTH(IMG_WIDTH), .IMG_HEIGHT(IMG_HEIGHT), .MOT_TH(64)
    ) u_motion (
        .I_clk(I_clk), .I_rst_n(I_rst_n),
        .I_valid(m_valid), .I_sof(m_sof), .I_eol(m_eol), .I_mask(m_mask),
        .O_mot_valid(md_hit), .O_mot_cnt(md_cnt),
        .O_mot_minx(md_minx), .O_mot_maxx(md_maxx),
        .O_mot_miny(md_miny), .O_mot_maxy(md_maxy)
    );
    assign O_mot_valid = md_hit;
    assign O_mot_cnt   = md_cnt;
    assign O_mot_minx  = md_minx; assign O_mot_maxx = md_maxx;
    assign O_mot_miny  = md_miny; assign O_mot_maxy = md_maxy;


endmodule
\end{Verbatim}
\subsection{gimbal\_ctrl.v}
\fileintro{云台控制律（50MHz 域）：误差->Q8.8 比例->死区->限速->限幅，双轴 task 复用。T3 修正：was\_lost 标志使恢复首帧 D 项=0（KD 默认关）；track\_en=0 当拍回中。控制节拍由 vision\_top 提供（=新包到达拍，T1 修正）。整定：Kp=20/18、SLEW=100。}
\begin{Verbatim}[numbers=left,numbersep=5pt,frame=single,fontsize=\footnotesize,breaklines=true,breakanywhere=true]
// =============================================================================
// gimbal_ctrl.v —— 云台闭环控制律（死区 + 比例 + 限速 + 限幅，双轴独立）
// -----------------------------------------------------------------------------
// 每帧一次（I_frame_tick，来自检测域经 pulse_sync 跨域）：
//   err   = invert ? (center - target) : (target - center)   [像素]
//   step  = (err * KP) >>> 8                                  [µs, Q8.8]
//   |err| <= DZ 时 step=0（死区）；|step|>SLEW 时截到 +/-SLEW（限速）
//   duty  = clamp(duty + step, SERVO_MIN_US, SERVO_MAX_US)     [µs]
// 丢失目标（valid=0）：保持当前脉宽，lost_cnt 累加（饱和 16bit）。
// I_track_en=0 时强制回中位（安全/演示用）。上电/复位默认中位。
// 定点算术与 golden_model.py::gimbal_ctrl_step 逐位一致。
// 时钟域：接 PWM/50MHz 域（1µs=50 拍整数）。KD 预留（默认 0=纯 P）。
// =============================================================================
module gimbal_ctrl #(
    parameter KP_X         = 20,     // Q8.8 ~= 0.156 µs/px（回路增益 0.15，稳定域）
    parameter KP_Y         = 18,     // Q8.8 ~= 0.137 µs/px
    parameter KD_X        = 0,      // 预留微分项（默认关）
    parameter KD_Y        = 0,
    parameter DZ_X        = 25,     // 死区（像素）
    parameter DZ_Y        = 20,
    parameter SLEW_US      = 100,    // 单帧最大脉宽增量（µs）
    parameter SERVO_MIN_US= 700,    // 软件限位（先实测机械行程再定）
    parameter SERVO_MAX_US= 2300,
    parameter SERVO_MID_US= 1500,
    parameter CENTER_X    = 640,
    parameter CENTER_Y    = 360,
    parameter INVERT_X    = 0,      // 极性（现场整定：方向反了改这里）
    parameter INVERT_Y    = 0
)(
    input  wire        I_clk,
    input  wire        I_rst_n,
    input  wire        I_frame_tick,     // 每帧 1 拍（已跨域）
    input  wire        I_track_en,
    input  wire        I_target_valid,
    input  wire [10:0] I_cx,
    input  wire [10:0] I_cy,
    output reg  [12:0] O_duty_x_us,
    output reg  [12:0] O_duty_y_us,
    output reg  [15:0] O_lost_cnt,
    output reg  [7:0]  O_status          // [0]dz_x [1]dz_y [2]lost [3]lim_x [4]lim_y
);

    // ---------------- 误差（组合） ----------------
    wire signed [12:0] err_x_c = INVERT_X ?
        ($signed({1'b0, CENTER_X[10:0]}) - $signed({1'b0, I_cx})) :
        ($signed({1'b0, I_cx})           - $signed({1'b0, CENTER_X[10:0]}));
    wire signed [12:0] err_y_c = INVERT_Y ?
        ($signed({1'b0, CENTER_Y[10:0]}) - $signed({1'b0, I_cy})) :
        ($signed({1'b0, I_cy})           - $signed({1'b0, CENTER_Y[10:0]}));

    localparam [12:0] SLEW_P  = SLEW_US;
    localparam [12:0] SMIN_P  = SERVO_MIN_US;
    localparam [12:0] SMAX_P  = SERVO_MAX_US;
    localparam [12:0] SMID_P  = SERVO_MID_US;

    // ---------------- 单轴步进（组合，供时序块锁存） ----------------
    task calc_axis;
        input  signed [12:0] err;
        input  [15:0]        kp;
        input  signed [12:0] err_prev;
        input  [15:0]        kd;
        input                resume;   // 恢复首帧: D 项强制 0（外脑评审 R5 修正）
        input  [11:0]        dz;
        input  [12:0]        duty_in;
        output [12:0]        step;
        output               dz_hit;
        output               lim_hit;
        output [12:0]        duty_out;
        reg signed [28:0]    prod;
        reg signed [28:0]    s;
        reg signed [28:0]    nx;
        begin
            prod = err * $signed({1'b0, kp});
            s    = prod >>> 8;
            if (kd != 16'd0 && !resume) begin
                prod = (err - err_prev) * $signed({1'b0, kd});
                s    = s + (prod >>> 8);
            end
            dz_hit = (err >= -$signed({1'b0, dz})) && (err <= $signed({1'b0, dz}));
            if (dz_hit)
                s = 29'sd0;                              // 死区
            if (s > $signed({16'd0, SLEW_P}))       s = $signed({16'd0, SLEW_P});
            else if (s < -$signed({16'd0, SLEW_P})) s = -$signed({16'd0, SLEW_P});
            step  = s[12:0];
            nx    = $signed({1'b0, duty_in}) + s;
            lim_hit = 1'b0;
            if (nx > $signed({16'd0, SMAX_P})) begin
                duty_out = SMAX_P; lim_hit = 1'b1;
            end else if (nx < $signed({16'd0, SMIN_P})) begin
                duty_out = SMIN_P; lim_hit = 1'b1;
            end else begin
                duty_out = nx[12:0];
            end
        end
    endtask

    reg [12:0] step_x, step_y;
    reg [12:0] duty_x, duty_y;
    reg [12:0] dnx_x, dnx_y;
    reg        dzx, dzy, limx, limy;
    reg signed [12:0] err_prev_x, err_prev_y;
    reg        lost;
    reg        was_lost;    // 上一帧无有效目标（恢复首帧 D=0）

    always @(posedge I_clk or negedge I_rst_n) begin
        if (!I_rst_n) begin
            duty_x <= SMID_P; duty_y <= SMID_P;
            O_duty_x_us <= SMID_P; O_duty_y_us <= SMID_P;
            err_prev_x <= 13'sd0; err_prev_y <= 13'sd0;
            O_lost_cnt <= 16'd0; O_status <= 8'h04;
            dzx <= 1'b0; dzy <= 1'b0; limx <= 1'b0; limy <= 1'b0; lost <= 1'b1;
            was_lost <= 1'b1;
        end else if (I_frame_tick) begin
            if (!I_track_en) begin
                duty_x <= SMID_P;
                duty_y <= SMID_P;
                lost   = 1'b1;
                dzx = 1'b0; dzy = 1'b0; limx = 1'b0; limy = 1'b0;
                was_lost <= 1'b1;
                O_lost_cnt <= 16'd0;
            end else if (I_target_valid) begin
                calc_axis(err_x_c, KP_X, err_prev_x, KD_X, was_lost, DZ_X,
                          duty_x, step_x, dzx, limx, dnx_x);
                calc_axis(err_y_c, KP_Y, err_prev_y, KD_Y, was_lost, DZ_Y,
                          duty_y, step_y, dzy, limy, dnx_y);
                was_lost <= 1'b0;
                duty_x <= dnx_x;
                duty_y <= dnx_y;
                err_prev_x <= err_x_c;    // 供微分项（KD=0 时不参与）
                err_prev_y <= err_y_c;
                lost  = 1'b0;
                O_lost_cnt <= 16'd0;
            end else begin
                lost = 1'b1;              // 保持脉宽，只计丢失
                dzx = 1'b0; dzy = 1'b0; limx = 1'b0; limy = 1'b0;
                was_lost <= 1'b1;         // 恢复首帧 D 项由 calc_axis 置 0
                if (O_lost_cnt != 16'hFFFF)
                    O_lost_cnt <= O_lost_cnt + 16'd1;
            end
            // 外脑评审: track_en=0 时输出当拍即回中位（原写法滞后一帧）
            O_duty_x_us <= !I_track_en ? SMID_P :
                           I_target_valid ? dnx_x : duty_x;
            O_duty_y_us <= !I_track_en ? SMID_P :
                           I_target_valid ? dnx_y : duty_y;
            O_status    <= {2'b00, limy, limx, lost, dzy, dzx};
        end
    end

endmodule
\end{Verbatim}
\subsection{pwm\_servo.v}
\fileintro{双路 50Hz 舵机 PWM：单一 1,000,000 拍共周期计数器 + 双比较器（严格共周期）。T4 修正：占空比改为周期起点边界换装（防帧拍落周期中途截断脉宽）；双道限幅（500~2500 硬件保险）；上电中位 1500µs。}
\begin{Verbatim}[numbers=left,numbersep=5pt,frame=single,fontsize=\footnotesize,breaklines=true,breakanywhere=true]
// =============================================================================
// pwm_servo.v —— 双路 50Hz 舵机 PWM（共周期，上电中位，输入限幅）
// -----------------------------------------------------------------------------
// 规格（与《环境配置与文件清单》5.2 节钉死）：
//   周期 20ms（50Hz），脉宽 0.5~2.5ms 对应 0~180°；
//   时钟 50MHz：周期 = 1,000,000 拍，1µs = 50 拍，分辨率 0.02µs(20ns)；
//   两路共用同一周期计数器（严格共周期，避免 X/Y 互相拉扯）；
//   输入 duty(µs) 在 I_update（帧拍）时锁存，内部再 clamp 到 500~2500µs
//   （第二道保险，第一道在 gimbal_ctrl 的 700~2300 软限位）；
//   复位后输出中位 1500µs 脉冲，上电即安全中位。
// 极性：高电平有效脉冲。时钟域：50MHz（I_clk）。
// =============================================================================
module pwm_servo #(
    parameter CLK_HZ        = 50_000_000,
    parameter SERVO_FREQ_HZ = 50,           // 50Hz => 周期 20ms
    parameter TICKS_PER_US  = 50            // CLK_HZ/1e6
)(
    input  wire        I_clk,
    input  wire        I_rst_n,
    input  wire        I_update,          // 帧拍：锁存新占空
    input  wire [12:0] I_duty_x_us,
    input  wire [12:0] I_duty_y_us,
    output wire        O_pwm_x,           // 高电平有效
    output wire        O_pwm_y,
    output wire        O_period_tick,     // 周期起点脉冲（测量用）
    output reg  [12:0] O_duty_x_dbg,
    output reg  [12:0] O_duty_y_dbg
);

    localparam PERIOD_TICKS = CLK_HZ / SERVO_FREQ_HZ;   // 1,000,000
    localparam PW_CNT_W = 20;

    // 输入限幅（第二道保险）
    wire [12:0] dx = (I_duty_x_us > 13'd2500) ? 13'd2500 :
                     (I_duty_x_us < 13'd500)  ? 13'd500  : I_duty_x_us;
    wire [12:0] dy = (I_duty_y_us > 13'd2500) ? 13'd2500 :
                     (I_duty_y_us < 13'd500)  ? 13'd500  : I_duty_y_us;

    // 共周期计数器
    reg [PW_CNT_W-1:0] period_cnt;
    always @(posedge I_clk or negedge I_rst_n) begin
        if (!I_rst_n)
            period_cnt <= {PW_CNT_W{1'b0}};
        else if (period_cnt >= PERIOD_TICKS - 1)
            period_cnt <= {PW_CNT_W{1'b0}};
        else
            period_cnt <= period_cnt + 20'd1;
    end

    reg [12:0] duty_x, duty_y;
    reg [12:0] req_x, req_y;    // 请求值（帧拍写入）
    reg        update_d;        // 外脑评审：gimbal 的 O_duty 与 update 同沿更新，
    always @(posedge I_clk or negedge I_rst_n) begin   // 同沿采样会拿到上一帧值
        if (!I_rst_n) begin
            req_x <= 13'd1500;   req_y <= 13'd1500;
            update_d <= 1'b0;
        end else begin
            update_d <= I_update;
            if (update_d) begin
                req_x <= dx;     req_y <= dy;   // 延后一拍采到刚更新的占空
            end
        end
    end

    // 占空在周期起点换装（外脑评审: 防止帧拍落在周期中途,
    // 本周期脉宽被截断产生毛刺；变化延迟 <=1 个 20ms 周期，舵机无感）
    always @(posedge I_clk or negedge I_rst_n) begin
        if (!I_rst_n) begin
            duty_x <= 13'd1500;  duty_y <= 13'd1500;
            O_duty_x_dbg <= 13'd1500; O_duty_y_dbg <= 13'd1500;
        end else if (period_cnt == {PW_CNT_W{1'b0}}) begin
            duty_x <= req_x;     duty_y <= req_y;
            O_duty_x_dbg <= req_x; O_duty_y_dbg <= req_y;
        end
    end

    // 脉宽换算：µs -> 拍（*50 = *32 + *16 + *2）
    wire [PW_CNT_W-1:0] pulse_x = duty_x * TICKS_PER_US;
    wire [PW_CNT_W-1:0] pulse_y = duty_y * TICKS_PER_US;

    assign O_period_tick = (period_cnt == {PW_CNT_W{1'b0}});
    assign O_pwm_x = (period_cnt < pulse_x);
    assign O_pwm_y = (period_cnt < pulse_y);

endmodule
\end{Verbatim}
\subsection{osd\_track.v}
\fileintro{OSD 叠加层 v3：影子寄存器组 + sof 换装（消撕裂）；三行文本（状态/遥测/跑马灯）+ 十字准星 +双框（目标红/运动橙）。135bit 目标包经 cdc\_4phase 接收。跑马灯索引 (pos>>3) mod 63（连减法，含参数变更保险分支）。文本构建用常量除法（帧率更新，非关键路径）。2026-10-03（\#N+9）：跑马灯路径结构性流水——地址+ROM 读前瞻 4 像素、4 级移位对齐，ROM 读获整拍预算，混合级仅 3~4 级（像素域 WNS +0.541->+2.441ns；此前"数组 ROM 化"并未真正消除 mux 树，属综合器估算运气，已诚实记录）；位精确由 tb\_osd\_track 复验。}
\begin{Verbatim}[numbers=left,numbersep=5pt,frame=single,fontsize=\footnotesize,breaklines=true,breakanywhere=true]
// =============================================================================
// osd_track.v —— 跟踪信息 OSD 叠加层 v3（三行文本 + 跑马灯 + 双框 + 坐标双缓冲）
// -----------------------------------------------------------------------------
// v3 变更（2026-09-20，消撕裂）：
//   CDC 送来的目标包（pack_valid）可能落在一帧中间。v2 及之前直接更新绘制
//   寄存器，导致同帧上半部分用旧坐标、下半部分用新坐标（可见的框撕裂）。
//   v3 引入影子寄存器组：pack_valid 只写影子；绘制寄存器仅在 I_sof 拍从
//   影子整组换装（vsync 语义），整帧恒用同一快照。
//   注：sof 拍本身（像素 (0,0)）仍用旧值绘制一拍，OSD 起始在 x=16，无视觉影响。
// v2 变更（配合 step5 拓展）：
//   目标包 68bit -> 135bit（IF-VISION v1.1.0）：
//     {valid, algo, preset[2:0], cx[10:0], cy[10:0], minx[10:0], maxx[10:0],
//      miny[9:0], maxy[9:0], fid[7:0], mot_valid, mot_cnt[13:0],
//      mot_minx[10:0], mot_maxx[10:0], mot_miny[9:0], mot_maxy[9:0]}
//   三行文本：
//     L1 y=16:  "<TRK|EDG><n> <LOCK|LOST> X####Y####"
//     L2 y=40:  "D%02d B%+03d C%03d M%04d"
//     L3 y=64:  跑马灯（64 字符消息逐像素滚动，2px/帧）
//   双框：绿色准星 + 红色目标框；运动有效时叠加橙色运动包围框
// 绘制优先级：文字 > 准星 > 目标框 > 运动框 > 底图
// =============================================================================
module osd_track #(
    parameter IMG_WIDTH  = 1280,
    parameter IMG_HEIGHT = 720,
    parameter TXT_X0     = 16,
    parameter TXT_Y0     = 16,
    parameter TXT2_Y0    = 40,
    parameter TXT3_Y0    = 64
)(
    input  wire        I_clk,
    input  wire        I_rst_n,
    input  wire        I_valid,
    input  wire        I_sof,
    input  wire        I_eol,
    input  wire [23:0] I_data,
    input  wire        I_det_clk,
    input  wire        I_det_rst_n,
    input  wire [134:0] I_det_pack,
    input  wire        I_det_load,
    output wire        O_det_busy,
    input  wire [6:0]  I_bright,       // -64..+63
    input  wire [7:0]  I_contrast,     // Q6, 64=1.0
    output reg         O_valid,
    output reg         O_sof,
    output reg         O_eol,
    output reg  [23:0] O_data
);

    wire [134:0] pack_dst;
    wire         pack_valid;
    cdc_4phase #(.WIDTH(135)) u_cdc (
        .I_clk_src(I_det_clk), .I_rst_src_n(I_det_rst_n),
        .I_src_data(I_det_pack), .I_src_load(I_det_load), .O_src_busy(O_det_busy),
        .I_clk_dst(I_clk), .I_rst_dst_n(I_rst_n),
        .O_dst_data(pack_dst), .O_dst_valid(pack_valid)
    );

    reg [7:0] disp_cnt;
    always @(posedge I_clk or negedge I_rst_n) begin
        if (!I_rst_n) disp_cnt <= 8'd0;
        else if (I_valid && I_sof) disp_cnt <= disp_cnt + 8'd1;
    end

    // ---------------- 影子寄存器组（CDC 域写入，sof 换装） ----------------
    reg        s_tv, s_algo_edg;
    reg [2:0]  s_preset;
    reg [10:0] s_cx, s_cy, s_bminx, s_bmaxx;
    reg [9:0]  s_bminy, s_bmaxy;
    reg [7:0]  s_lat;
    reg        s_mv;
    reg [13:0] s_mcnt;
    reg [10:0] s_mminx, s_mmaxx;
    reg [9:0]  s_mminy, s_mmaxy;
    reg [7:0]  txt1_s [0:19];
    reg [7:0]  txt2_s [0:19];

    // ---------------- 绘制寄存器组（sof 拍整组换装，帧内恒定） ----------------
    reg        tv, algo_edg;
    reg [2:0]  preset_q;
    reg [10:0] cx, cy, bminx, bmaxx;
    reg [9:0]  bminy, bmaxy;
    reg [7:0]  lat_q;
    reg        mv;
    reg [13:0] mcnt;
    reg [10:0] mminx, mmaxx;
    reg [9:0]  mminy, mmaxy;
    reg [7:0]  txt1 [0:19];
    reg [7:0]  txt2 [0:19];
    integer    k;

    function [31:0] to_digits4;
        input [13:0] val;
        reg [13:0] r;
        begin
            to_digits4[31:24] = 8'd48 + val / 14'd1000;
            r = val % 14'd1000;
            to_digits4[23:16] = 8'd48 + r / 14'd100;
            r = r % 14'd100;
            to_digits4[15:8]  = 8'd48 + r / 14'd10;
            to_digits4[7:0]   = 8'd48 + r % 14'd10;
        end
    endfunction
    function [15:0] to_digits2;
        input [7:0] val;
        begin
            to_digits2[15:8] = 8'd48 + val / 8'd10;
            to_digits2[7:0]   = 8'd48 + val % 8'd10;
        end
    endfunction

    reg [31:0] d_x, d_y, d_mot;
    reg [15:0] d_lat;
    reg [7:0]  b_abs;
    // 外脑评审修复：D 字段应使用本包延迟（原读旧 s_lat，文本滞后一包）
    wire [7:0] lat_new = disp_cnt - pack_dst[65:58];

    // 影子写入：CDC 包到达时（可能在帧中间）
    always @(posedge I_clk or negedge I_rst_n) begin
        if (!I_rst_n) begin
            s_tv <= 0; s_algo_edg <= 0; s_preset <= 0;
            s_cx <= 0; s_cy <= 0;
            s_bminx <= 0; s_bmaxx <= 0; s_bminy <= 0; s_bmaxy <= 0;
            s_lat <= 0; s_mv <= 0; s_mcnt <= 0;
            s_mminx <= 0; s_mmaxx <= 0; s_mminy <= 0; s_mmaxy <= 0;
            for (k = 0; k < 20; k = k + 1) begin
                txt1_s[k] <= 8'd32; txt2_s[k] <= 8'd32;
            end
        end else if (pack_valid) begin
            s_tv       <= pack_dst[134];
            s_algo_edg <= pack_dst[133];
            s_preset   <= pack_dst[132:130];
            s_cx       <= pack_dst[129:119];
            s_cy       <= pack_dst[118:108];
            s_bminx    <= pack_dst[107:97];
            s_bmaxx    <= pack_dst[96:86];
            s_bminy    <= pack_dst[85:76];
            s_bmaxy    <= pack_dst[75:66];
            s_lat      <= lat_new;
            s_mv       <= pack_dst[57];
            s_mcnt     <= pack_dst[56:43];
            s_mminx    <= pack_dst[42:32];
            s_mmaxx    <= pack_dst[31:21];
            s_mminy    <= pack_dst[20:11];
            s_mmaxy    <= pack_dst[10:1];

            txt1_s[0] <= pack_dst[133] ? "E" : "T";
            txt1_s[1] <= pack_dst[133] ? "D" : "R";
            txt1_s[2] <= pack_dst[133] ? "G" : "K";
            txt1_s[3] <= 8'd48 + {5'd0, pack_dst[132:130]};
            txt1_s[4] <= " ";
            txt1_s[5] <= "L"; txt1_s[6] <= "O";
            if (pack_dst[134]) begin
                txt1_s[7] <= "C"; txt1_s[8] <= "K";
            end else begin
                txt1_s[7] <= "S"; txt1_s[8] <= "T";
            end
            txt1_s[9] <= " "; txt1_s[10] <= "X";
            d_x = to_digits4(pack_dst[134] ? pack_dst[129:119] : 11'd0);
            d_y = to_digits4(pack_dst[134] ? pack_dst[118:108] : 11'd0);
            if (!pack_dst[134]) begin
                d_x = 32'h2D2D2D2D; d_y = 32'h2D2D2D2D;
            end
            txt1_s[11] <= d_x[31:24]; txt1_s[12] <= d_x[23:16];
            txt1_s[13] <= d_x[15:8];  txt1_s[14] <= d_x[7:0];
            txt1_s[15] <= "Y";
            txt1_s[16] <= d_y[31:24]; txt1_s[17] <= d_y[23:16];
            txt1_s[18] <= d_y[15:8];  txt1_s[19] <= d_y[7:0];

            d_lat = to_digits2((lat_new > 8'd99) ? 8'd99 : lat_new);
            txt2_s[0] <= "D";
            txt2_s[1] <= d_lat[15:8]; txt2_s[2] <= d_lat[7:0];
            txt2_s[3] <= " ";
            txt2_s[4] <= "B";
            txt2_s[5] <= I_bright[6] ? "-" : "+";
            b_abs = I_bright[6] ? (~I_bright + 1'b1) : I_bright;
            d_lat = to_digits2(b_abs);
            txt2_s[6] <= d_lat[15:8]; txt2_s[7] <= d_lat[7:0];
            txt2_s[8] <= " ";
            txt2_s[9] <= "C";
            d_lat = to_digits2(I_contrast[7:3]);
            txt2_s[10] <= d_lat[15:8]; txt2_s[11] <= d_lat[7:0];
            txt2_s[12] <= " ";
            txt2_s[13] <= "M";
            // 外脑评审覆盖测试抓出的真 bug：原读 s_mcnt（同块 NBA 旧值）,
            // M 字段永远滞后一包（与 D 字段同类的"读旧值"错误）。
            // 改用本包 live 值。
            d_mot = to_digits4((pack_dst[56:43] > 14'd9999) ? 14'd9999 : pack_dst[56:43]);
            txt2_s[14] <= d_mot[31:24]; txt2_s[15] <= d_mot[23:16];
            txt2_s[16] <= d_mot[15:8];  txt2_s[17] <= d_mot[7:0];
            txt2_s[18] <= " "; txt2_s[19] <= " ";
        end
    end

    // vsync 换装：sof 拍整组影子 -> 绘制（帧内恒定，消撕裂）
    always @(posedge I_clk or negedge I_rst_n) begin
        if (!I_rst_n) begin
            tv <= 0; algo_edg <= 0; preset_q <= 0;
            cx <= 0; cy <= 0;
            bminx <= 0; bmaxx <= 0; bminy <= 0; bmaxy <= 0;
            lat_q <= 0; mv <= 0; mcnt <= 0;
            mminx <= 0; mmaxx <= 0; mminy <= 0; mmaxy <= 0;
            for (k = 0; k < 20; k = k + 1) begin
                txt1[k] <= 8'd32; txt2[k] <= 8'd32;
            end
        end else if (I_valid && I_sof) begin
            tv       <= s_tv;
            algo_edg <= s_algo_edg;
            preset_q <= s_preset;
            cx       <= s_cx;
            cy       <= s_cy;
            bminx    <= s_bminx;
            bmaxx    <= s_bmaxx;
            bminy    <= s_bminy;
            bmaxy    <= s_bmaxy;
            lat_q    <= s_lat;
            mv       <= s_mv;
            mcnt     <= s_mcnt;
            mminx    <= s_mminx;
            mmaxx    <= s_mmaxx;
            mminy    <= s_mminy;
            mmaxy    <= s_mmaxy;
            for (k = 0; k < 20; k = k + 1) begin
                txt1[k] <= txt1_s[k];
                txt2[k] <= txt2_s[k];
            end
        end
    end

    reg [10:0] px;
    reg [9:0]  py;
    always @(posedge I_clk or negedge I_rst_n) begin
        if (!I_rst_n) begin
            px <= 0; py <= 0;
        end else if (I_valid) begin
            if (I_sof) begin
                px <= 11'd1; py <= 10'd0;
            end else if (I_eol) begin
                px <= 11'd0; py <= py + 10'd1;
            end else begin
                px <= px + 11'd1;
            end
        end
    end

    // 跑马灯：消息 "2026 FPGA CONTEST - ANLOGIC PH1P35 VISION TRACKER - " 共 63 字符
    // 偏移 = disp_cnt x 2 px/帧，取模 504（63x8），字符索引 = (px-TXT_X0+off)>>3
    wire [9:0]  mar_off  = {disp_cnt[7:0], 1'b0};               // 0..510
    // 外脑评审: 原 mar_pos[8:3] 截断高位——mar_pos>=512 后字符索引错位。
    // 正确 (mar_pos>>3) mod 63: 63*3=189, 连减三次(索引最大 221)
    wire [11:0] mar_pos  = (px >= TXT_X0) ? (px - TXT_X0 + {2'b00, mar_off}) : 12'd0;
    wire [2:0]  mar_bit  = 3'd7 - mar_pos[2:0];
    wire [3:0]  mar_row  = py - TXT3_Y0;
    wire        in_mar   = (py >= TXT3_Y0) && (py < TXT3_Y0 + 16) &&
                           (px >= TXT_X0) && (px < TXT_X0 + 1280);

    // ---- 跑马灯流水（2026-10-03：TD 将 1024x8 ROM 映射为 ~10 级 LUT3 mux 链，
    //      是 74.25MHz 域的唯一关键路径。此处把"地址计算 + ROM 读"整体提前 4 像素，
    //      ROM 输出经 4 级移位与当前像素对齐——ROM 获得整拍预算，混合级仅剩
    //      位选 + 区域门控 3~4 级逻辑。功能位精确等价（tb_osd_track 逐像素复验）。----
    // 前瞻位姿 (px+4) 需处理行尾跨越：越过 1279 -> 下一行的 px/py。
    // 语义核对：n = px - TXT_X0 + mar_off（虚拟文本坐标，TXT_X0 为 8 的倍数）；
    //           字符 = (n>>3) mod 63；位选 = 7 - (n&7)。
    wire [11:0] px_4  = px + 12'd4;
    wire        xw    = (px_4 > 12'd1279) && (py < 10'd719);   // 行内跨越才换行
    wire [11:0] px_c  = xw ? (px_4 - 12'd1280) : px_4;
    wire [9:0]  py_c  = xw ? (py + 10'd1) : py;

    wire [11:0] mar_pos_c = (px_c >= TXT_X0) ? (px_c - TXT_X0 + {2'b00, mar_off}) : 12'd0;
    wire [7:0]  ci_full_c = {1'b0, mar_pos_c[10:3]};
    wire [7:0]  ci_r1_c   = (ci_full_c >= 8'd252) ? ci_full_c - 8'd252 :
                            (ci_full_c >= 8'd189) ? ci_full_c - 8'd189 :
                            (ci_full_c >= 8'd126) ? ci_full_c - 8'd126 :
                            (ci_full_c >= 8'd63)  ? ci_full_c - 8'd63  : ci_full_c;
    wire [5:0]  mar_ci_c  = ci_r1_c[5:0];
    wire [3:0]  mar_row_c = py_c - TXT3_Y0;
    wire [7:0]  mar_glyph_c;
    marquee_rom_64 u_marq (.I_addr({mar_ci_c, mar_row_c}), .O_bits(mar_glyph_c));

    // 4 级移位对齐：g3 在当拍恰为"本像素"的 ROM 输出（仅在有效像素拍推进）
    reg [7:0] mg0, mg1, mg2, mg3;
    always @(posedge I_clk or negedge I_rst_n) begin
        if (!I_rst_n) begin
            mg0 <= 8'd0; mg1 <= 8'd0; mg2 <= 8'd0; mg3 <= 8'd0;
        end else if (I_valid) begin
            mg0 <= mar_glyph_c;
            mg1 <= mg0;
            mg2 <= mg1;
            mg3 <= mg2;
        end
    end
    wire on_mar = in_mar && mg3[mar_bit];

    // L1 状态行
    wire in_txt1 = (py >= TXT_Y0) && (py < TXT_Y0 + 16) &&
                   (px >= TXT_X0) && (px < TXT_X0 + 160);
    wire [4:0]  ch1_idx = (px - TXT_X0) >> 3;
    wire [3:0]  ch1_row = py - TXT_Y0;
    wire [10:0] dxp     = px - TXT_X0;
    wire [2:0]  bit_sel = 3'd7 - dxp[2:0];
    wire [7:0]  g1;
    font_rom_8x16 u_f1 (.I_addr({txt1[ch1_idx], ch1_row}), .O_bits(g1));
    wire on_txt1 = in_txt1 && g1[bit_sel];

    // L2 遥测行
    wire in_txt2 = (py >= TXT2_Y0) && (py < TXT2_Y0 + 16) &&
                   (px >= TXT_X0) && (px < TXT_X0 + 160);
    wire [4:0]  ch2_idx = (px - TXT_X0) >> 3;
    wire [3:0]  ch2_row = py - TXT2_Y0;
    wire [7:0]  g2;
    font_rom_8x16 u_f2 (.I_addr({txt2[ch2_idx], ch2_row}), .O_bits(g2));
    wire on_txt2 = in_txt2 && g2[bit_sel];

    // 准星/框
    wire signed [12:0] dx_c = $signed({1'b0, px}) - $signed({1'b0, cx});
    wire signed [12:0] dy_c = $signed({1'b0, py}) - $signed({1'b0, cy});
    wire ax = (dx_c >= -13'sd12) && (dx_c <= 13'sd12) &&
              ((dx_c <= -13'sd3) || (dx_c >= 13'sd3));
    wire ay = (dy_c >= -13'sd12) && (dy_c <= 13'sd12) &&
              ((dy_c <= -13'sd3) || (dy_c >= 13'sd3));
    wire on_cross = tv && ((ay && (dx_c == 13'sd0)) || (ax && (dy_c == 13'sd0)));

    wire x_edge = (px == bminx) || (px == bminx + 11'd1) ||
                  (px == bmaxx - 11'd1) || (px == bmaxx);
    wire y_edge = (py == bminy) || (py == bminy + 10'd1) ||
                  (py == bmaxy - 10'd1) || (py == bmaxy);
    wire in_bx = (px >= bminx) && (px <= bmaxx);
    wire in_by = (py >= bminy) && (py <= bmaxy);
    wire on_box = tv && ((x_edge && in_by) || (y_edge && in_bx));

    wire mx_e = (px == mminx) || (px == mminx + 11'd1) ||
                (px == mmaxx - 11'd1) || (px == mmaxx);
    wire my_e = (py == mminy) || (py == mminy + 10'd1) ||
                (py == mmaxy - 10'd1) || (py == mmaxy);
    wire in_mx = (px >= mminx) && (px <= mmaxx);
    wire in_my = (py >= mminy) && (py <= mmaxy);
    wire on_mbox = mv && ((mx_e && in_my) || (my_e && in_mx));

    always @(posedge I_clk or negedge I_rst_n) begin
        if (!I_rst_n) begin
            O_valid <= 0; O_sof <= 0; O_eol <= 0; O_data <= 0;
        end else begin
            O_valid <= I_valid;
            O_sof   <= I_sof;
            O_eol   <= I_eol;
            O_data  <= on_txt1 ? 24'hFF_FF_FF :
                       on_txt2 ? 24'hFF_FF_FF :
                       on_mar  ? 24'h00_FF_FF :
                       on_cross ? 24'h00_FF_00 :
                       on_box  ? 24'hFF_00_00 :
                       on_mbox ? 24'hFF_80_00 : I_data;
        end
    end

endmodule
\end{Verbatim}
\subsection{vision\_top.v}
\fileintro{整机顶层 v2+：直通零扰动；检测/边缘双分支 + 帧首锁存 algo mux；video\_adj 画质调节（显示路径）；OSD 叠加；声光反馈；控制路径（T1 修正：tick=新包到达拍，控制滞后 1 帧）+ 云台 + PWM。目标状态（O\_frame\_tick/O\_target\_valid/O\_tgt\_cx/cy）端口化直出替代层次引用。}
\begin{Verbatim}[numbers=left,numbersep=5pt,frame=single,fontsize=\footnotesize,breaklines=true,breakanywhere=true]
// =============================================================================
// vision_top.v —— 视觉整机顶层 v2（step5 拓展后）
// -----------------------------------------------------------------------------
// v2 新增（step5）：
//   进阶① 多算法切换：edge_detect_96 并行 + algo mux（帧边界平稳切换）
//   进阶③ 延迟可视化：fid 随目标包传递，OSD 显示 D+NN
//   进阶④ 帧差运动检测：motion_detect 分支，运动包围框 OSD 橙色叠加
//   扩展④ 亮度/对比度：video_adj（显示路径，检测分支不受影响）
//   声光反馈：led_buzzer（嵌入式方向）
// 目标包 v2 = 135bit（IF-VISION v1.1.0）
// =============================================================================
module vision_top #(
    parameter IMG_WIDTH    = 1280,
    parameter IMG_HEIGHT   = 720,
    parameter MIN_PIXELS   = 30,
    parameter KP_X         = 20,
    parameter KP_Y         = 18,
    parameter DZ_X         = 25,
    parameter DZ_Y         = 20,
    parameter SLEW_US      = 100,
    parameter SERVO_MIN_US = 700,
    parameter SERVO_MAX_US = 2300,
    parameter INVERT_X     = 0,
    parameter INVERT_Y     = 0
)(
    input  wire        I_pix_clk,
    input  wire        I_pix_rst_n,
    input  wire        I_pwm_clk,
    input  wire        I_pwm_rst_n,
    input  wire [95:0] I_tdata,
    input  wire        I_tvalid,
    input  wire        I_tuser,
    input  wire        I_tlast,
    output wire        O_tready,
    output wire [95:0] O_tdata,
    output wire        O_tvalid,
    output wire        O_tuser,
    output wire        O_tlast,
    output wire        O_osd_valid,
    output wire        O_osd_sof,
    output wire        O_osd_eol,
    output wire [23:0] O_osd_data,
    input  wire [2:0]  I_preset_idx,
    input  wire        I_track_en,
    input  wire        I_algo_edge,      // 进阶①：显示算法 0=跟踪 1=边缘
    input  wire [6:0]  I_bright,         // 扩展④ 亮度
    input  wire [7:0]  I_contrast,       // 扩展④ 对比度 Q6
    output wire        O_pwm_x,
    output wire        O_pwm_y,
    output wire [12:0] O_duty_x_us,
    output wire [12:0] O_duty_y_us,
    output wire [15:0] O_lost_cnt,
    output wire [7:0]  O_status,
    output wire        O_led_g,
    output wire        O_led_r,
    output wire        O_buzzer,
    // 目标状态直出（外脑评审: 供顶层 UART 遥测等使用, 替代层次化引用）
    output wire        O_frame_tick,
    output wire        O_target_valid,
    output wire [10:0] O_tgt_cx,
    output wire [10:0] O_tgt_cy
);

    // ---------------- 复位同步释放（异步置位 / 各域同步释放） ----------------
    // 外脑评审（两版一致）：直接透传引脚复位在释放沿存在 recovery/removal 风险
    // 与域间释放偏斜；此处每域两拍同步释放，全部分子模块统一使用。
    reg [1:0] pix_rst_sync, pwm_rst_sync;
    always @(posedge I_pix_clk or negedge I_pix_rst_n)
        if (!I_pix_rst_n) pix_rst_sync <= 2'b00;
        else              pix_rst_sync <= {pix_rst_sync[0], 1'b1};
    always @(posedge I_pwm_clk or negedge I_pwm_rst_n)
        if (!I_pwm_rst_n) pwm_rst_sync <= 2'b00;
        else              pwm_rst_sync <= {pwm_rst_sync[0], 1'b1};
    wire pix_rst_n_s = pix_rst_sync[1];
    wire pwm_rst_n_s = pwm_rst_sync[1];

    wire        det_s_valid, det_s_sof, det_s_eol;
    wire [23:0] det_s_data;
    wire        det_sync_err;      // v4.1 串化器帧同步错误（检测支路）

    // ---------------- 检测分支 ----------------
    wire         t_valid, t_upd, t_tick;
    wire [10:0]  t_cx, t_cy, t_minx, t_maxx;
    wire [9:0]   t_miny, t_maxy;
    wire [20:0]  t_cnt;
    wire [15:0]  lost_det;
    wire [7:0]   t_fid;
    wire         mo_valid;
    wire [13:0]  mo_cnt;
    wire [10:0]  mo_minx, mo_maxx;
    wire [9:0]   mo_miny, mo_maxy;

    detect_top #(
        .IMG_WIDTH(IMG_WIDTH), .IMG_HEIGHT(IMG_HEIGHT), .MIN_PIXELS(MIN_PIXELS)
    ) u_det (
        .I_clk(I_pix_clk), .I_rst_n(pix_rst_n_s),
        .I_tdata(I_tdata), .I_tvalid(I_tvalid), .I_tuser(I_tuser), .I_tlast(I_tlast),
        .O_tready(O_tready),
        .O_tdata(O_tdata), .O_tvalid(O_tvalid), .O_tuser(O_tuser), .O_tlast(O_tlast),
        .I_preset_idx(I_preset_idx),
        .O_target_valid(t_valid), .O_cx(t_cx), .O_cy(t_cy), .O_count(t_cnt),
        .O_min_x(t_minx), .O_max_x(t_maxx), .O_min_y(t_miny), .O_max_y(t_maxy),
        .O_lost_cnt(lost_det), .O_results_upd(t_upd), .O_frame_tick(t_tick),
        .O_fid(t_fid),
        .O_mot_valid(mo_valid), .O_mot_cnt(mo_cnt),
        .O_mot_minx(mo_minx), .O_mot_maxx(mo_maxx),
        .O_mot_miny(mo_miny), .O_mot_maxy(mo_maxy),
        .O_dbg_mask_valid(), .O_dbg_mask(),
        .O_s_valid(det_s_valid), .O_s_sof(det_s_sof),
        .O_s_eol(det_s_eol), .O_s_data(det_s_data),
        .O_sync_err(det_sync_err)
    );

    // ---------------- 进阶①：边缘检测并行分支 ----------------
    wire        e_tlast, e_tuser, e_tvalid;
    wire [95:0] e_tdata;
    wire [7:0]  gimbal_status;     // gimbal 原始状态字节（[6] 由顶层并入同步错误）

    edge_detect_96 #(
        .IMG_WIDTH(IMG_WIDTH), .IMG_HEIGHT(IMG_HEIGHT), .EDGE_TH(80)
    ) u_edge (
        .I_clk(I_pix_clk), .I_rst_n(pix_rst_n_s),
        .I_tlast(I_tlast), .I_tuser(I_tuser),
        .I_tdata(I_tdata), .I_tvalid(I_tvalid), .I_tready(I_tvalid),
        .O_tlast(e_tlast), .O_tuser(e_tuser), .O_tdata(e_tdata),
        .O_tvalid(e_tvalid), .O_tready()
    );

    wire        e_s_valid, e_s_sof, e_s_eol, e_sync_err;
    wire [23:0] e_s_data;
    ser96to24 u_ser_edge (
        .I_clk(I_pix_clk), .I_rst_n(pix_rst_n_s),
        .I_tdata(e_tdata), .I_tvalid(e_tvalid), .I_tuser(e_tuser), .I_tlast(e_tlast),
        .O_tready(),
        .O_valid(e_s_valid), .O_sof(e_s_sof), .O_eol(e_s_eol), .O_data(e_s_data),
        .O_sync_err(e_sync_err)
    );

    // algo mux：输出侧帧边界提交（三版外部评审共识修复）
    // 原实现按"输入 SOF"切换：mux 立即换选择，而被选分支仍在输出其上一帧
    // 尾部，造成切换帧内新旧算法内容拼接。现改为：输入 SOF 处请求，
    // 等**目标分支自己的输出 SOF** 到达时才换装——每个输出帧严格单算法，
    // 旧算法帧尾被截短（下游由新帧 SOF 复位，无污染）。
    reg        algo_q;       // 当前生效算法
    reg        algo_req;     // 有未决切换
    reg        algo_pend;    // 目标算法
    reg        algo_req_d;   // 请求延迟一拍（超时判定用）
    // 提交/超时条件与选择信号（外脑评审 GPT-3 实锤修复）：
    // 原实现里 mux 由独立 always 在"提交沿之后"才看到新 algo_q，
    // 而目标分支的 SOF 仅 1 拍脉冲 -> 新帧 SOF 被吃掉（切换窗 sof=0，
    // tb_algo_switch 已在修复前复现）。现在让提交沿**当拍**直接控制选择。
    wire algo_commit = algo_req && (algo_pend ? e_s_sof  : det_s_sof);
    wire algo_force  = algo_req && algo_req_d && I_tvalid && I_tuser;
    wire algo_sel    = (algo_commit || algo_force) ? algo_pend : algo_q;

    always @(posedge I_pix_clk or negedge pix_rst_n_s) begin
        if (!pix_rst_n_s) begin
            algo_q <= 1'b0; algo_req <= 1'b0; algo_pend <= 1'b0; algo_req_d <= 1'b0;
        end else begin
            algo_req_d <= algo_req;
            if (I_tvalid && I_tuser && (I_algo_edge != algo_q) && !algo_req) begin
                algo_req  <= 1'b1;
                algo_pend <= I_algo_edge;
            end else if (algo_commit || algo_force) begin
                // 正常：目标分支输出 SOF 当拍提交；
                // 超时：请求挂满一帧仍未提交（目标分支异常）-> 输入帧首强制提交，
                // 退化为输入侧切换语义（仅损失该帧的无拼接保证，绝不卡死）。
                algo_q   <= algo_pend;
                algo_req <= 1'b0;
            end
        end
    end

    reg        m_valid, m_sof, m_eol;
    reg [23:0] m_data;
    always @(posedge I_pix_clk or negedge pix_rst_n_s) begin
        if (!pix_rst_n_s) begin
            m_valid <= 0; m_sof <= 0; m_eol <= 0; m_data <= 0;
        end else begin
            // 注意：这里用组合的 algo_sel（含提交沿当拍）而非 algo_q，
            // 保证提交拍就选出新分支的 SOF/首像素（见上方修复注释）。
            if (!algo_sel) begin
                m_valid <= det_s_valid; m_sof <= det_s_sof;
                m_eol <= det_s_eol;     m_data <= det_s_data;
            end else begin
                m_valid <= e_s_valid;   m_sof <= e_s_sof;
                m_eol <= e_s_eol;       m_data <= e_s_data;
            end
        end
    end

    // ---------------- 扩展④：亮度/对比度（显示路径） ----------------
    // 外脑评审修复（显示链断接）：原实现 OSD 底图接 det_s_*，导致
    // 算法切换（edge）与亮度/对比度只进 video_adj 后再无下游（死分支）、
    // 画面上永远只有跟踪视频。现改为 mux -> video_adj -> OSD 完整串联；
    // sof/eol 与 video_adj 同拍延迟对齐。
    wire        a_valid;
    wire [23:0] a_data;
    reg         a_sof, a_eol;
    video_adj u_adj (
        .I_clk(I_pix_clk), .I_rst_n(pix_rst_n_s),
        .I_valid(m_valid), .I_data(m_data),
        .I_bright(I_bright), .I_contrast(I_contrast),
        .O_valid(a_valid), .O_data(a_data)
    );
    always @(posedge I_pix_clk or negedge pix_rst_n_s) begin
        if (!pix_rst_n_s) begin a_sof <= 1'b0; a_eol <= 1'b0; end
        else begin        a_sof <= m_sof;      a_eol <= m_eol; end
    end

    // ---------------- 目标包 v2（135bit，bit0 保留位，发送端必须补齐。
    // 缺位时整包上移一位，OSD 解析 valid 恒 0（"永远 LOST"事故根因，2026-09-20 修） ----
    // 外脑评审: algo 位用帧首锁存的 algo_q（与该帧实际使用的算法一致）
    wire [134:0] pack_vision = {t_valid, algo_q, I_preset_idx, t_cx, t_cy,
                                t_minx, t_maxx, t_miny, t_maxy, t_fid,
                                mo_valid, mo_cnt[13:0],
                                mo_minx, mo_maxx, mo_miny, mo_maxy, 1'b0};

    // ---------------- OSD 路径 ----------------
    osd_track #(
        .IMG_WIDTH(IMG_WIDTH), .IMG_HEIGHT(IMG_HEIGHT)
    ) u_osd (
        .I_clk(I_pix_clk), .I_rst_n(pix_rst_n_s),
        .I_valid(a_valid), .I_sof(a_sof),
        .I_eol(a_eol), .I_data(a_data),
        .I_det_clk(I_pix_clk), .I_det_rst_n(pix_rst_n_s),
        .I_det_pack(pack_vision), .I_det_load(t_upd), .O_det_busy(),
        .I_bright(I_bright), .I_contrast(I_contrast),
        .O_valid(O_osd_valid), .O_sof(O_osd_sof), .O_eol(O_osd_eol),
        .O_data(O_osd_data)
    );

    // ---------------- 声光反馈 ----------------
    // 200ms @74.25MHz = 14,850,000（原 100000 仅 1.35ms，与文档"200ms"不符）
    led_buzzer #(.BEEP_CYCLES(14850000)) u_led (
        .I_clk(I_pix_clk), .I_rst_n(pix_rst_n_s),
        .I_frame_tick(t_tick), .I_target_valid(t_valid),
        .O_led_g(O_led_g), .O_led_r(O_led_r), .O_buzzer(O_buzzer)
    );

    // ---------------- 云台控制路径 ----------------
    // 外脑评审 D1 实锤：原 tick=pulse_sync(t_tick) 比 cdc 数据早到约 0.9µs，
    // 云台在 tick 拍拿到的是上一帧甚至更旧的包（实际 2 帧滞后）。
    // 现改为"新包到达的下一拍"做节拍——tick 与数据同源，滞后=1 帧。
    wire [22:0] ctrl_pack_src = {t_valid, t_cx, t_cy};
    wire [22:0] ctrl_pack_dst;
    wire        ctrl_valid_dst;
    reg  [22:0] ctrl_hold;
    reg         ctrl_vld_hold;
    cdc_4phase #(.WIDTH(23)) u_ctrl_cdc (
        .I_clk_src(I_pix_clk), .I_rst_src_n(pix_rst_n_s),
        .I_src_data(ctrl_pack_src), .I_src_load(t_upd), .O_src_busy(),
        .I_clk_dst(I_pwm_clk), .I_rst_dst_n(pwm_rst_n_s),
        .O_dst_data(ctrl_pack_dst), .O_dst_valid(ctrl_valid_dst)
    );
    // 控制节拍 = ctrl_valid_dst 打一拍（保证本拍消费的就是刚换装的包）
    reg tick_pwm_r;
    always @(posedge I_pwm_clk or negedge pwm_rst_n_s) begin
        if (!pwm_rst_n_s) tick_pwm_r <= 1'b0;
        else              tick_pwm_r <= ctrl_valid_dst;
    end
    wire tick_pwm = tick_pwm_r;

    always @(posedge I_pwm_clk or negedge pwm_rst_n_s) begin
        if (!pwm_rst_n_s) begin
            ctrl_hold <= 23'd0; ctrl_vld_hold <= 1'b0;
        end else begin
            // 外脑评审: 新包与 tick 同拍时新包优先（原写法后清零覆盖新包,
            // 云台会丢一帧目标）
            if (ctrl_valid_dst) begin
                ctrl_hold     <= ctrl_pack_dst;
                ctrl_vld_hold <= ctrl_pack_dst[22];
            end else if (tick_pwm) begin
                ctrl_vld_hold <= 1'b0;
            end
        end
    end

    wire [12:0] duty_x, duty_y;
    gimbal_ctrl #(
        .KP_X(KP_X), .KP_Y(KP_Y), .DZ_X(DZ_X), .DZ_Y(DZ_Y),
        .SLEW_US(SLEW_US), .SERVO_MIN_US(SERVO_MIN_US), .SERVO_MAX_US(SERVO_MAX_US),
        .CENTER_X(IMG_WIDTH/2), .CENTER_Y(IMG_HEIGHT/2),
        .INVERT_X(INVERT_X), .INVERT_Y(INVERT_Y)
    ) u_gimbal (
        .I_clk(I_pwm_clk), .I_rst_n(pwm_rst_n_s),
        .I_frame_tick(tick_pwm), .I_track_en(I_track_en),
        .I_target_valid(ctrl_vld_hold),
        .I_cx(ctrl_hold[21:11]), .I_cy(ctrl_hold[10:0]),
        .O_duty_x_us(duty_x), .O_duty_y_us(duty_y),
        .O_lost_cnt(O_lost_cnt), .O_status(gimbal_status)
    );

    // O_status 位定义：gimbal [5:0]（dz/lim/lost）+ [6] 串化器帧同步错误（v4.1 新）
    //                    + [7] 保留（0）。位 6 为"帧同步契约破坏"粘滞指示。
    assign O_status = {1'b0, (det_sync_err | e_sync_err), gimbal_status[5:0]};

    pwm_servo u_pwm (
        .I_clk(I_pwm_clk), .I_rst_n(pwm_rst_n_s),
        .I_update(tick_pwm),
        .I_duty_x_us(duty_x), .I_duty_y_us(duty_y),
        .O_pwm_x(O_pwm_x), .O_pwm_y(O_pwm_y),
        .O_period_tick(), .O_duty_x_dbg(O_duty_x_us), .O_duty_y_dbg(O_duty_y_us)
    );

    // ---------------- 目标状态直出 ----------------
    assign O_frame_tick   = t_tick;
    assign O_target_valid = t_valid;
    assign O_tgt_cx       = t_cx;
    assign O_tgt_cy       = t_cy;

endmodule
\end{Verbatim}
\subsection{video\_adj.v}
\fileintro{亮度/对比度（1 拍流水）：out = clamp(((c-128)xC>>6)+128+B)。C 为 Q6（64=1.0x）；B 为 7bit 有符号——修复记：原 {1'b0,bright} 零扩展把负亮度当正数，已改 \$signed(bright) 并手工验证（-16->-16）。}
\begin{Verbatim}[numbers=left,numbersep=5pt,frame=single,fontsize=\footnotesize,breaklines=true,breakanywhere=true]
// =============================================================================
// video_adj.v —— 显示参数调节（扩展④）：亮度/对比度
// out = clamp( ((c - 128) * C >> 6) + 128 + B )
//   C: 对比度 64..192（64 = 1.0 倍，Q6 定点）；B: 亮度 -64..+64
// 应用点：显示路径（OSD 之前），检测分支不受影响（扩展④"调节显示参数"）。
// 1 拍流水；参数同域直连（vision_top 内 pix 域）。
// =============================================================================
module video_adj (
    input  wire        I_clk,
    input  wire        I_rst_n,
    input  wire        I_valid,
    input  wire [23:0] I_data,        // {R,G,B}
    input  wire [6:0]  I_bright,      // 有符号：-64..+63（二补码）
    input  wire [7:0]  I_contrast,    // 无符号：64..192（Q6）
    output reg         O_valid,
    output reg  [23:0] O_data
);
    // 单通道处理函数
    function [7:0] adj_ch;
        input [7:0] c;
        input [6:0] bright;
        input [7:0] contrast;
        reg signed [16:0] v;
        begin
            v = $signed({1'b0, c}) - 17'sd128;       // 居中
            v = (v * $signed({1'b0, contrast})) >>> 6;
            // 外脑评审修复: bright 是 7bit 有符号(-64..+63), 原 {1'b0,bright}
            // 零扩展会把负数当正数(-64 -> +64); 用 $signed 直接符号扩展
            v = v + 17'sd128 + $signed(bright);
            if (v < 0)        adj_ch = 8'd0;
            else if (v > 255) adj_ch = 8'd255;
            else              adj_ch = v[7:0];
        end
    endfunction

    wire [7:0] r_adj = adj_ch(I_data[23:16], I_bright, I_contrast);
    wire [7:0] g_adj = adj_ch(I_data[15:8],  I_bright, I_contrast);
    wire [7:0] b_adj = adj_ch(I_data[7:0],   I_bright, I_contrast);

    always @(posedge I_clk or negedge I_rst_n) begin
        if (!I_rst_n) begin
            O_valid <= 1'b0;
            O_data  <= 24'd0;
        end else begin
            O_valid <= I_valid;
            O_data  <= {r_adj, g_adj, b_adj};
        end
    end
endmodule
\end{Verbatim}
\subsection{motion\_detect.v}
\fileintro{帧差运动检测 v3.1：整帧 1bit 前帧缓存（115KB ERAM）。同步读一级流水（TD 可实现），写地址延迟一拍与读结构性错开（与 ERAM 写模式解耦——防 WRITE\_FIRST 下 diff 恒 0）；首帧 prev\_ready 门控；addr 越界保护；cnt==0 时 bbox 全 0（与金标准一致）。每帧末锁存 cnt/bbox + 1 拍选通，亚门限帧照常更新（valid=cnt>=MOT\_TH）。}
\begin{Verbatim}[numbers=left,numbersep=5pt,frame=single,fontsize=\footnotesize,breaklines=true,breakanywhere=true]
// =============================================================================
// motion_detect.v —— 帧差运动检测（进阶④，v3 流水版 2026-10-03）
// 输入：平滑后的二值掩码流（与 centroid 同源同拍）。
// 结构：整帧 1bit 前帧缓存（W*H bit，TD 映射 ERAM），同步读一级流水：
//       第 1 拍发地址/写新值，第 2 拍取读出旧值与延迟一拍的掩码比对。
//       —— v2 用组合读 prev_mem[addr]，行为仿真可行但 TD ERAM(Block RAM)
//          不支持组合读，上板综合会失败；v3 改为同步读 + 信号对齐流水。
// 输出（语义定稿，2026-10-03 评审措辞澄清）：
//   O_mot_cnt/O_mot_bbox 每帧末锁存（含亚门限帧，cnt==0 时 bbox 全 0）；
//   O_mot_valid **本身就是阈值标志**=(cnt >= MOT_TH)：仅该帧判为运动时
//   拉高 1 拍（帧末结果与阈值判决合为一个输出，无独立逐帧选通信号）；
//   OSD 运动框以它为门控（mv），亚门限帧不画框。
//
// 版本要点：
//   v2 [1] sof 拍写地址强制 0（不再用上一帧遗留 row_base 算地址）；
//      [2] 首帧门控用 prev_ready（首帧不比对、不输出）；
//   v3 [3] 同步读流水（TD 可实现）；[4] addr 越界保护；
//   v3.1 写地址延迟一拍（与 ERAM 写模式解耦，防 WRITE_FIRST 下 diff 恒 0）；
//      [5] cnt==0 时 bbox 锁存 0（与金标准 frame_diff_motion 语义一致）。
// 已知约定：(0,0) 像素每帧只写缓存、不参与比对（sof 拍结构，1/921600）。
// =============================================================================
module motion_detect #(
    parameter IMG_WIDTH  = 1280,
    parameter IMG_HEIGHT = 720,
    parameter MOT_TH     = 64
)(
    input  wire        I_clk,
    input  wire        I_rst_n,
    input  wire        I_valid,
    input  wire        I_sof,
    input  wire        I_eol,
    input  wire        I_mask,
    output reg         O_mot_valid,
    output reg  [13:0] O_mot_cnt,
    output reg  [10:0] O_mot_minx,
    output reg  [10:0] O_mot_maxx,
    output reg  [9:0]  O_mot_miny,
    output reg  [9:0]  O_mot_maxy
);
    localparam TOTAL = IMG_WIDTH * IMG_HEIGHT;

    (* ram_style = "block" *) reg prev_mem [0:TOTAL-1];

    // ---------------- 坐标计数（输入拍） ----------------
    reg [10:0] x;
    reg [9:0]  y;
    reg [20:0] row_base;
    reg        prev_ready;   // 首帧缓存填满后置位；首帧不比对、不输出

    // sof 拍像素是 (0,0)：写地址强制 0（见 v2[1]）
    wire [20:0] addr = I_sof ? 21'd0 : (row_base + {10'd0, x});
    wire [20:0] addr_safe = (addr < TOTAL) ? addr : 21'd0;   // v3[4] 越界保护

    always @(posedge I_clk or negedge I_rst_n) begin
        if (!I_rst_n) begin
            x <= 11'd0; y <= 10'd0; row_base <= 21'd0; prev_ready <= 1'b0;
        end else if (I_valid) begin
            if (I_sof) begin
                x <= 11'd1; y <= 10'd0; row_base <= 21'd0;
            end else if (I_eol) begin
                x <= 11'd0;
                y <= (y == IMG_HEIGHT - 1) ? 10'd0 : y + 10'd1;
                row_base <= row_base + IMG_WIDTH;
                if (y == IMG_HEIGHT - 1) prev_ready <= 1'b1;  // 首帧末行：缓存已填满
            end else begin
                x <= x + 11'd1;
            end
        end
    end

    // ---------------- 一级流水：ERAM 同步读 + 信号对齐（v3[3]） ----------------
    // v3.1 外脑评审：写地址**延迟一拍**——杜绝同拍同址读写。
    // 若 ERAM 被综合为 WRITE_FIRST 模式，同址读写会让 old_r 直接读到新掩码，
    // diff 恒 0（运动检测静默失效）；地址错开一拍后与写模式无关（结构性读旧）。
    reg        old_r;                          // 前帧掩码（同步读输出）
    reg [20:0] addr_d;                         // 写地址（延迟一拍）
    reg        mask_wr;                        // 写数据（延迟一拍）
    reg        wr_en_d;                        // 写使能（延迟一拍）
    reg        v_d, sof_d, eol_d, mask_d, ready_d;
    reg [10:0] x_d;
    reg [9:0]  y_d;

    always @(posedge I_clk or negedge I_rst_n) begin
        if (!I_rst_n) begin
            old_r <= 1'b0;
            addr_d <= 21'd0; mask_wr <= 1'b0; wr_en_d <= 1'b0;
            v_d <= 1'b0; sof_d <= 1'b0; eol_d <= 1'b0;
            mask_d <= 1'b0; ready_d <= 1'b0;
            x_d <= 11'd0; y_d <= 10'd0;
        end else begin
            old_r <= prev_mem[addr_safe];              // 读：本拍地址（读旧）
            // 写：上一拍的地址/数据 -> 与读地址恒不同址（流内地址逐一递增）
            if (wr_en_d && (addr_d < TOTAL))
                prev_mem[addr_d] <= mask_wr;
            addr_d  <= addr_safe;
            mask_wr <= I_mask;
            wr_en_d <= I_valid;
            v_d   <= I_valid; sof_d <= I_sof; eol_d <= I_eol;
            mask_d <= I_mask; ready_d <= prev_ready;
            x_d   <= x; y_d <= y;
        end
    end

    wire diff = mask_d ^ old_r;    // 与 x_d/y_d/v_d 同拍对齐

    // ---------------- 帧内累加（延迟域） ----------------
    reg [13:0] cnt;
    reg        seen;
    reg [10:0] bmin_x, bmax_x;
    reg [9:0]  bmin_y, bmax_y;

    always @(posedge I_clk or negedge I_rst_n) begin
        if (!I_rst_n) begin
            cnt <= 14'd0; seen <= 1'b0;
            bmin_x <= 11'd0; bmax_x <= 11'd0; bmin_y <= 10'd0; bmax_y <= 10'd0;
        end else if (v_d && sof_d) begin
            cnt <= 14'd0; seen <= 1'b0;            // 帧首清零（sof 像素不比对）
        end else if (v_d && ready_d && diff) begin
            cnt <= (cnt == 14'h3FFF) ? cnt : cnt + 14'd1;   // 饱和（14bit 接口）
            if (!seen) begin
                seen   <= 1'b1;
                bmin_x <= x_d; bmax_x <= x_d; bmin_y <= y_d; bmax_y <= y_d;
            end else begin
                if (x_d < bmin_x) bmin_x <= x_d;
                if (x_d > bmax_x) bmax_x <= x_d;
                if (y_d < bmin_y) bmin_y <= y_d;
                if (y_d > bmax_y) bmax_y <= y_d;
            end
        end
    end

    // ---------------- 帧末锁存（延迟域；首帧 prev_ready=0 不输出） ----------------
    // 外脑评审修复：锁存再延迟一拍——末像素的 diff 在 last_eol_d 同拍累加，
    // 同拍锁存读到的是旧 cnt/bbox（每帧漏掉 (W-1,H-1) 一个像素的计数）。
    wire last_eol_d = v_d && eol_d && (y_d == IMG_HEIGHT - 1);
    reg  last_eol_dd;
    always @(posedge I_clk or negedge I_rst_n) begin
        if (!I_rst_n) last_eol_dd <= 1'b0;
        else          last_eol_dd <= last_eol_d;
    end

    always @(posedge I_clk or negedge I_rst_n) begin
        if (!I_rst_n) begin
            O_mot_valid <= 1'b0;
            O_mot_cnt <= 14'd0;
            O_mot_minx <= 11'd0; O_mot_maxx <= 11'd0;
            O_mot_miny <= 10'd0; O_mot_maxy <= 10'd0;
        end else begin
            O_mot_valid <= 1'b0;
            if (last_eol_dd && ready_d) begin
                O_mot_valid <= (cnt >= MOT_TH);
                O_mot_cnt   <= cnt;
                // v3[5] cnt==0 时 bbox 全 0（与金标准一致）
                O_mot_minx  <= (cnt == 14'd0) ? 11'd0 : bmin_x;
                O_mot_maxx  <= (cnt == 14'd0) ? 11'd0 : bmax_x;
                O_mot_miny  <= (cnt == 14'd0) ? 10'd0 : bmin_y;
                O_mot_maxy  <= (cnt == 14'd0) ? 10'd0 : bmax_y;
            end
        end
    end

endmodule
\end{Verbatim}
\subsection{led\_buzzer.v}
\fileintro{声光反馈：锁定绿灯、丢失红灯 1Hz 闪、锁定沿蜂鸣 200ms（BEEP\_CYCLES 参数化）。}
\begin{Verbatim}[numbers=left,numbersep=5pt,frame=single,fontsize=\footnotesize,breaklines=true,breakanywhere=true]
// =============================================================================
// led_buzzer.v —— 目标状态声光反馈（扩展"外部控制"演示）
// 锁定：绿灯亮；丢失：红灯亮（1Hz 闪烁提示）；锁定沿：蜂鸣器短鸣 200ms。
// 板载双色 LED / 蜂鸣器为 GPIO 直驱（若为有源蜂鸣器低有效，接 BEEP_N 低有效参数）。
// =============================================================================
module led_buzzer #(
    parameter integer BEEP_CYCLES = 20000000   // 200ms @100MHz，参数化
)(
    input  wire I_clk,
    input  wire I_rst_n,
    input  wire I_frame_tick,       // 每帧一拍
    input  wire I_target_valid,     // 本帧检测有效
    output reg  O_led_g,            // 锁定绿灯
    output reg  O_led_r,            // 丢失红灯
    output reg  O_buzzer            // 蜂鸣器（高有效脉冲）
);
    reg [25:0] beep_cnt;
    reg        valid_d;
    reg [25:0] blink_cnt;

    always @(posedge I_clk or negedge I_rst_n) begin
        if (!I_rst_n) begin
            O_led_g <= 1'b0; O_led_r <= 1'b0; O_buzzer <= 1'b0;
            valid_d <= 1'b0; beep_cnt <= 0; blink_cnt <= 0;
        end else begin
            valid_d <= I_target_valid;

            // 锁定沿：蜂鸣 200ms
            if (I_frame_tick && I_target_valid && !valid_d)
                beep_cnt <= BEEP_CYCLES[25:0];
            if (beep_cnt != 0) begin
                O_buzzer <= 1'b1;
                beep_cnt <= beep_cnt - 26'd1;
            end else begin
                O_buzzer <= 1'b0;
            end

            // LED：锁定常亮绿；丢失红灯 1Hz 闪烁（远看醒目）
            blink_cnt <= blink_cnt + 26'd1;
            O_led_g <= I_target_valid;
            if (I_target_valid)
                O_led_r <= 1'b0;
            else
                O_led_r <= blink_cnt[24];   // ~0.67s 半周期 @100MHz
        end
    end
endmodule
\end{Verbatim}
\subsection{hsv\_inv\_lut.v（生成文件 · 节选）}
\fileintro{倒数表 ROM（生成文件，禁止手改）：inv[d]=floor(3840/d)，256x12bit，由 gen\_vectors.py 生成 case 表（综合为 Block ROM）。}
\begin{Verbatim}[numbers=left,numbersep=5pt,frame=single,fontsize=\footnotesize,breaklines=true,breakanywhere=true]
// 自动生成：tools/gen_vectors.py  inv[d]=floor(3840/d) d=1..255
// 供 rgb2hsv.v 求色调用，综合为 Block ROM
module hsv_inv_lut (
    input  wire [7:0]  I_d,
    output reg  [11:0] O_inv
);
always @(*) begin
    case (I_d)
        8'd1: O_inv = 12'd3840;
        8'd2: O_inv = 12'd1920;
        8'd3: O_inv = 12'd1280;
        8'd4: O_inv = 12'd960;
        8'd5: O_inv = 12'd768;
        8'd6: O_inv = 12'd640;
        8'd7: O_inv = 12'd548;
        8'd8: O_inv = 12'd480;
        8'd9: O_inv = 12'd426;
        8'd10: O_inv = 12'd384;
        8'd11: O_inv = 12'd349;
        8'd12: O_inv = 12'd320;
        8'd13: O_inv = 12'd295;
        8'd14: O_inv = 12'd274;
        8'd15: O_inv = 12'd256;
        8'd16: O_inv = 12'd240;
\end{Verbatim}
\begin{notebox}{节选说明}
{\small 该文件共 268 行（完整内容由 gencode 脚本可复现生成）。}
\end{notebox}
\subsection{font\_rom\_8x16.v（生成文件 · 节选）}
\fileintro{8x16 点阵字库（生成文件，禁止手改）：40 字符 ASCII 子集，由 gen\_vectors.py 从 cv2 Hershey 字体渲染生成；地址 {char[7:0], row[3:0]}，MSB=最左像素。}
\begin{Verbatim}[numbers=left,numbersep=5pt,frame=single,fontsize=\footnotesize,breaklines=true,breakanywhere=true]
// 自动生成：tools/gen_vectors.py (cv2 Hershey->8x16 点阵)
// 地址 {char[7:0], row[3:0]} -> 该行 8 位点阵（MSB=最左像素）
module font_rom_8x16 (
    input  wire [11:0] I_addr,
    output reg  [7:0]  O_bits
);
always @(*) begin
    case (I_addr)
        12'h200: O_bits = 8'h00;  // ' ' row0
        12'h201: O_bits = 8'h00;  // ' ' row1
        12'h202: O_bits = 8'h00;  // ' ' row2
        12'h203: O_bits = 8'h00;  // ' ' row3
        12'h204: O_bits = 8'h00;  // ' ' row4
        12'h205: O_bits = 8'h00;  // ' ' row5
        12'h206: O_bits = 8'h00;  // ' ' row6
        12'h207: O_bits = 8'h00;  // ' ' row7
        12'h208: O_bits = 8'h00;  // ' ' row8
        12'h209: O_bits = 8'h00;  // ' ' row9
\end{Verbatim}
\begin{notebox}{节选说明}
{\small 该文件共 685 行（完整内容由 gencode 脚本可复现生成）。}
\end{notebox}
\subsection{marquee\_rom\_64.v（生成文件 · 节选）}
\fileintro{跑马灯字库（生成文件，禁止手改）：63 字符消息 x 16 行 case 表，由 gen\_vectors.py 生成。}
\begin{Verbatim}[numbers=left,numbersep=5pt,frame=single,fontsize=\footnotesize,breaklines=true,breakanywhere=true]
// 自动生成：tools/gen_marquee_rom.py（由 case 版逐值转为数组 ROM，位值不变）
// 地址 {char_idx[5:0], row[3:0]} -> 该行 8 位点阵（MSB=最左像素）
// 2026-10-03 转换缘由：case 版被 TD 映射为 ~10 级 mux 树，是 74.25MHz 域
//   WNS -0.113ns 的唯一关键路径（gate 报告 8 个违规端点全部同源）；
//   数组 ROM 异步读映射为 LUTRAM（1-2 级），功能零变化（tb_osd_track 逐像素复验）。
module marquee_rom_64 (
    input  wire [9:0]  I_addr,
    output wire [7:0]  O_bits
);

    (* rom_style = "distributed" *) reg [7:0] mem [0:1023];

    initial begin
        mem[10'h000] = 8'h00;  //'A' row0
        mem[10'h001] = 8'h00;  //'A' row1
        mem[10'h002] = 8'h00;  //'A' row2
        mem[10'h003] = 8'h0E;  //'A' row3
        mem[10'h004] = 8'h0E;  //'A' row4
\end{Verbatim}
\begin{notebox}{节选说明}
{\small 该文件共 1028 行（完整内容由 gencode 脚本可复现生成）。}
\end{notebox}
\section{视觉链路验证 TB（18 项回归全集）}
\subsection{tb\_ser96to24.v}
\begin{Verbatim}[numbers=left,numbersep=5pt,frame=single,fontsize=\footnotesize,breaklines=true,breakanywhere=true]
// =============================================================================
// tb_ser96to24.v —— 96->24 串化器 TB：随机空拍 + 帧边界标志 + 队列自检
// 验收：输入侧 4 像素/拍全展开，sof/eol 标志逐像素正确，无丢失/乱序。
// 阶段 2（v4.1 加固验证，2026-10-03 第 7 轮复审）：
//   故障注入——连续 back-to-back 入队制造溢出 + 全程 tuser 缺失（帧同步契约破坏），
//   验证：看门狗置 O_sync_err；tlast 门控的合成帧首恢复入队（恰 1 个 sof）；
//   恢复后普通 beat 持续可写（drop 停止增长）。SYNC_TIMEOUT 缩短为 2000 便于仿真。
// =============================================================================
`timescale 1ns/1ps
module tb_ser96to24;
    localparam W = 64, H = 8, BEATS = W/4, FRAME_BEATS = BEATS*H, FRAMES = 3;

    reg clk = 0, rst_n = 0;
    always #7 clk = ~clk;

    reg [95:0] tdata;
    reg        tvalid, tuser, tlast;
    wire       tready;
    wire       o_valid, o_sof, o_eol;
    wire [23:0] o_data;
    wire       o_sync_err;

    ser96to24 #(.SYNC_TIMEOUT(24'd2000)) dut (.I_clk(clk), .I_rst_n(rst_n),
        .I_tdata(tdata), .I_tvalid(tvalid), .I_tuser(tuser), .I_tlast(tlast),
        .O_tready(tready),
        .O_valid(o_valid), .O_sof(o_sof), .O_eol(o_eol), .O_data(o_data),
        .O_sync_err(o_sync_err));

    // 期望队列
    reg [23:0] exp_d [0:FRAMES*W*H-1];
    reg        exp_s [0:FRAMES*W*H-1];
    reg        exp_e [0:FRAMES*W*H-1];
    integer    push_n = 0, pop_n = 0, mism = 0;
    integer    f, b, k, i;
    reg [15:0] lfsr = 16'hBEEF;
    reg        drv_valid;

    // 阶段 2（故障注入）观测
    reg        phase2 = 0;
    integer    p2_pop = 0, p2_sof = 0, p2_mism = 0, p2_after_sof = 0;
    reg        p2_empty_seen = 0;
    reg        p2_busy_after = 0;

    // 输入驱动
    initial begin
        tdata = 0; tvalid = 0; tuser = 0; tlast = 0;
        repeat (4) @(posedge clk);
        rst_n <= 1;
        for (f = 0; f < FRAMES; f = f + 1) begin
            for (b = 0; b < FRAME_BEATS; b = b + 1) begin
                // 总线契约：beat 率 = 像素率/4（4 拍 1 beat），另有少量随机额外空闲
                repeat (3) begin
                    @(posedge clk);
                    tvalid <= 1'b0;
                end
                lfsr = {lfsr[14:0], lfsr[15]^lfsr[13]^lfsr[12]^lfsr[10]};
                drv_valid = (lfsr[2:0] != 3'b000);
                while (!drv_valid) begin
                    @(posedge clk);
                    tvalid <= 1'b0;
                    lfsr = {lfsr[14:0], lfsr[15]^lfsr[13]^lfsr[12]^lfsr[10]};
                    drv_valid = (lfsr[2:0] != 3'b000);
                end
                @(posedge clk);
                tvalid <= 1'b1;
                for (k = 0; k < 4; k = k + 1) begin
                    tdata[95-24*k -: 24] <= 24'h100000 * f + b * 4 + k;
                    exp_d[push_n + k] = 24'h100000 * f + b * 4 + k;
                    exp_s[push_n + k] = (b == 0) && (k == 0);
                    exp_e[push_n + k] = (k == 3) && ((b % BEATS) == BEATS-1);
                end
                tuser <= (b == 0);
                tlast <= ((b % BEATS) == BEATS-1);
                push_n = push_n + 4;
            end
        end
        @(posedge clk);
        tvalid <= 1'b0; tuser <= 1'b0; tlast <= 1'b0;
        repeat (64) @(posedge clk);

        // ---------------- 阶段 2：故障注入（溢出 + tuser 缺失 + 看门狗恢复） ----------------
        phase2 = 1;
        // 2a: 连续 back-to-back 入队至溢出（960 beat @1/clk -> 净积压 2880 > 2048），
        //     全程 tuser=0（帧同步契约破坏）；tlast 保留行边界
        for (i = 0; i < 960; i = i + 1) begin
            @(posedge clk);
            tdata  <= 96'h0;
            tvalid <= 1'b1; tuser <= 1'b0; tlast <= ((i % 320) == 319);
        end
        // 2b: 无 tuser 继续喂 > SYNC_TIMEOUT(2000) beat —— 看门狗应触发
        for (i = 0; i < 2600; i = i + 1) begin
            @(posedge clk);
            tvalid <= 1'b1; tuser <= 1'b0; tlast <= 1'b0;
        end
        @(posedge clk);
        if (!o_sync_err) begin p2_mism = p2_mism + 1; $display("MISM P2: sync_err 未置位"); end
        else $display("P2: 看门狗触发 sync_err=1；缓冲已排空 o_valid=%b", o_valid);
        if (o_valid) p2_busy_after = 1;     // 此刻缓冲应已排空（o_valid 应为 0）
        // 2c: 给一个 tlast 拍 -> 合成帧首恢复（仍无 tuser）
        @(posedge clk);
        tdata  <= {24'hB00001, 24'hB00002, 24'hB00003, 24'hB00004};
        tvalid <= 1'b1; tuser <= 1'b0; tlast <= 1'b1;
        // 2d: 恢复后继续喂普通 beat（仍无 tuser），验证持续可写
        for (i = 0; i < 16; i = i + 1) begin
            @(posedge clk);
            tdata  <= {24'hB00005, 24'hB00006, 24'hB00007, 24'hB00008};
            tvalid <= 1'b1; tuser <= 1'b0; tlast <= 1'b0;
        end
        @(posedge clk);
        tvalid <= 1'b0;
        repeat (64) @(posedge clk);
        if ((mism == 0) && (p2_mism == 0) && (pop_n == push_n) &&
            (p2_sof == 1) && (p2_after_sof == 3) && (p2_pop >= 12) &&
            (p2_empty_seen) && !p2_busy_after)
            $display("TB_RESULT: PASS (p1 pixels=%0d, p2 recovered sof=%0d pop=%0d)",
                     pop_n, p2_sof, p2_pop);
        else
            $display("TB_RESULT: FAIL p1 mism=%0d pop=%0d/%0d | p2 mism=%0d sof=%0d after_sof=%0d pop=%0d empty_seen=%b",
                     mism, pop_n, push_n, p2_mism, p2_sof, p2_after_sof, p2_pop, p2_empty_seen);
        $finish;
    end

    // 输出比对（阶段 1：逐像素期望；阶段 2：恢复段语义检查）
    always @(posedge clk) begin
        if (rst_n && o_valid) begin
            if (!phase2) begin
                if ((o_data !== exp_d[pop_n]) || (o_sof !== exp_s[pop_n]) ||
                    (o_eol !== exp_e[pop_n])) begin
                    if (mism < 8)
                        $display("MISM idx=%0d got d=%h s=%b e=%b exp d=%h s=%b e=%b",
                                 pop_n, o_data, o_sof, o_eol,
                                 exp_d[pop_n], exp_s[pop_n], exp_e[pop_n]);
                    mism = mism + 1;
                end
                pop_n = pop_n + 1;
            end else begin
                p2_pop = p2_pop + 1;
                if (o_sof) begin
                    p2_sof = p2_sof + 1;
                    if (o_data !== 24'hB00001) begin
                        p2_mism = p2_mism + 1;
                        $display("MISM P2: sof 像素数据=%h（应 B00001）", o_data);
                    end
                end else if (p2_sof > 0 && p2_after_sof < 3) begin
                    // 合成帧首之后的 3 个像素必须为 B00002..B00004
                    if (o_data !== (24'hB00002 + p2_after_sof)) begin
                        p2_mism = p2_mism + 1;
                        $display("MISM P2: 恢复段数据=%h（应 %h）",
                                 o_data, 24'hB00002 + p2_after_sof);
                    end
                    p2_after_sof = p2_after_sof + 1;
                end
            end
        end
        if (phase2 && o_sync_err && !o_valid) p2_empty_seen = 1;   // 看门狗后的排空态
    end

    initial begin
        #20000000;
        $display("TB_RESULT: FAIL timeout");
        $finish;
    end
endmodule
\end{Verbatim}
\subsection{tb\_tpg.v}
\begin{Verbatim}[numbers=left,numbersep=5pt,frame=single,fontsize=\footnotesize,breaklines=true,breakanywhere=true]
// =============================================================================
// tb_tpg.v —— 测试图案发生器 TB：几何计数自检 + 帧转储交 Python 比对
// 验收：每帧 sof 恰 1 次、eol 恰 720 次、像素数恰 921600；frame_0 交给
// Python 与 gen_vectors 生成的 tpg_exp.hex 逐像素比对（run_all.py）。
// =============================================================================
`timescale 1ns/1ps
module tb_tpg;
    reg clk = 0, rst_n = 0;
    always #7 clk = ~clk;

    reg  [10:0] box_x = 11'd500;
    reg  [9:0]  box_y = 10'd300;
    wire        tvalid, tuser, tlast;
    wire [95:0] tdata;

    tpg_720p #(.BOX_EN(1)) dut (
        .I_clk(clk), .I_rst_n(rst_n),
        .I_box_x(box_x), .I_box_y(box_y),
        .O_tvalid(tvalid), .O_tdata(tdata), .O_tuser(tuser), .O_tlast(tlast),
        .I_tready(1'b1));

    wire        s_valid, s_sof, s_eol;
    wire [23:0] s_data;
    ser96to24 u_ser (.I_clk(clk), .I_rst_n(rst_n),
        .I_tdata(tdata), .I_tvalid(tvalid), .I_tuser(tuser), .I_tlast(tlast),
        .O_tready(), .O_valid(s_valid), .O_sof(s_sof), .O_eol(s_eol), .O_data(s_data));

    frame_dump24 #(.BASE_NAME("tpg_act"), .MAX_FRAMES(1)) u_dump (
        .I_clk(clk), .I_rst_n(rst_n),
        .I_valid(s_valid), .I_eol(s_eol), .I_data(s_data), .O_frame_done());

    integer sof_cnt = 0, eol_cnt = 0, px_cnt = 0, err_cnt = 0;
    integer frames = 0;

    always @(posedge clk) begin
        if (rst_n && s_valid) begin
            px_cnt = px_cnt + 1;
            if (s_sof) begin
                sof_cnt = sof_cnt + 1;
                if (px_cnt != 1) err_cnt = err_cnt + 1;  // sof 必须是帧首像素
            end
            if (s_eol) eol_cnt = eol_cnt + 1;
            if (px_cnt == 1280*720) begin
                if ((sof_cnt != 1) || (eol_cnt != 720)) begin
                    $display("FRAME_GEO_ERR frame=%0d sof=%0d eol=%0d",
                             frames, sof_cnt, eol_cnt);
                    err_cnt = err_cnt + 1;
                end
                frames = frames + 1;
                sof_cnt = 0; eol_cnt = 0; px_cnt = 0;
                if (frames == 2) begin
                    if (err_cnt == 0)
                        $display("TB_RESULT: PASS (geometry 2 frames)");
                    else
                        $display("TB_RESULT: FAIL err_cnt=%0d", err_cnt);
                    $finish;
                end
            end
        end
    end

    initial begin
        repeat (4) @(posedge clk);
        rst_n <= 1;
        #30000000;
        $display("TB_RESULT: FAIL timeout frames=%0d", frames);
        $finish;
    end
endmodule
\end{Verbatim}
\subsection{tb\_pulse\_sync.v}
\begin{Verbatim}[numbers=left,numbersep=5pt,frame=single,fontsize=\footnotesize,breaklines=true,breakanywhere=true]
// =============================================================================
// tb_pulse_sync.v —— 脉冲跨域 TB：异步双时钟，100 个随机脉冲 1:1 传递
// =============================================================================
`timescale 1ns/1ps
module tb_pulse_sync;
    reg clk_s = 0, clk_d = 0, rst_s = 0, rst_d = 0;
    always #7  clk_s = ~clk_s;    // ~71.4MHz 源
    always #10 clk_d = ~clk_d;    // 50MHz 目标

    reg  pulse_s = 0;
    wire pulse_d;

    pulse_sync dut (
        .I_clk_src(clk_s), .I_rst_src_n(rst_s), .I_pulse_src(pulse_s),
        .I_clk_dst(clk_d), .I_rst_dst_n(rst_d), .O_pulse_dst(pulse_d));

    integer sent = 0, got = 0, i;
    reg [15:0] lfsr = 16'h1234;

    always @(posedge clk_d) begin
        if (rst_d && pulse_d)
            got = got + 1;
    end

    initial begin
        repeat (3) @(posedge clk_s);
        rst_s <= 1;
        repeat (3) @(posedge clk_d);
        rst_d <= 1;
        for (i = 0; i < 100; i = i + 1) begin
            repeat ($urandom_range(4, 20)) @(posedge clk_s);
            pulse_s <= 1'b1;
            @(posedge clk_s);
            pulse_s <= 1'b0;
            sent = sent + 1;
        end
        #5000;
        if (got == sent)
            $display("TB_RESULT: PASS (sent=%0d got=%0d)", sent, got);
        else
            $display("TB_RESULT: FAIL sent=%0d got=%0d", sent, got);
        $finish;
    end
endmodule
\end{Verbatim}
\subsection{tb\_cdc4phase.v}
\begin{Verbatim}[numbers=left,numbersep=5pt,frame=single,fontsize=\footnotesize,breaklines=true,breakanywhere=true]
// =============================================================================
// tb_cdc4phase.v —— 4 相握手数据跨域 TB：200 个随机数据完整性校验
// =============================================================================
`timescale 1ns/1ps
module tb_cdc4phase;
    localparam WIDTH = 23;
    reg clk_s = 0, clk_d = 0, rst_s = 0, rst_d = 0;
    always #7  clk_s = ~clk_s;
    always #10 clk_d = ~clk_d;

    reg  [WIDTH-1:0] s_data = 0;
    reg              s_load = 0;
    wire             s_busy;
    wire [WIDTH-1:0] d_data;
    wire             d_valid;

    cdc_4phase #(.WIDTH(WIDTH)) dut (
        .I_clk_src(clk_s), .I_rst_src_n(rst_s),
        .I_src_data(s_data), .I_src_load(s_load), .O_src_busy(s_busy),
        .I_clk_dst(clk_d), .I_rst_dst_n(rst_d),
        .O_dst_data(d_data), .O_dst_valid(d_valid));

    integer sent = 0, got = 0, mism = 0;
    reg [WIDTH-1:0] exp_q [0:255];
    reg [7:0] wp = 0, rp = 0;

    always @(posedge clk_d) begin
        if (rst_d && d_valid) begin
            if (d_data !== exp_q[rp]) begin
                if (mism < 5)
                    $display("MISM got=%h exp=%h", d_data, exp_q[rp]);
                mism = mism + 1;
            end
            rp = rp + 8'd1;
            got = got + 1;
        end
    end

    integer i;
    reg [WIDTH-1:0] data_v;
    initial begin
        repeat (3) @(posedge clk_s); rst_s <= 1;
        repeat (3) @(posedge clk_d); rst_d <= 1;
        for (i = 0; i < 200; i = i + 1) begin
            wait (!s_busy);
            repeat ($urandom_range(0, 6)) @(posedge clk_s);
            data_v = WIDTH'($urandom());
            s_data <= data_v;
            exp_q[wp] <= data_v;
            wp = wp + 8'd1;
            s_load <= 1'b1;
            @(posedge clk_s);
            s_load <= 1'b0;
            // 等 DUT 受理（busy 拉高）并完成握手（busy 拉低），避免同沿竞态
            @(posedge clk_s);
            wait (s_busy === 1'b1);
            wait (s_busy === 1'b0);
            sent = sent + 1;
        end
        #5000;
        if ((mism == 0) && (got == sent))
            $display("TB_RESULT: PASS (sent=%0d got=%0d)", sent, got);
        else
            $display("TB_RESULT: FAIL sent=%0d got=%0d mism=%0d", sent, got, mism);
        $finish;
    end
endmodule
\end{Verbatim}
\subsection{tb\_cdc\_reset\_skew.v}
\fileintro{CDC 复位偏斜验证（三版评审共识 Gate）：两域复位同拉低、错开 0/3/7/10 拍释放，静默期断言无伪 valid，随后两笔事务计数与数据逐一比对。}
\begin{Verbatim}[numbers=left,numbersep=5pt,frame=single,fontsize=\footnotesize,breaklines=true,breakanywhere=true]
`timescale 1ns/1ps
// =============================================================================
// tb_cdc_reset_skew.v —— cdc_4phase 复位偏斜验证（三版外部评审共识的硬 Gate）
// -----------------------------------------------------------------------------
// 架构约束：CDC 两侧复位**同源**（同一板级复位网络），仅释放时刻存在域间偏斜。
// 本 TB 验证该约束下的行为：
//   每轮：两域复位同时拉低 -> 错开 K 拍释放（K = 0/3/7/10）
//   1) 释放后静默 60 拍：断言无任何 O_dst_valid（无伪事务）
//   2) 连续两笔事务：断言恰好两次 valid、数据与顺序正确（无重复/无丢失）
// 说明：运行时"单侧复位"不在支持范围（需 epoch/re-arm 协议），已在
// cdc_4phase.v 头部契约注明。
// =============================================================================
module tb_cdc_reset_skew;
    reg src_clk = 0, dst_clk = 0;
    always #7  src_clk = ~src_clk;    // ~71.4MHz（源域）
    always #10 dst_clk = ~dst_clk;    // 50MHz（目标域）

    reg        src_rst_n = 0, dst_rst_n = 0;
    reg [22:0] src_data = 23'd0;
    reg        src_load = 0;
    wire [22:0] dst_data;
    wire        dst_valid, src_busy;

    cdc_4phase #(.WIDTH(23)) u_dut (
        .I_clk_src(src_clk), .I_rst_src_n(src_rst_n),
        .I_src_data(src_data), .I_src_load(src_load), .O_src_busy(src_busy),
        .I_clk_dst(dst_clk), .I_rst_dst_n(dst_rst_n),
        .O_dst_data(dst_data), .O_dst_valid(dst_valid)
    );

    integer    valid_cnt = 0, mism = 0, rounds = 0;
    reg [22:0] cap;

    always @(posedge dst_clk) begin
        if (dst_valid) begin
            valid_cnt = valid_cnt + 1;
            cap = dst_data;
        end
    end

    // 一轮：错开 skew 拍释放 -> 静默窗 -> 两笔事务
    task run_round;
        input integer skew;
        input [22:0] d0;
        input [22:0] d1;
        begin
            // 两域复位同时拉低
            src_rst_n = 1'b0; dst_rst_n = 1'b0;
            repeat (5) @(posedge src_clk);
            repeat (2) @(posedge dst_clk);
            valid_cnt = 0;

            // 错开释放：源域先、目标域晚 skew 个目标周期
            src_rst_n = 1'b1;
            repeat (skew) @(posedge dst_clk);
            dst_rst_n = 1'b1;

            // 静默窗：不得出现任何 valid
            repeat (60) @(posedge dst_clk);
            if (valid_cnt != 0) begin
                $display("MISM skew=%0d 静默期伪 valid x%0d", skew, valid_cnt);
                mism = mism + 1;
            end

            // 第一笔
            @(posedge src_clk); src_data = d0; src_load = 1'b1;
            @(posedge src_clk); src_load = 1'b0;
            repeat (150) @(posedge dst_clk);
            if (valid_cnt != 1) begin
                $display("MISM skew=%0d 第一笔 valid 次数=%0d（期望 1）", skew, valid_cnt);
                mism = mism + 1;
            end else if (cap !== d0) begin
                $display("MISM skew=%0d 第一笔数据 %h != %h", skew, cap, d0);
                mism = mism + 1;
            end

            // 第二笔
            @(posedge src_clk); src_data = d1; src_load = 1'b1;
            @(posedge src_clk); src_load = 1'b0;
            repeat (150) @(posedge dst_clk);
            if (valid_cnt != 2) begin
                $display("MISM skew=%0d 第二笔后累计 valid=%0d（期望 2）", skew, valid_cnt);
                mism = mism + 1;
            end else if (cap !== d1) begin
                $display("MISM skew=%0d 第二笔数据 %h != %h", skew, cap, d1);
                mism = mism + 1;
            end

            rounds = rounds + 1;
            $display("ROUND skew=%0d OK (valid_cnt=%0d)", skew, valid_cnt);
        end
    endtask

    initial begin
        run_round(0,  23'h012345, 23'h0ABCDE);
        run_round(3,  23'h100001, 23'h200002);
        run_round(7,  23'h3FFFFF, 23'h000001);
        run_round(10, 23'h155555, 23'h2AAAAA);
        if (mism == 0)
            $display("TB_RESULT: PASS rounds=%0d", rounds);
        else
            $display("TB_RESULT: FAIL mism=%0d", mism);
        $finish;
    end

    initial begin
        #2000000;
        $display("TB_RESULT: FAIL timeout mism=%0d rounds=%0d", mism, rounds);
        $finish;
    end
endmodule
\end{Verbatim}
\subsection{tb\_rgb2hsv.v}
\begin{Verbatim}[numbers=left,numbersep=5pt,frame=single,fontsize=\footnotesize,breaklines=true,breakanywhere=true]
// =============================================================================
// tb_rgb2hsv.v —— RGB->HSV 位精确校验：10010 个向量（含角点）逐个比对
// 期望值：tools/gen_vectors.py 由 golden_model 位精确生成（rgb2hsv_exp.hex）
// =============================================================================
`timescale 1ns/1ps
module tb_rgb2hsv;
    localparam N = 10010;
    reg clk = 0, rst_n = 0;
    always #7 clk = ~clk;

    reg        in_valid;
    reg  [23:0] in_data;
    wire       o_valid, o_sof, o_eol;
    wire [8:0] o_h;
    wire [7:0] o_s, o_v;

    rgb2hsv dut (.I_clk(clk), .I_rst_n(rst_n),
        .I_valid(in_valid), .I_sof(1'b0), .I_eol(1'b0), .I_data(in_data),
        .O_valid(o_valid), .O_sof(o_sof), .O_eol(o_eol),
        .O_h(o_h), .O_s(o_s), .O_v(o_v));

    reg [23:0] vin  [0:N-1];
    reg [27:0] vexp [0:N-1];   // {H[11:0], S[7:0], V[7:0]} 28bit
    integer i, idx = 0, mism = 0;

    initial begin
        $readmemh("rgb2hsv_in.hex", vin);
        $readmemh("rgb2hsv_exp.hex", vexp);
    end

    // 1:1 有序流：第 oc 个输出即第 oc 个输入的期望
    integer oc = 0;
    always @(posedge clk) begin
        if (rst_n && o_valid) begin
            if ({3'b000, o_h, o_s, o_v} !== vexp[oc]) begin
                if (mism < 6)
                    $display("MISM idx=%0d got h=%0d s=%0d v=%0d exp=%h",
                             oc, o_h, o_s, o_v, vexp[oc]);
                mism = mism + 1;
            end
            oc = oc + 1;
        end
    end

    initial begin
        repeat (3) @(posedge clk);
        rst_n <= 1;
        @(posedge clk);
        for (i = 0; i < N; i = i + 1) begin
            in_valid <= 1'b1;
            in_data  <= vin[i];
            idx = idx + 1;
            @(posedge clk);
        end
        in_valid <= 1'b0;
        repeat (8) @(posedge clk);
        if ((mism == 0) && (oc == N))
            $display("TB_RESULT: PASS (vectors=%0d)", N);
        else
            $display("TB_RESULT: FAIL mism=%0d/%0d oc=%0d", mism, N, oc);
        $finish;
    end
endmodule
\end{Verbatim}
\subsection{tb\_hsv\_threshold.v}
\begin{Verbatim}[numbers=left,numbersep=5pt,frame=single,fontsize=\footnotesize,breaklines=true,breakanywhere=true]
// =============================================================================
// tb_hsv_threshold.v —— 阈值分割位精确校验：4000 向量 x 8 组预置随机覆盖
// =============================================================================
`timescale 1ns/1ps
module tb_hsv_threshold;
    localparam N = 4000;
    reg clk = 0, rst_n = 0;
    always #7 clk = ~clk;

    reg        in_valid;
    reg  [8:0] in_h;
    reg  [7:0] in_s, in_v;
    reg  [2:0] in_p;
    wire       o_valid, o_sof, o_eol, o_mask;

    hsv_threshold dut (.I_clk(clk), .I_rst_n(rst_n),
        .I_valid(in_valid), .I_sof(1'b0), .I_eol(1'b0),
        .I_h(in_h), .I_s(in_s), .I_v(in_v), .I_preset_idx(in_p),
        .O_valid(o_valid), .O_sof(o_sof), .O_eol(o_eol), .O_mask(o_mask));

    reg [31:0] vin  [0:N-1];   // {H[11:0], S[15:8], V[7:0], preset[3:0]}
    reg [2:0]  vpre [0:N-1];
    reg        vexp [0:N-1];
    integer i, idx = 0, mism = 0;

    initial begin
        $readmemh("threshold_in.hex", vin);
        $readmemh("threshold_exp.hex", vexp);
    end

    integer oc = 0;
    always @(posedge clk) begin
        if (rst_n && o_valid) begin
            if (o_mask !== vexp[oc]) begin
                if (mism < 6)
                    $display("MISM idx=%0d got=%b exp=%b", oc, o_mask, vexp[oc]);
                mism = mism + 1;
            end
            oc = oc + 1;
        end
    end

    initial begin
        repeat (3) @(posedge clk);
        rst_n <= 1;
        @(posedge clk);
        for (i = 0; i < N; i = i + 1) begin
            in_valid <= 1'b1;
            {in_h, in_s, in_v} <= vin[i][31:4];   // 28bit {H,S,V}
            in_p     <= vin[i][3:0];      // 文件自带预置组（=i%8）
            idx = idx + 1;
            @(posedge clk);
        end
        in_valid <= 1'b0;
        repeat (6) @(posedge clk);
        if (mism == 0)
            $display("TB_RESULT: PASS (vectors=%0d)", N);
        else
            $display("TB_RESULT: FAIL mism=%0d/%0d oc=%0d", mism, N, oc);
        $finish;
    end
endmodule
\end{Verbatim}
\subsection{tb\_bin\_smooth.v}
\begin{Verbatim}[numbers=left,numbersep=5pt,frame=single,fontsize=\footnotesize,breaklines=true,breakanywhere=true]
// =============================================================================
// tb_bin_smooth.v —— 3x3 多数表决平滑位精确校验（64x48 迷你几何 x 4 帧）
// 输入带 sof/eol 帧行定界；逐像素比对 smooth_exp.hex。
// =============================================================================
`timescale 1ns/1ps
module tb_bin_smooth;
    localparam W = 64, H = 48, FRAMES = 4, TOTAL = W*H*FRAMES;
    reg clk = 0, rst_n = 0;
    always #7 clk = ~clk;

    reg        in_valid, in_sof, in_eol, in_mask;
    wire       o_valid, o_sof, o_eol, o_mask;

    bin_smooth3x3 #(.IMG_WIDTH(W), .IMG_HEIGHT(H)) dut (
        .I_clk(clk), .I_rst_n(rst_n),
        .I_valid(in_valid), .I_sof(in_sof), .I_eol(in_eol), .I_mask(in_mask),
        .O_valid(o_valid), .O_sof(o_sof), .O_eol(o_eol), .O_mask(o_mask));

    reg        vin  [0:TOTAL-1];
    reg        vexp [0:TOTAL-1];
    integer    i = 0, o = 0, mism = 0, x = 0, y = 0;
    reg        x_sof, x_eol;

    initial begin
        $readmemh("smooth_in.hex", vin);
        $readmemh("smooth_exp.hex", vexp);
    end

    // 驱动：1 像素/拍，sof/eol 按几何
    always @(posedge clk) begin
        if (rst_n && (i < TOTAL+1)) begin           // 多喂 1 个冲刷令牌
            in_valid <= 1'b1;
            in_mask  <= (i == TOTAL) ? 1'b0 : vin[i];
            x_sof    = (i % (W*H) == 0);
            x_eol    = (x == W-1);
            in_sof   <= x_sof;
            in_eol   <= x_eol;
            if (x == W-1) begin
                x <= 0;
                y <= y + 1;
            end else begin
                x <= x + 1;
            end
            i = i + 1;
        end else begin
            in_valid <= 1'b0; in_sof <= 1'b0; in_eol <= 1'b0;
        end
    end

    // 比对
    always @(posedge clk) begin
        if (rst_n && o_valid) begin
            if (o_mask !== vexp[o]) begin
                if (mism < 8)
                    $display("MISM px=%0d got=%b exp=%b", o, o_mask, vexp[o]);
                mism = mism + 1;
            end
            o = o + 1;
            if (o == TOTAL) begin
                if (mism == 0)
                    $display("TB_RESULT: PASS (pixels=%0d)", o);
                else
                    $display("TB_RESULT: FAIL mism=%0d/%0d", mism, o);
                $finish;
            end
        end
    end

    initial begin
        repeat (3) @(posedge clk);
        rst_n <= 1;
        #20000000;
        $display("TB_RESULT: FAIL timeout o=%0d/%0d", o, TOTAL);
        $finish;
    end
endmodule
\end{Verbatim}
\subsection{tb\_centroid.v}
\begin{Verbatim}[numbers=left,numbersep=5pt,frame=single,fontsize=\footnotesize,breaklines=true,breakanywhere=true]
// =============================================================================
// tb_centroid.v —— 质心/包围框/除零保护校验（128x72 迷你几何，3 帧）
// 帧0 实心块 / 帧1 噪声+块 / 帧2 微块(<MIN_PIXELS -> 无效+lost_cnt)
// =============================================================================
`timescale 1ns/1ps
module tb_centroid;
    localparam W = 128, H = 72, FRAMES = 3, TOTAL = W*H*FRAMES;
    reg clk = 0, rst_n = 0;
    always #7 clk = ~clk;

    reg        in_valid, in_sof, in_eol, in_mask;
    wire       t_valid, r_upd;
    wire [10:0] cx, cy, minx, maxx;
    wire [9:0]  miny, maxy;
    wire [20:0] cnt;
    wire [15:0] lost;

    centroid_accum #(.IMG_WIDTH(W), .IMG_HEIGHT(H), .MIN_PIXELS(30)) dut (
        .I_clk(clk), .I_rst_n(rst_n),
        .I_valid(in_valid), .I_sof(in_sof), .I_eol(in_eol), .I_mask(in_mask),
        .O_target_valid(t_valid), .O_cx(cx), .O_cy(cy), .O_count(cnt),
        .O_min_x(minx), .O_max_x(maxx), .O_min_y(miny), .O_max_y(maxy),
        .O_lost_cnt(lost), .O_results_upd(r_upd), .O_frame_tick());

    reg vin [0:TOTAL-1];
    integer fd, i = 0, x = 0, y = 0, frame = 0, mism = 0, scanned;
    reg exp_v; integer exp_cnt, exp_cx, exp_cy, exp_x0, exp_x1, exp_y0, exp_y1;

    initial begin
        $readmemh("centroid_in.hex", vin);
        fd = $fopen("centroid_exp.txt", "r");
    end

    // 驱动
    always @(posedge clk) begin
        if (rst_n && (i < TOTAL)) begin
            in_valid <= 1'b1;
            in_mask  <= vin[i];
            in_sof   <= (i % (W*H) == 0);
            in_eol   <= (x == W-1);
            if (x == W-1) begin
                x <= 0; y <= y + 1;
            end else begin
                x <= x + 1;
            end
            i = i + 1;
        end else begin
            in_valid <= 1'b0; in_sof <= 1'b0; in_eol <= 1'b0;
        end
    end

    // 结果比对（每帧 results_upd 时读一行期望）
    task check_line;
        begin
            scanned = $fscanf(fd, "%1x%5x%3x%3x%3x%3x%3x%3x",
                              exp_v, exp_cnt, exp_cx, exp_cy,
                              exp_x0, exp_x1, exp_y0, exp_y1);
            if (scanned != 8) begin
                $display("TB_RESULT: FAIL exp-file-short");
                $finish;
            end
            if ((t_valid !== (exp_v == 1)) ||
                (exp_v == 1 && ((cnt !== exp_cnt[20:0]) || (cx !== exp_cx[10:0]) ||
                 (cy !== exp_cy[10:0]) || (minx !== exp_x0[10:0]) ||
                 (maxx !== exp_x1[10:0]) || (miny !== exp_y0[9:0]) ||
                 (maxy !== exp_y1[9:0])))) begin
                $display("MISM frame=%0d got v=%b cnt=%0d cx=%0d cy=%0d x[%0d..%0d] y[%0d..%0d] | exp v=%0d cnt=%0d cx=%0d cy=%0d x[%0d..%0d] y[%0d..%0d]",
                         frame, t_valid, cnt, cx, cy, minx, maxx, miny, maxy,
                         exp_v, exp_cnt, exp_cx, exp_cy, exp_x0, exp_x1, exp_y0, exp_y1);
                $display("                 exp v=%0d cnt=%0d cx=%0d cy=%0d x[%0d..%0d] y[%0d..%0d]",
                         exp_v, exp_cnt, exp_cx, exp_cy, exp_x0, exp_x1, exp_y0, exp_y1);
                mism = mism + 1;
            end else begin
                $display("FRAME%0d OK v=%0d cnt=%0d cx=%0d cy=%0d lost=%0d",
                         frame, t_valid, cnt, cx, cy, lost);
            end
            frame = frame + 1;
            if (frame == FRAMES) begin
                if (mism == 0)
                    $display("TB_RESULT: PASS (frames=%0d)", FRAMES);
                else
                    $display("TB_RESULT: FAIL mism=%0d", mism);
                $finish;
            end
        end
    endtask

    always @(posedge clk) begin
        if (rst_n && r_upd)
            check_line;
    end

    initial begin
        repeat (3) @(posedge clk);
        rst_n <= 1;
        #30000000;
        $display("TB_RESULT: FAIL timeout frame=%0d", frame);
        $finish;
    end
endmodule
\end{Verbatim}
\subsection{tb\_detect\_chain.v}
\begin{Verbatim}[numbers=left,numbersep=5pt,frame=single,fontsize=\footnotesize,breaklines=true,breakanywhere=true]
// =============================================================================
// tb_detect_chain.v —— 检测链整链 TB（真实场景 3 帧，全尺寸 1280x720）
// 场景：红块+噪声 / 红块+品红干扰物 / 小目标。
// 验收：1) 每帧 results_upd 与 golden（chain_exp.txt）一致；
//       2) 窃流直通输出与输入逐拍一致（不干扰主线）。
// =============================================================================
`timescale 1ns/1ps
module tb_detect_chain;
    reg clk = 0, rst_n = 0;
    always #7 clk = ~clk;

    wire       src_valid, src_user, src_last;
    wire [95:0] src_data;

    sim_video_src #(
        .BEAT_FILE("chain_beats.hex"), .FLAG_FILE("chain_flags.hex"),
        .BEATS_TOTAL(3*720*320), .FRAMES(3)
    ) u_src (
        .I_clk(clk), .I_rst_n(rst_n),
        .O_tvalid(src_valid), .O_tdata(src_data),
        .O_tuser(src_user), .O_tlast(src_last), .I_tready(1'b1));

    wire       pt_valid, pt_user, pt_last;
    wire [95:0] pt_data;
    wire       t_valid, r_upd;
    wire [10:0] cx, cy, minx, maxx;
    wire [9:0]  miny, maxy;
    wire [20:0] cnt;
    wire [15:0] lost;

    detect_top dut (
        .I_clk(clk), .I_rst_n(rst_n),
        .I_tdata(src_data), .I_tvalid(src_valid), .I_tuser(src_user), .I_tlast(src_last),
        .O_tready(), .O_tdata(pt_data), .O_tvalid(pt_valid),
        .O_tuser(pt_user), .O_tlast(pt_last),
        .I_preset_idx(3'd0),
        .O_target_valid(t_valid), .O_cx(cx), .O_cy(cy), .O_count(cnt),
        .O_min_x(minx), .O_max_x(maxx), .O_min_y(miny), .O_max_y(maxy),
        .O_lost_cnt(lost), .O_results_upd(r_upd), .O_frame_tick(),
        .O_dbg_mask_valid(), .O_dbg_mask(), .O_s_valid(), .O_s_sof(),
        .O_s_eol(), .O_s_data());

    // 直通一致性
    integer pass_mism = 0, frames = 0, mism = 0;
    always @(posedge clk) begin
        if (rst_n && src_valid) begin
            if ((pt_data !== src_data) || (pt_valid !== 1'b1) ||
                (pt_user !== src_user) || (pt_last !== src_last)) begin
                if (pass_mism < 5)
                    $display("PASSTHRU_MISM t=%0t", $time);
                pass_mism = pass_mism + 1;
            end
        end
    end

    // 每帧结果比对
    integer fd, scanned;
    reg exp_v; integer exp_cnt, exp_cx, exp_cy, exp_x0, exp_x1, exp_y0, exp_y1;
    task check_line;
        begin
            scanned = $fscanf(fd, "%1x%5x%3x%3x%3x%3x%3x%3x",
                              exp_v, exp_cnt, exp_cx, exp_cy,
                              exp_x0, exp_x1, exp_y0, exp_y1);
            if ((t_valid !== (exp_v == 1)) ||
                (exp_v == 1 && ((cnt !== exp_cnt[20:0]) || (cx !== exp_cx[10:0]) || (cy !== exp_cy[10:0]) ||
                 (minx !== exp_x0[10:0]) || (maxx !== exp_x1[10:0]) ||
                 (miny !== exp_y0[9:0]) || (maxy !== exp_y1[9:0])))) begin
                $display("MISM frame=%0d got v=%b cnt=%0d cx=%0d cy=%0d x[%0d..%0d] y[%0d..%0d] | exp v=%0d cnt=%0d cx=%0d cy=%0d x[%0d..%0d] y[%0d..%0d]",
                         frames, t_valid, cnt, cx, cy, minx, maxx, miny, maxy,
                         exp_v, exp_cnt, exp_cx, exp_cy, exp_x0, exp_x1, exp_y0, exp_y1);
                mism = mism + 1;
            end else begin
                $display("FRAME%0d OK v=%0d cnt=%0d cx=%0d cy=%0d",
                         frames, t_valid, cnt, cx, cy);
            end
            frames = frames + 1;
        end
    endtask

    always @(posedge clk) begin
        if (rst_n && r_upd && frames < 3)
            check_line;
    end

    initial begin
        fd = $fopen("chain_exp.txt", "r");
        repeat (3) @(posedge clk);
        rst_n <= 1;
        #80000000;
        if ((mism == 0) && (pass_mism == 0) && (frames == 3))
            $display("TB_RESULT: PASS (frames=%0d passthru_ok)", frames);
        else
            $display("TB_RESULT: FAIL frames=%0d mism=%0d passthru_mism=%0d",
                     frames, mism, pass_mism);
        $finish;
    end
endmodule
\end{Verbatim}
\subsection{tb\_gimbal.v}
\begin{Verbatim}[numbers=left,numbersep=5pt,frame=single,fontsize=\footnotesize,breaklines=true,breakanywhere=true]
// =============================================================================
// tb_gimbal.v —— 云台控制律表驱动校验（死区/限速/限幅/丢失保持）
// 激励与期望由 golden_model.gimbal_ctrl_step 位精确生成。
// =============================================================================
`timescale 1ns/1ps
module tb_gimbal;
    localparam N = 10;
    reg clk = 0, rst_n = 0;
    always #10 clk = ~clk;   // 50MHz 域

    reg        tick = 0, track_en = 1, t_valid = 0;
    reg [10:0] cx = 0, cy = 0;
    wire [12:0] duty_x, duty_y;
    wire [15:0] lost;
    wire [7:0]  status;

    gimbal_ctrl dut (
        .I_clk(clk), .I_rst_n(rst_n),
        .I_frame_tick(tick), .I_track_en(track_en),
        .I_target_valid(t_valid), .I_cx(cx), .I_cy(cy),
        .O_duty_x_us(duty_x), .O_duty_y_us(duty_y),
        .O_lost_cnt(lost), .O_status(status));

    integer fd_s, fd_e, i, scanned, scanned_e, mism = 0;
    reg stim_v; integer stim_cx, stim_cy;
    integer e_dx, e_dy, e_lost, e_dzx, e_dzy;

    initial begin
        fd_s = $fopen("gimbal_stim.txt", "r");
        fd_e = $fopen("gimbal_exp.txt", "r");
        repeat (3) @(posedge clk);
        rst_n <= 1;
        repeat (2) @(posedge clk);
        for (i = 0; i < N; i = i + 1) begin
            scanned = $fscanf(fd_s, "%1x%3x%3x", stim_v, stim_cx, stim_cy);
            scanned_e = $fscanf(fd_e, "%3x%3x%4x%1x%1x",
                                e_dx, e_dy, e_lost, e_dzx, e_dzy);
            t_valid <= stim_v;
            cx <= stim_cx[10:0];
            cy <= stim_cy[10:0];
            @(posedge clk);
            tick <= 1'b1;
            @(posedge clk);
            tick <= 1'b0;
            repeat (2) @(posedge clk);   // 等 O_duty 锁存
            if ((duty_x !== e_dx[12:0]) || (duty_y !== e_dy[12:0]) ||
                (lost !== e_lost[15:0]) ||
                (status[0] !== e_dzx[0]) || (status[1] !== e_dzy[0])) begin
                $display("MISM i=%0d got x=%0d y=%0d lost=%0d st=%b | exp x=%0d y=%0d lost=%0d dz=%b%b",
                         i, duty_x, duty_y, lost, status, e_dx, e_dy, e_lost, e_dzy, e_dzx);
                mism = mism + 1;
            end else begin
                $display("STEP%0d OK x=%0d y=%0d lost=%0d", i, duty_x, duty_y, lost);
            end
        end
        if (mism == 0)
            $display("TB_RESULT: PASS (steps=%0d)", N);
        else
            $display("TB_RESULT: FAIL mism=%0d", mism);
        $finish;
    end
endmodule
\end{Verbatim}
\subsection{tb\_pwm.v}
\begin{Verbatim}[numbers=left,numbersep=5pt,frame=single,fontsize=\footnotesize,breaklines=true,breakanywhere=true]
// =============================================================================
// tb_pwm.v —— 双路 50Hz PWM 校验（v2：连续监视器，无事件同步竞态）
// 方案：自由运转的计数器统计每周期脉宽/长度，period_tick 拍锁存结果；
//       按 tick 序号在两段窗口各校验一次（上电中位 / 新占空）。
// =============================================================================
`timescale 1ns/1ps
module tb_pwm;
    reg clk = 0, rst_n = 0;
    always #10 clk = ~clk;   // 50MHz

    reg         update = 0;
    reg [12:0]  dx = 13'd1500, dy = 13'd1500;
    wire        pwm_x, pwm_y, tick;

    pwm_servo dut (.I_clk(clk), .I_rst_n(rst_n), .I_update(update),
        .I_duty_x_us(dx), .I_duty_y_us(dy),
        .O_pwm_x(pwm_x), .O_pwm_y(pwm_y),
        .O_period_tick(tick), .O_duty_x_dbg(), .O_duty_y_dbg());

    // ---------------- 连续监视器 ----------------
    reg [20:0] c_x, c_y, c_len;
    reg [20:0] l_x, l_y, l_len;
    reg [15:0] tick_no;
    integer    mism = 0;

    always @(posedge clk or negedge rst_n) begin
        if (!rst_n) begin
            c_x <= 0; c_y <= 0; c_len <= 0;
            l_x <= 0; l_y <= 0; l_len <= 0;
            tick_no <= 0;
        end else if (tick) begin
            // 含当拍锁存：latched 窗口 = 恰好一个完整周期
            l_x <= c_x   + (pwm_x ? 21'd1 : 21'd0);
            l_y <= c_y   + (pwm_y ? 21'd1 : 21'd0);
            l_len <= c_len + 21'd1;
            c_x <= 0;   c_y <= 0;   c_len <= 0;
            tick_no <= tick_no + 16'd1;
        end else begin
            c_x   <= c_x   + (pwm_x ? 21'd1 : 21'd0);
            c_y   <= c_y   + (pwm_y ? 21'd1 : 21'd0);
            c_len <= c_len + 21'd1;
        end
    end

    // ---------------- 校验调度（按 tick 序号） ----------------
    // 期望表：tick 2 为上电中位 1500µs（75000 拍）；tick 6 为新占空
    // 1200/2000µs（60000/100000 拍）。tick 3..5 为换挡过渡（RTL 采用
    // 周期边界换装：tick3 施加请求 -> tick5 起点生效 -> tick6 窗口纯新占空）。
    task check_tick;
        input [15:0] no;
        input [20:0] exp_x;
        input [20:0] exp_y;
        begin
            if ((l_len != 21'd1000000) || (l_x != exp_x) || (l_y != exp_y)) begin
                $display("MISM tick=%0d got x=%0d y=%0d len=%0d | exp x=%0d y=%0d len=1000000",
                         no, l_x, l_y, l_len, exp_x, exp_y);
                mism = mism + 1;
            end else begin
                $display("MEAS tick=%0d OK x=%0d y=%0d len=%0d", no, l_x, l_y, l_len);
            end
        end
    endtask

    always @(posedge clk) begin
        if (rst_n && tick) begin
            case (tick_no)
                16'd2: check_tick(16'd2, 21'd75000, 21'd75000);  // 首个完整周期（上电中位）
                16'd3: begin
                    // 换占空：延后 2 拍施加载荷（req 寄存器），
                    // 周期 5 起点边界换装，周期 5->6 成为首个纯新占空完整周期
                    dx     <= 13'd1200; dy     <= 13'd2000;
                    update <= 1'b1;
                end
                16'd4: update <= 1'b0;
                16'd6: begin
                    check_tick(16'd6, 21'd60000, 21'd100000);    // 纯新占空完整周期
                    if (mism == 0)
                        $display("TB_RESULT: PASS");
                    else
                        $display("TB_RESULT: FAIL mism=%0d", mism);
                    $finish;
                end
                default: ;
            endcase
        end
    end

    initial begin
        repeat (3) @(posedge clk);
        rst_n <= 1;
        #150000000;   // 6 个完整周期(120ms) + 余量
        $display("TB_RESULT: FAIL timeout mism=%0d", mism);
        $finish;
    end
endmodule
\end{Verbatim}
\subsection{tb\_osd\_track.v}
\begin{Verbatim}[numbers=left,numbersep=5pt,frame=single,fontsize=\footnotesize,breaklines=true,breakanywhere=true]
// =============================================================================
// tb_osd_track.v —— OSD 叠加层帧级校验 v2（双缓冲消撕裂验证）
// 时序：帧 0 中段（y=360）注入目标包 -> CDC 落入影子寄存器（帧 0 画面不受影响），
//       帧 1 sof 换装 -> 帧 1/帧 2 整帧使用新状态。
// 比对：dump osd_act_f0/f1/f2.hex 与 Python 金标准 osd_exp_f0/f1/f2.hex 逐位比对
//       （由 run_all 的 frame check 完成）。
//       f0 = 空状态（无文本无框，跑马灯滚动）；f1/f2 = 完整 OSD（跑马灯偏移递增）。
//       若双缓冲失效（包到达即生效），f0 将被撕裂，不可能匹配金标准。
// =============================================================================
`timescale 1ns/1ps
module tb_osd_track;
    localparam FLEN = 1280*720, FRAMES = 3;
    reg clk = 0, rst_n = 0;
    always #7 clk = ~clk;

    reg        in_valid, in_sof, in_eol;
    reg [23:0] in_data;
    wire       o_valid, o_sof, o_eol;
    wire [23:0] o_data;

    reg  [134:0] pack;
    reg         pack_load = 0;
    wire        pack_busy;

    osd_track dut (
        .I_clk(clk), .I_rst_n(rst_n),
        .I_valid(in_valid), .I_sof(in_sof), .I_eol(in_eol), .I_data(in_data),
        .I_det_clk(clk), .I_det_rst_n(rst_n),
        .I_det_pack(pack), .I_det_load(pack_load), .O_det_busy(pack_busy),
        .I_bright(7'd0), .I_contrast(8'd64),
        .O_valid(o_valid), .O_sof(o_sof), .O_eol(o_eol), .O_data(o_data));

    reg [23:0] base [0:FLEN-1];
    integer f0, f1, f2, i = 0, x = 0, y = 0;
    integer fdc, scanned;
    integer got = 0;
    integer w0 = 0, w1 = 0, w2 = 0;
    reg v_in; integer c_cx, c_cy, c_x0, c_x1, c_y0, c_y1;

    initial begin
        // 句柄必须先于首像素打开：若在首个 o_valid 边沿 fopen，与 $fwrite 同拍
        // 竞态，像素 0 会写进 fd0（stdout）且整个文件前移一格（实测教训）
        f0 = $fopen("osd_act_f0.hex", "w");
        f1 = $fopen("osd_act_f1.hex", "w");
        f2 = $fopen("osd_act_f2.hex", "w");
        $readmemh("osd_base.hex", base);
        fdc = $fopen("osd_coord.txt", "r");
        scanned = $fscanf(fdc, "%1x%3x%3x%3x%3x%3x%3x",
                          v_in, c_cx, c_cy, c_x0, c_x1, c_y0, c_y1);
        // 运动框字段填入与金标准一致的值（外脑评审补覆盖）：
        // mot_valid=1, mot_cnt=64, bbox=(200,200)-(260,240)
        pack = {v_in, 1'b0, 3'd0, c_cx[10:0], c_cy[10:0], c_x0[10:0], c_x1[10:0],
                c_y0[9:0], c_y1[9:0], 8'd0, 1'b1, 14'd64,
                11'd200, 11'd260, 10'd200, 10'd240, 1'b0};   // 135bit：bit0 保留位必须补齐
        $display("SCAN scanned=%0d v_in=%b cx=%0d cy=%0d x0=%0d x1=%0d y0=%0d y1=%0d",
                 scanned, v_in, c_cx, c_cy, c_x0, c_x1, c_y0, c_y1);
    end

    // 帧 0 中段（y=360 的行中点）注入目标包：影子寄存器当帧更新，画面不变
    always @(posedge clk) begin
        if (rst_n && (y == 360) && (x == 640))
            pack_load <= 1'b1;
        else if (pack_load)
            pack_load <= 1'b0;
    end

    // 驱动底图
    always @(posedge clk) begin
        if (rst_n && (i < FLEN*FRAMES)) begin
            in_valid <= 1'b1;
            in_data  <= base[i % FLEN];
            in_sof   <= (i % FLEN == 0);
            in_eol   <= (x == 1280-1);
            if (x == 1280-1) begin
                x <= 0; y <= y + 1;
            end else begin
                x <= x + 1;
            end
            i = i + 1;
        end else begin
            in_valid <= 1'b0; in_sof <= 1'b0; in_eol <= 1'b0;
        end
    end

    // 采集输出（按帧分文件）。快照+延迟一拍写入，避免与 NBA 同沿读写的耦合。
    reg [23:0] od_d;
    reg        ov_d;
    always @(posedge clk) begin
        od_d <= o_data;
        ov_d <= o_valid;
    end

    always @(posedge clk) begin
        if (rst_n && ov_d) begin
            case (got / FLEN)
                0: begin
                    $fwrite(f0, "%02X%02X%02X\n", od_d[23:16], od_d[15:8], od_d[7:0]);
                    w0 = w0 + 1;
                end
                1: begin
                    $fwrite(f1, "%02X%02X%02X\n", od_d[23:16], od_d[15:8], od_d[7:0]);
                    w1 = w1 + 1;
                end
                default: begin
                    $fwrite(f2, "%02X%02X%02X\n", od_d[23:16], od_d[15:8], od_d[7:0]);
                    w2 = w2 + 1;
                end
            endcase
            got = got + 1;
            if (got == FLEN*FRAMES) begin
                $fflush(f0); $fclose(f0);
                $fflush(f1); $fclose(f1);
                $fflush(f2); $fclose(f2);
                $display("TB_RESULT: PASS (dumped=%0d w0=%0d w1=%0d w2=%0d)",
                         got, w0, w1, w2);
                $finish;
            end
        end
    end

    initial begin
        repeat (3) @(posedge clk);
        rst_n <= 1;
        #130000000;
        $display("TB_RESULT: FAIL timeout got=%0d", got);
        $finish;
    end
endmodule
\end{Verbatim}
\subsection{tb\_vision\_top.v}
\begin{Verbatim}[numbers=left,numbersep=5pt,frame=single,fontsize=\footnotesize,breaklines=true,breakanywhere=true]
// =============================================================================
// tb_vision_top.v —— 整机联调 TB：tpg(红块@500,300) -> vision_top
// 采集：第 2 帧 OSD 叠加图（frame_1）+ 云台脉宽方向。
// Python 侧：帧像素比对 + 质心/方向校验（run_all.py）。
// =============================================================================
`timescale 1ns/1ps
module tb_vision_top;
    reg pix_clk = 0, pwm_clk = 0, pix_rst = 0, pwm_rst = 0;
    always #7  pix_clk = ~pix_clk;
    always #10 pwm_clk = ~pwm_clk;

    reg  [10:0] box_x = 11'd500;
    reg  [9:0]  box_y = 10'd300;
    wire        tvalid, tuser, tlast;
    wire [95:0] tdata;

    tpg_720p #(.BOX_EN(1)) u_tpg (
        .I_clk(pix_clk), .I_rst_n(pix_rst),
        .I_box_x(box_x), .I_box_y(box_y),
        .O_tvalid(tvalid), .O_tdata(tdata), .O_tuser(tuser), .O_tlast(tlast),
        .I_tready(1'b1));

    wire        osd_valid, osd_sof, osd_eol;
    wire [23:0] osd_data;
    wire [95:0] pt_data;
    wire        pt_valid, pt_user, pt_last, tready;
    wire        pwm_x, pwm_y;
    wire [12:0] duty_x, duty_y;
    wire [15:0] lost;
    wire [7:0]  status;

    vision_top dut (
        .I_pix_clk(pix_clk), .I_pix_rst_n(pix_rst),
        .I_pwm_clk(pwm_clk), .I_pwm_rst_n(pwm_rst),
        .I_tdata(tdata), .I_tvalid(tvalid), .I_tuser(tuser), .I_tlast(tlast),
        .O_tready(tready),
        .O_tdata(pt_data), .O_tvalid(pt_valid), .O_tuser(pt_user), .O_tlast(pt_last),
        .O_osd_valid(osd_valid), .O_osd_sof(osd_sof),
        .O_osd_eol(osd_eol), .O_osd_data(osd_data),
        .I_preset_idx(3'd0), .I_track_en(1'b1),
        .I_algo_edge(1'b0), .I_bright(7'd0), .I_contrast(8'd64),
        .O_pwm_x(pwm_x), .O_pwm_y(pwm_y),
        .O_duty_x_us(duty_x), .O_duty_y_us(duty_y),
        .O_lost_cnt(lost), .O_status(status),
        .O_led_g(), .O_led_r(), .O_buzzer());

    frame_dump24 #(.BASE_NAME("vision_frame"), .MAX_FRAMES(5)) u_dump (
        .I_clk(pix_clk), .I_rst_n(pix_rst),
        .I_valid(osd_valid), .I_eol(osd_eol), .I_data(osd_data),
        .O_frame_done());

    integer done_frames = 0;
    integer fd;
    initial fd = $fopen("vision_duty.txt", "w");

    // 直通一致性
    integer pm = 0;
    always @(posedge pix_clk) begin
        if (pix_rst && tvalid && ((pt_data !== tdata) || (pt_valid !== 1'b1)))
            pm = pm + 1;
    end

    always @(posedge pix_clk) begin
        if (pix_rst && osd_valid && osd_eol && (done_frames == 1)) begin
            // 第 2 帧末（dump 帧序号 1 完成）
        end
    end

    always @(posedge pwm_clk) begin
        if (pwm_rst && dut.u_gimbal.O_lost_cnt === 16'd0) begin
        end
    end

    // 帧计数（用 dump 模块内部 frame_idx 层次引用）
    always @(posedge pix_clk) begin
        if (pix_rst && (u_dump.frame_idx == 5)) begin   // 第5帧末：frame_4.hex 已写完
            $fwrite(fd, "%0d %0d %0d %0d\n", duty_x, duty_y, lost, status);
            $fflush(fd);
            $fclose(fd);
            $display("DUTY_REPORT: x=%0d y=%0d lost=%0d status=%b passthru_mism=%0d",
                     duty_x, duty_y, lost, status, pm);
            if (pm == 0)
                $display("TB_RESULT: PASS (duty_x=%0d, expect python frame compare)", duty_x);
            else
                $display("TB_RESULT: FAIL passthru_mism=%0d", pm);
            $finish;
        end
    end

    initial begin
        repeat (4) @(posedge pix_clk);
        pix_rst <= 1;
        repeat (4) @(posedge pwm_clk);
        pwm_rst <= 1;
        #170000000;
        $display("TB_RESULT: FAIL timeout frames=%0d", u_dump.frame_idx);
        $finish;
    end
endmodule
\end{Verbatim}
\subsection{tb\_closed\_loop.v}
\begin{Verbatim}[numbers=left,numbersep=5pt,frame=single,fontsize=\footnotesize,breaklines=true,breakanywhere=true]
// =============================================================================
// tb_closed_loop.v —— 闭环收敛性系统 TB（RTL 在环）
// 相机模型（TB 内行为级）：云台角度 theta 由脉宽换算（0.18°/µs），
// 带一阶速率限制（1.5°/帧）；目标世界角 40° 固定；
// 目标像素 x = 640 + (40-theta)*20（20 px/°），TPG 方块实时跟随。
// 验收：60 帧内误差进入 +/-25px 死区并保持 10 帧（无振荡/发散）。
// =============================================================================
`timescale 1ns/1ps
module tb_closed_loop;
    reg pix_clk = 0, pwm_clk = 0, pix_rst = 0, pwm_rst = 0;
    always #7  pix_clk = ~pix_clk;
    always #10 pwm_clk = ~pwm_clk;

    reg  [10:0] box_x = 11'd600;
    reg  [9:0]  box_y = 10'd300;
    wire        tvalid, tuser, tlast;
    wire [95:0] tdata;

    tpg_720p #(.BOX_EN(1)) u_tpg (
        .I_clk(pix_clk), .I_rst_n(pix_rst),
        .I_box_x(box_x), .I_box_y(box_y),
        .O_tvalid(tvalid), .O_tdata(tdata), .O_tuser(tuser), .O_tlast(tlast),
        .I_tready(1'b1));

    wire [23:0] osd_data;
    wire        osd_valid, osd_sof, osd_eol;
    wire [12:0] duty_x;
    vision_top dut (
        .I_pix_clk(pix_clk), .I_pix_rst_n(pix_rst),
        .I_pwm_clk(pwm_clk), .I_pwm_rst_n(pwm_rst),
        .I_tdata(tdata), .I_tvalid(tvalid), .I_tuser(tuser), .I_tlast(tlast),
        .O_tready(), .O_tdata(), .O_tvalid(), .O_tuser(), .O_tlast(),
        .O_osd_valid(osd_valid), .O_osd_sof(osd_sof),
        .O_osd_eol(osd_eol), .O_osd_data(osd_data),
        .I_preset_idx(3'd0), .I_track_en(1'b1),
        .I_algo_edge(1'b0), .I_bright(7'd0), .I_contrast(8'd64),
        .O_pwm_x(), .O_pwm_y(),
        .O_duty_x_us(duty_x), .O_duty_y_us(),
        .O_lost_cnt(), .O_status(),
        .O_led_g(), .O_led_r(), .O_buzzer());

    // ---------------- 相机-云台行为模型 ----------------
    real theta = 30.0;          // 当前相机指向（°）
    real theta_cmd;
    real delta;
    integer frame = 0, locked_run = 0, settle_frame = -1, err_int;
    integer fd;
    initial fd = $fopen("closed_loop_trace.txt", "w");

    // 每帧末（detect 域 frame_tick）采样脉宽并推进模型
    wire frame_tick = dut.t_tick;
    always @(posedge pix_clk) begin
        if (pix_rst && frame_tick) begin
            theta_cmd = 30.0 + ($signed(duty_x) - 1500) * 0.18;  // 位置舵机：1500=保持 theta0（有符号！）
            delta = theta_cmd - theta;
            if (delta > 1.5)      delta = 1.5;
            else if (delta < -1.5) delta = -1.5;
            theta = theta + delta;
            err_int = $rtoi((40.0 - theta) * 20.0);
            $fwrite(fd, "%0d %0d %0d %f\n", frame, duty_x, err_int, theta);
            if ((err_int >= -25) && (err_int <= 25)) begin
                locked_run = locked_run + 1;
                if ((settle_frame < 0) && (locked_run >= 10)) begin
                    settle_frame = frame - 9;
                    $display("LOCKED: settle_frame=%0d", settle_frame);
                end
            end else begin
                locked_run = 0;
            end
            frame = frame + 1;
            // 更新下一帧 TPG 方块位置（相机视角）
            box_x <= 640 + err_int - 40;   // 隐式截断到 11 位（范围已保证非负）
            if (frame == 60) begin
                $fflush(fd);
                $fclose(fd);
                if (settle_frame >= 0)
                    $display("TB_RESULT: PASS settle_frame=%0d final_err=%0d",
                             settle_frame, err_int);
                else
                    $display("TB_RESULT: FAIL not locked, final_err=%0d", err_int);
                $finish;
            end
        end
    end

    initial begin
        repeat (4) @(posedge pix_clk);
        pix_rst <= 1;
        repeat (4) @(posedge pwm_clk);
        pwm_rst <= 1;
        #900000000;
        $display("TB_RESULT: FAIL timeout frame=%0d", frame);
        $finish;
    end
endmodule
\end{Verbatim}
\subsection{tb\_algo\_switch.v}
\fileintro{算法切换帧专项 TB（第 18 项，GPT-3 复审后新增）：TPG->vision\_top，第 3/5 帧输入algo\_edge 翻转，输出 921600 像素分窗断言"切换后每窗 SOF=1、EOL=720、分支内容符合切换后算法"——首跑即钉死"提交沿扇出窗口 SOF 丢失"（MISM win=5 sof=0），修复后 PASS。判据含 sof<=2/eol>=719 容忍切换窗固有跨窗语义。2026-10-03 复审响应再升级为 	extbf{SOF->SOF 严格 scoreboard}：SB1 每条 EOL 行恰 1280 像素；SB2 非切换跨度恰 720 行/921600 像素；SB3 切换跨度允许旧帧截断（<=719 行+残行）；SB4 截断事件 <= 命令切换数。复跑 PASS：SB frames=9 partial=2 sb\_fail=0。}
\begin{Verbatim}[numbers=left,numbersep=5pt,frame=single,fontsize=\footnotesize,breaklines=true,breakanywhere=true]
`timescale 1ns/1ps
// =============================================================================
// tb_algo_switch.v —— 算法切换帧对齐专项 TB（三版外部评审要求，2026-10-03）
// -----------------------------------------------------------------------------
// 场景：TPG(灰阶柱 + 彩色柱 + 静态红块) -> vision_top。
//   I_algo_edge 在输入帧 3 / 5 期间切换（0->1->0），请求在下一输入帧首被锁存、
//   由输出侧 SOF 提交逻辑换装。
// 判据（两层）：
//  【窗判据】对 O_osd_* 输出流按 921600 像素划分"输出帧窗"，跳过前 2 个启动窗：
//   1) 任何窗 sof >= 1（**关键**：绝不允许丢帧首——修复前切换窗 sof=0 即此弊）；
//   2) sof <= 2（切换跨越窗允许 2 次：旧帧截尾 + 新帧帧首，且此类窗 <= 2 个）；
//   3) eol  in  [719,720]（跨切换窗 719，稳态 720）；
//   4) valid 像素恰好 921600（含跨切换窗——输出流连续无洞）；
//   5) 分支内容真实生效：>=2 窗采样点全灰（edge 输出）且 >=2 窗采样点含彩色
//      （det 输出）——采样点取彩色柱区且避开 OSD 叠加/目标框。
//  【严格帧 scoreboard】（2026-10-03 复审要求，SOF -> 下一 SOF）：
//   SB1 每条 EOL 结束的行必须恰好 1280 valid 像素；SB2 非切换帧必须 720 行/921600 像素；
//   SB3 切换帧允许 719 行 + 1 条残行（或恰在行边界按完整帧）；SB4 残行事件 <= 命令切换数（2）。
// 语义备注：切换发生在目标分支输出 SOF 当拍（组合选择），旧分支帧被截短、
//   新分支从自身 SOF/首像素开始——下游各模块由每次 SOF 复位内部计数。
// =============================================================================
module tb_algo_switch;
    reg pix_clk = 0, pwm_clk = 0, pix_rst = 0, pwm_rst = 0;
    always #7  pix_clk = ~pix_clk;   // ~71.4MHz
    always #10 pwm_clk = ~pwm_clk;   // 50MHz

    reg [10:0] box_x = 11'd600;
    reg [9:0]  box_y = 10'd300;
    wire        tvalid, tuser, tlast;
    wire [95:0] tdata;

    tpg_720p #(.BOX_EN(1)) u_tpg (
        .I_clk(pix_clk), .I_rst_n(pix_rst),
        .I_box_x(box_x), .I_box_y(box_y),
        .O_tvalid(tvalid), .O_tdata(tdata), .O_tuser(tuser), .O_tlast(tlast),
        .I_tready(1'b1));

    reg         algo_edge = 1'b0;
    wire [23:0] osd_data;
    wire        osd_valid, osd_sof, osd_eol;

    vision_top dut (
        .I_pix_clk(pix_clk), .I_pix_rst_n(pix_rst),
        .I_pwm_clk(pwm_clk), .I_pwm_rst_n(pwm_rst),
        .I_tdata(tdata), .I_tvalid(tvalid), .I_tuser(tuser), .I_tlast(tlast),
        .O_tready(), .O_tdata(), .O_tvalid(), .O_tuser(), .O_tlast(),
        .O_osd_valid(osd_valid), .O_osd_sof(osd_sof),
        .O_osd_eol(osd_eol), .O_osd_data(osd_data),
        .I_preset_idx(3'd0), .I_track_en(1'b1),
        .I_algo_edge(algo_edge), .I_bright(7'd0), .I_contrast(8'd64),
        .O_pwm_x(), .O_pwm_y(),
        .O_duty_x_us(), .O_duty_y_us(),
        .O_lost_cnt(), .O_status(),
        .O_led_g(), .O_led_r(), .O_buzzer(),
        .O_frame_tick(), .O_target_valid(), .O_tgt_cx(), .O_tgt_cy());

    // ---------------- 切换调度（按输入帧号） ----------------
    integer in_frame = 0;
    always @(posedge pix_clk) begin
        if (pix_rst && tvalid && tuser) in_frame = in_frame + 1;
        if (in_frame == 3) algo_edge <= 1'b1;   // 帧 4 首被锁存 -> 请求切边缘
        if (in_frame == 5) algo_edge <= 1'b0;   // 帧 6 首被锁存 -> 请求切回
    end

    // ---------------- 输出流分窗监视 ----------------
    localparam FRAME_PX = 1280 * 720;
    integer x = 0, y = 0;
    integer px_cnt = 0, sof_cnt = 0, eol_cnt = 0, col_smpl = 0;
    integer win = 0, mism = 0, gray_wins = 0, color_wins = 0, switch_wins = 0;
    integer eval_wins = 0;

    // ---------------- 严格帧 scoreboard（2026-10-03 复审要求：SOF -> 下一 SOF） ----------------
    // 判据（比 921600 像素分窗严格；"切换跨度"= 结束 SOF 落在行中，即旧帧被提交瞬时截断）：
    //   SB1 每一条以 EOL 结束的行必须恰好 1280 个 valid 像素（禁止 1279）；
    //   SB2 非切换跨度（结束 SOF 紧接行尾）：必须恰好 720 行、921600 像素；
    //   SB3 切换跨度：允许 <=719 行 + 1 条残行（残行像素 1..1279，切点任意——
    //       实测两次切换分别为 0 行+1px 与 719 行+1279px，源于两分支流水延迟差）；
    //   SB4 切换跨度总数 <= 本 TB 命令的切换次数（2）——多余即未命令的截断，FAIL。
    integer sb_active = 0, sb_lines = 0, sb_px_line = 0, sb_px_total = 0;
    integer sb_frames = 0, sb_fail = 0, sb_partial = 0, sb_partial_px = 0;
    localparam SWITCH_CMD = 2;   // 本 TB 命令的切换次数（帧 3、帧 5）

    // 采样点：彩色柱区（x=700 蓝、x=880 绿），y=600/650 避开 OSD/目标框
    task sample_px;
        input integer sx, sy;
        begin
            if ((x == sx) && (y == sy)) begin
                if (!(osd_data[23:16] == osd_data[15:8] &&
                      osd_data[15:8]  == osd_data[7:0]))
                    col_smpl = col_smpl + 1;
            end
        end
    endtask

    always @(posedge pix_clk) begin
        if (pix_rst && osd_valid) begin
            if (osd_sof) sof_cnt = sof_cnt + 1;
            if (osd_eol) begin
                x = 0; y = y + 1;
            end else begin
                x = x + 1;
            end
            sample_px(700, 600);
            sample_px(880, 650);
            px_cnt = px_cnt + 1;
            if (osd_eol) eol_cnt = eol_cnt + 1;

            // ---- 严格帧 scoreboard（按"本跨度如何结束"判定） ----
            if (osd_sof) begin
                if (sb_active) begin
                    sb_frames = sb_frames + 1;
                    if (sb_px_line == 0) begin                 // SB2 非切换跨度：必须完整
                        if ((sb_lines != 720) || (sb_px_total != FRAME_PX)) begin
                            sb_fail = sb_fail + 1;
                            $display("SB_MISM frame=%0d lines=%0d px=%0d（非切换跨度不完整）",
                                     sb_frames, sb_lines, sb_px_total);
                        end
                    end else begin                              // SB3 切换跨度：旧帧截断
                        sb_partial = sb_partial + 1;
                        sb_partial_px = sb_px_line;
                        if (sb_lines > 719) begin
                            sb_fail = sb_fail + 1;
                            $display("SB_MISM frame=%0d lines=%0d（切换跨度行数非法）",
                                     sb_frames, sb_lines);
                        end
                        if ((sb_px_total != sb_lines*1280 + sb_px_line)) begin
                            sb_fail = sb_fail + 1;              // 恒真性自检（防计数错）
                            $display("SB_MISM frame=%0d px=%0d != lines*1280+partial",
                                     sb_frames, sb_px_total);
                        end
                    end
                end
                sb_active = 1;
                sb_lines = 0; sb_px_line = 0; sb_px_total = 0;
            end
            sb_px_line  = sb_px_line + 1;
            sb_px_total = sb_px_total + 1;
            if (osd_eol) begin
                if (sb_active && (sb_px_line != 1280)) begin          // SB1 行宽严格
                    sb_fail = sb_fail + 1;
                    $display("SB_MISM line_len=%0d（frame=%0d line=%0d）",
                             sb_px_line, sb_frames+1, sb_lines);
                end
                sb_lines = sb_lines + 1;
                sb_px_line = 0;
            end

            if (px_cnt == FRAME_PX) begin    // 一个输出帧窗闭合
                win = win + 1;
                if (win > 2) begin           // 跳过启动窗
                    eval_wins = eval_wins + 1;
                    if (sof_cnt == 0) begin                          // 关键判据
                        $display("MISM win=%0d sof=0（帧首丢失！）", win);
                        mism = mism + 1;
                    end
                    if (sof_cnt > 2) begin
                        $display("MISM win=%0d sof=%0d（>2）", win, sof_cnt);
                        mism = mism + 1;
                    end
                    if (sof_cnt == 2) switch_wins = switch_wins + 1;
                    if (eol_cnt < 719) begin
                        $display("MISM win=%0d eol=%0d（<719）", win, eol_cnt);
                        mism = mism + 1;
                    end
                    if (col_smpl == 0) gray_wins = gray_wins + 1;
                    else               color_wins = color_wins + 1;
                    $display("WIN %0d: sof=%0d eol=%0d 彩色采样=%0d", win, sof_cnt, eol_cnt, col_smpl);
                end
                px_cnt = 0; sof_cnt = 0; eol_cnt = 0; col_smpl = 0;
                y = 0;
            end
        end
    end

    // ---------------- 终判 ----------------
    integer guard;
    initial begin
        repeat (5) @(posedge pix_clk);
        pix_rst = 1;
        repeat (4) @(posedge pwm_clk);
        pwm_rst = 1;
        guard = 1;
    end

    always @(posedge pix_clk) begin
        if (pix_rst && guard && (in_frame >= 11) && (win >= 8)) begin
            if (mism == 0 && sb_fail == 0 && (sb_partial <= SWITCH_CMD) &&
                (sb_frames >= 6) && gray_wins >= 2 && color_wins >= 2 && switch_wins >= 1)
                $display("TB_RESULT: PASS windows=%0d gray=%0d color=%0d switch=%0d | SB frames=%0d partial=%0d(px=%0d) sb_fail=0",
                         eval_wins, gray_wins, color_wins, switch_wins,
                         sb_frames, sb_partial, sb_partial_px);
            else
                $display("TB_RESULT: FAIL mism=%0d gray=%0d color=%0d switch=%0d | SB frames=%0d partial=%0d sb_fail=%0d",
                         mism, gray_wins, color_wins, switch_wins,
                         sb_frames, sb_partial, sb_fail);
            $finish;
        end
    end

    initial begin
        #300000000;
        $display("TB_RESULT: FAIL timeout in_frame=%0d win=%0d mism=%0d",
                 in_frame, win, mism);
        $finish;
    end
endmodule
\end{Verbatim}
\subsection{tb\_motion\_detect.v}
\begin{Verbatim}[numbers=left,numbersep=5pt,frame=single,fontsize=\footnotesize,breaklines=true,breakanywhere=true]
// =============================================================================
// tb_motion_detect.v —— 帧差运动检测位精确校验（64x48 x 5 帧）
// 场景：f0 静止背景(填缓存) -> f1 出现8x8块 -> f2 移动 -> f3 消失 -> f4 5px亚门限
//       -> f5 仅末像素 (63,47) 差分（帧末锁存边界）
// 期望：golden_model.frame_diff_motion 位精确生成（motion_exp.txt，
//       每行 15 位 hex = {cnt[13:0], valid, minx[10:0], maxx[10:0], miny[9:0], maxy[9:0]}）
// 全部 5 帧各有一条期望（f0 = 全 0：prev_ready 未置位，输出保持复位值）。
// 采样方法：O_mot_valid 仅 1 拍，直接抓会竞态。驱动器在每个帧末拍置 frame_end，
//       监视器经 frame_end_d 同步后延时读取输出总线快照（snap 每拍 NBA 采样）。
//       注：motion_detect v3（ERAM 同步读流水）输出脉冲比 v2 后移 1 拍，
//       采样窗口相应 +1 拍。
// =============================================================================
`timescale 1ns/1ps
module tb_motion_detect;
    localparam W = 64, H = 48, TH = 12, FRAMES = 6, TOTAL = W*H*FRAMES;
    reg clk = 0, rst_n = 0;
    always #7 clk = ~clk;

    reg        in_valid, in_sof, in_eol, in_mask;
    wire       o_v;
    wire [13:0] o_cnt;
    wire [10:0] o_mnx, o_mxx;
    wire [9:0]  o_mny, o_mxy;

    motion_detect #(.IMG_WIDTH(W), .IMG_HEIGHT(H), .MOT_TH(TH)) dut (
        .I_clk(clk), .I_rst_n(rst_n),
        .I_valid(in_valid), .I_sof(in_sof), .I_eol(in_eol), .I_mask(in_mask),
        .O_mot_valid(o_v), .O_mot_cnt(o_cnt),
        .O_mot_minx(o_mnx), .O_mot_maxx(o_mxx),
        .O_mot_miny(o_mny), .O_mot_maxy(o_mxy));

    reg        vin  [0:TOTAL-1];
    reg [56:0] vexp [0:FRAMES-1];
    reg [56:0] snap;
    integer    i = 0, x = 0, y = 0, mism = 0, f = 0;
    integer    pending = -1, pend_r = -1, cmp_cnt = 0;
    reg        frame_end;
    reg        x_sof, x_eol;

    // 输出总线每拍快照（NBA）
    always @(posedge clk) snap <= {o_cnt, o_v, o_mnx, o_mxx, o_mny, o_mxy};

    // 单帧比对任务
    // snap 位段：cnt=[56:43] valid=[42] minx=[41:31] maxx=[30:20] miny=[19:10] maxy=[9:0]
    task check_frame;
        input integer idx;
        begin
            if (snap !== vexp[idx]) begin
                if (mism < 8)
                    $display("MISM frame=%0d got=%h exp=%h", idx, snap, vexp[idx]);
                mism = mism + 1;
            end else begin
                $display("MEAS f%0d: cnt=%0d valid=%b bbox=(%0d,%0d)-(%0d,%0d)",
                         idx, snap[56:43], snap[42], snap[41:31], snap[19:10],
                         snap[30:20], snap[9:0]);
            end
        end
    endtask

    initial begin
        $readmemh("motion_in.hex", vin);
        $readmemh("motion_exp.txt", vexp);
    end

    // 驱动：1 像素/拍，sof/eol 按几何；帧末拍置 frame_end/pending
    always @(posedge clk) begin
        if (rst_n && (i < TOTAL+1)) begin           // 多喂 1 个冲刷令牌
            in_valid <= 1'b1;
            in_mask  <= (i == TOTAL) ? 1'b0 : vin[i];
            x_sof    = (i % (W*H) == 0);
            x_eol    = (x == W-1);
            in_sof   <= x_sof;
            in_eol   <= x_eol;
            if (x == W-1) begin
                x <= 0;
                y <= y + 1;
            end else begin
                x <= x + 1;
            end
            i = i + 1;
            frame_end <= (i % (W*H) == 0) && (i <= TOTAL);
            if (i % (W*H) == 0)
                pending = i/(W*H) - 1;
        end else begin
            in_valid <= 1'b0; in_sof <= 1'b0; in_eol <= 1'b0;
            frame_end <= 1'b0;
        end
    end

    // 监视器：frame_end_d 同步 -> 延时 2 拍 -> 读快照比对
    reg frame_end_d = 1'b0;
    always @(posedge clk) begin
        frame_end_d <= frame_end;
        if (frame_end_d) begin
            pend_r  <= pending;
            cmp_cnt <= 3;   // v3.1 帧末锁存再延一拍：脉冲在 frame_end_d 后第 3 拍
        end else if (cmp_cnt != 0) begin
            cmp_cnt <= cmp_cnt - 1;
            if (cmp_cnt == 1) begin
                check_frame(pend_r);
                f = f + 1;
                if (f == FRAMES) begin
                    if (mism == 0)
                        $display("TB_RESULT: PASS (frames=%0d)", FRAMES);
                    else
                        $display("TB_RESULT: FAIL mism=%0d/%0d", mism, FRAMES);
                    $finish;
                end
            end
        end
    end

    initial begin
        repeat (3) @(posedge clk);
        rst_n <= 1;
        #20000000;
        $display("TB_RESULT: FAIL timeout f=%0d/%0d", f, FRAMES);
        $finish;
    end
endmodule
\end{Verbatim}
\section{嵌入式对接 RTL 与 TB}
\subsection{system\_top.v}
\fileintro{嵌入式集成顶层（仿真/TD 双用）：官方 ISP 全链 + vision\_top 挂载 + 按键消抖 glue + UART 遥测。2026-10-03 修复（\#N+7）：视觉抽头改 AWB 96bit 直抽（isp\_top 新增 tap\_* 端口；弃用 128bit 重排——其 tuser 与 valid 结构性错位）；I\_tvalid 直用 RAW 有效、O\_tready 显式 1'b1（原悬空网）。}
\begin{Verbatim}[numbers=left,numbersep=5pt,frame=single,fontsize=\footnotesize,breaklines=true,breakanywhere=true]
// =============================================================================
// system_top.v —— 嵌入式方向系统集成顶层（仿真/TD 双用）
// -----------------------------------------------------------------------------
// 挂载关系（对应 06_interfaces/current 三份接口文档）：
//   虚拟/真实 RAW10 流 ─> isp_top（官方，CSI 解包在 PHY 侧已完成，此处从
//                          raw10_unpacket 输出界接入）─> 128bit DDR 写流
//   128bit ─data128_96(官方重排)─> 96bit ─> vision_top（视觉目标包/OSD/云台）
//   按键消抖 ─> 阈值组切换；UART 遥测输出目标状态
// 与官方例程的差异（全部为"增加"，无官方文件改动）：
//   1. 不含 DDR2/MIPI-PHY/HDMI 加密核（仿真以帧转储替代 HDMI；DDR 写流空置）
//   2. 新增 uart_telemetry 遥测
// 时钟：单域 axi4s_video_aclk（74.25M 级）；pwm_clk 50M 由引脚给入
// =============================================================================
module system_top #(
    parameter IMG_WIDTH  = 1280,
    parameter IMG_HEIGHT = 720,
    parameter CLK_DIV_UART = 646          // 74.25MHz/115200 ~= 645
)(
    input  wire        I_clk,             // axi4s_video_aclk
    input  wire        I_rst_n,
    input  wire        I_pwm_clk,         // 50MHz
    input  wire        I_pwm_rst_n,
    // 虚拟/真实 RAW10 传感器流（真实板 = raw10_unpacket_2lane 输出）
    input  wire [39:0] I_raw10_data,
    input  wire        I_raw10_valid,
    input  wire        I_raw10_sof,
    input  wire        I_raw10_eol,
    // 按键（低有效，经消抖）
    input  wire [1:0]  I_btn_n,
    input  wire        I_algo_edge,
    input  wire [6:0]  I_bright,
    input  wire [7:0]  I_contrast,
    output wire        O_led_g,
    output wire        O_led_r,
    output wire        O_buzzer,
    // 云台 PWM
    output wire        O_pwm_x,
    output wire        O_pwm_y,
    output wire [12:0] O_duty_x_us,
    output wire [12:0] O_duty_y_us,
    output wire [15:0] O_lost_cnt,
    // UART 遥测
    output wire        O_uart_txd,
    // 调试观察口（外脑评审：OSD 流接口补全 sof/eol，供显示侧正式消费）
    output wire [23:0] O_osd_data,
    output wire        O_osd_valid,
    output wire        O_osd_sof,
    output wire        O_osd_eol
);

    // ---------------- 官方 ISP（RAW10 -> RGB） ----------------
    // 2026-10-03 修复（迭代记录 #N+6）：
    //  (1) 握手：isp_top.O_tready 是"下游就绪输入"（官方 wrapper 也悬空该网，
    //      有效行为=无背压）。原代码把该浮动网串进 I_tvalid —— 仿真为 Z、综合后
    //      会毁掉输入使能。本整合顶层无显示主路，显式恒 1（与有效行为一致）；
    //      I_tvalid 直接用 RAW 有效（与官方 design_top_wrapper 同款接法）。
    //  (2) 视觉抽头改 AWB 96bit（tap_*，4px RGB888，tuser/tlast 实测正确）；
    //      弃用 128bit 重排路径——其 tuser 与 valid 结构性错位（frame_start 直通
    //      延迟 1 拍 vs valid 受 4:3 门控，实测 tuser&&valid 恒为 0）。
    wire [127:0] isp_tdata;
    wire         isp_tvalid, isp_tuser, isp_tlast;
    wire         isp_up_tready;
    wire [95:0]  tap_tdata;
    wire         tap_tvalid, tap_tuser, tap_tlast;

    isp_top u_isp (
        .axi4s_video_aclk (I_clk),
        .I_rst_n          (I_rst_n),
        .I_tlast          (I_raw10_eol),
        .I_tuser          (I_raw10_sof),
        .I_tdata          (I_raw10_data),
        .I_tvalid         (I_raw10_valid),
        .I_tdest          (10'd0),
        .O_tready         (1'b1),            // 无显示主路，恒就绪（原悬空网已修）
        .O_tdata          (isp_tdata),
        .O_tlast          (isp_tlast),
        .O_tuser          (isp_tuser),
        .O_tvalid         (isp_tvalid),
        .I_tready         (isp_up_tready),   // 对上游就绪指示（输出）
        .tap_tdata        (tap_tdata),
        .tap_tvalid       (tap_tvalid),
        .tap_tuser        (tap_tuser),
        .tap_tlast        (tap_tlast)
    );

    // ---------------- 按键消抖（官方模块）-> 阈值组/跟踪使能 ----------------
    wire key0_trig, key1_trig;
    key_remove_shakes u_key0 (.I_clk(I_clk), .I_rst_n(I_rst_n),
        .I_key_in(~I_btn_n[0]), .O_key_trig_out(key0_trig));
    key_remove_shakes u_key1 (.I_clk(I_clk), .I_rst_n(I_rst_n),
        .I_key_in(~I_btn_n[1]), .O_key_trig_out(key1_trig));

    reg [2:0]  preset_idx  = 3'd0;
    reg        track_en    = 1'b1;
    always @(posedge I_clk or negedge I_rst_n) begin
        if (!I_rst_n) begin
            preset_idx <= 3'd0;
            track_en   <= 1'b1;
        end else begin
            if (key0_trig) preset_idx <= preset_idx + 3'd1;   // KEY0: 切阈值组
            if (key1_trig) track_en   <= ~track_en;           // KEY1: 跟踪使能
        end
    end

    // ---------------- 外部输入 2FF 同步 ----------------
    // 外脑评审：按键/拨码为异步引脚输入，直接使用存在亚稳风险
    // （尤其 I_algo_edge 在帧首采样）。每输入两组触发器同步。
    reg [1:0] syn_algo, syn_bright0, syn_bright1, syn_ctr0, syn_ctr1;
    always @(posedge I_clk or negedge I_rst_n) begin
        if (!I_rst_n) begin
            syn_algo <= 2'b00; syn_bright0 <= 7'd0; syn_bright1 <= 7'd0;
            syn_ctr0 <= 8'd64; syn_ctr1 <= 8'd64;
        end else begin
            syn_algo    <= {syn_algo[0], I_algo_edge};
            syn_bright0 <= I_bright;      syn_bright1 <= syn_bright0;
            syn_ctr0    <= I_contrast;    syn_ctr1    <= syn_ctr0;
        end
    end
    wire        algo_s     = syn_algo[1];
    wire [6:0]  bright_s   = syn_bright1;
    wire [7:0]  contrast_s = syn_ctr1;

    // ---------------- 视觉方向 vision_top（检测/OSD/云台） ----------------
    wire [95:0]  vis_pt_data;   // 视觉直通（本顶层不使用）
    wire         vis_pt_valid, vis_pt_user, vis_pt_last, vis_tready;
    wire         osd_valid, osd_sof, osd_eol;
    wire [23:0]  osd_data;
    wire [12:0]  duty_x, duty_y;
    wire [15:0]  lost_cnt;
    wire [7:0]   status;
    // 遥测用目标状态（外脑评审: 端口化替代层次化引用, 可综合可约束）
    wire         frame_tick, target_valid;
    wire [10:0]  tgt_cx, tgt_cy;

    vision_top #(
        .IMG_WIDTH(IMG_WIDTH), .IMG_HEIGHT(IMG_HEIGHT)
    ) u_vision (
        .I_pix_clk    (I_clk),
        .I_pix_rst_n  (I_rst_n),
        .I_pwm_clk    (I_pwm_clk),
        .I_pwm_rst_n  (I_pwm_rst_n),
        .I_tdata      (tap_tdata),
        .I_tvalid     (tap_tvalid),
        .I_tuser      (tap_tuser),
        .I_tlast      (tap_tlast),
        .O_tready     (vis_tready),
        .O_tdata      (vis_pt_data),
        .O_tvalid     (vis_pt_valid),
        .O_tuser      (vis_pt_user),
        .O_tlast      (vis_pt_last),
        .O_osd_valid  (osd_valid),
        .O_osd_sof    (osd_sof),
        .O_osd_eol    (osd_eol),
        .O_osd_data   (osd_data),
        .I_preset_idx (preset_idx),
        .I_track_en   (track_en),
        .I_algo_edge  (algo_s),
        .I_bright     (bright_s),
        .I_contrast   (contrast_s),
        .O_pwm_x      (O_pwm_x),
        .O_pwm_y      (O_pwm_y),
        .O_duty_x_us  (O_duty_x_us),
        .O_duty_y_us  (O_duty_y_us),
        .O_lost_cnt   (O_lost_cnt),
        .O_status     (status),
        .O_led_g      (O_led_g),
        .O_led_r      (O_led_r),
        .O_buzzer     (O_buzzer),
        .O_frame_tick (frame_tick),
        .O_target_valid (target_valid),
        .O_tgt_cx     (tgt_cx),
        .O_tgt_cy     (tgt_cy)
    );
    assign O_osd_valid = osd_valid;
    assign O_osd_data  = osd_data;
    assign O_osd_sof   = osd_sof;
    assign O_osd_eol   = osd_eol;
    // vis_tready 由 vision_top 内部恒 1 驱动（窃流不反压），
    // 顶层不得再 assign 同一网络（外脑评审: 多驱动会在 TD 综合报错）。

    // ---------------- UART 遥测（新增 RTL；信号经 vision_top 端口引出） ----------------
    // 外脑评审修复：原用层次化引用 u_vision.t_* 取跨域信号——TD 综合下
    // 层次引用不可靠且无法施加 CDC 约束，现改走 vision_top 显式端口。
    uart_telemetry #(.CLK_DIV(CLK_DIV_UART)) u_uart (
        .I_clk          (I_clk),
        .I_rst_n        (I_rst_n),
        .I_frame_tick   (frame_tick),
        .I_target_valid (target_valid),
        .I_cx           (tgt_cx),
        .I_cy           (tgt_cy),
        .I_preset       (preset_idx),
        .O_txd          (O_uart_txd)
    );

    // 视觉内部信号已端口化（见 vision_top O_frame_tick/O_target_valid/O_tgt_cx/cy）
endmodule
\end{Verbatim}
\subsection{uart\_telemetry.v}
\fileintro{目标状态 UART 遥测（115200-8N1，每帧一包 9 字节）：A5|valid|cx(2)|cy(2)|preset|sum|5A。帧拍起包（发送中跳过）；单 always 状态机（外脑评审修复：原两块同驱寄存器）。}
\begin{Verbatim}[numbers=left,numbersep=5pt,frame=single,fontsize=\footnotesize,breaklines=true,breakanywhere=true]
// =============================================================================
// uart_telemetry.v —— 目标状态 UART 遥测（115200-8N1，每帧一包）
// -----------------------------------------------------------------------------
// 帧格式（9 字节）：A5 | valid | cx_lo cx_hi | cy_lo cy_hi | preset | checksum | 5A
// checksum = 低 8 位累加(A5..preset)。上位机按 0xA5 ... 0x5A 定界。
// 波特率 = I_clk / CLK_DIV（74.25MHz/646 ~= 114938，标准容差内）。
// 发送期间新帧到达则跳过（下一帧再发），绝不打断当前包。
// =============================================================================
module uart_telemetry #(
    parameter integer CLK_DIV = 646
)(
    input  wire        I_clk,
    input  wire        I_rst_n,
    input  wire        I_frame_tick,     // 每帧一拍（检测域）
    input  wire        I_target_valid,
    input  wire [10:0] I_cx,
    input  wire [10:0] I_cy,
    input  wire [2:0]  I_preset,
    output reg         O_txd
);
    localparam NBYTES = 9;
    reg [7:0] payload [0:NBYTES-1];
    reg [3:0] byte_idx;
    reg [3:0] bit_idx;                    // 0=起始 1..8=数据 9=停止
    reg [9:0] div_cnt;
    reg       sending;
    reg [7:0] shifter;
    reg       tick_d;

    // 快照 + 发送状态机（单一 always —— 外脑评审: 原两块同驱
    // byte_idx/bit_idx/div_cnt/sending 属多驱动寄存器，TD 综合会报错；
    // 合并后时序行为与原设计逐拍一致）
    always @(posedge I_clk or negedge I_rst_n) begin
        if (!I_rst_n) begin
            O_txd   <= 1'b1;
            tick_d <= 1'b0;
            sending <= 1'b0;
            div_cnt <= 10'd0;
            bit_idx <= 4'd0;
            byte_idx <= 4'd0;
            shifter <= 8'hFF;
        end else begin
            tick_d <= I_frame_tick;
            if (sending) begin
            if (div_cnt == CLK_DIV - 1) begin
                div_cnt <= 10'd0;
                case (bit_idx)
                    4'd0: begin
                        O_txd   <= 1'b0;                 // 起始位
                        shifter <= payload[byte_idx];    // 装载
                    end
                    4'd9: begin
                        O_txd <= 1'b1;                   // 停止位
                        if (byte_idx == NBYTES - 1)
                            sending <= 1'b0;
                        else
                            byte_idx <= byte_idx + 4'd1;
                    end
                    default: begin                       // 数据位 LSB 先行
                        O_txd   <= shifter[0];
                        shifter <= {1'b0, shifter[7:1]};
                    end
                endcase
                bit_idx <= (bit_idx == 4'd9) ? 4'd0 : bit_idx + 4'd1;
            end else begin
                div_cnt <= div_cnt + 10'd1;
            end
        end else begin
            O_txd <= 1'b1;
            // 帧拍上升沿: 快照目标状态并启动一包发送
            if (I_frame_tick && !tick_d) begin
                payload[0] <= 8'hA5;
                payload[1] <= {7'd0, I_target_valid};
                payload[2] <= I_cx[7:0];
                payload[3] <= {5'd0, I_cx[10:8]};
                payload[4] <= I_cy[7:0];
                payload[5] <= {5'd0, I_cy[10:8]};
                payload[6] <= {5'd0, I_preset};
                payload[7] <= 8'hA5 + {7'd0, I_target_valid} + I_cx[7:0]
                            + {5'd0, I_cx[10:8]} + I_cy[7:0] + {5'd0, I_cy[10:8]}
                            + {5'd0, I_preset};
                payload[8] <= 8'h5A;
                byte_idx <= 4'd0;
                bit_idx  <= 4'd0;
                div_cnt  <= 10'd0;
                sending  <= 1'b1;
            end
        end
    end
end

endmodule
\end{Verbatim}
\subsection{i2c\_bfm\_sc500.v}
\fileintro{SC500CS I2C 从机 BFM（仿真）：开漏建模（inout + 上拉）、7bit 地址 0x36、ACK 位相、每 4 字节一个寄存器写事务。已知遗留：ACK 位相未收敛（列 TD 阶段项，需真机逻辑分析仪对照）——寄存器序列本身已由 ROM 侧逐条确认。}
\begin{Verbatim}[numbers=left,numbersep=5pt,frame=single,fontsize=\footnotesize,breaklines=true,breakanywhere=true]
// =============================================================================
// i2c_bfm_sc500.v —— SC500CS I2C 从机 BFM（仅供仿真）
// -----------------------------------------------------------------------------
// 行为级开漏 I2C 从机：7bit 地址 0x36（8bit 写 0x6C），记录主控写入字节序列。
// 每 4 字节为一次寄存器写事务：{0x6C, reg_hi, reg_lo, data}。
// 时序：数据位在 SCL 高电平期稳定，下降沿采样；ACK 在第 9 个 SCL 上升沿拉低、
// 第 9 个下降沿释放；START/STOP 在 SCL 高电平期按 SDA 边沿识别。
// =============================================================================
module i2c_bfm_sc500 #(
    parameter integer MAX_BYTES = 2048
)(
    input  wire I_clk,
    input  wire I_rst_n,
    input  wire I_scl,
    inout  wire I_sda,
    output reg  [15:0] O_byte_cnt,
    output reg         O_addr_seen,
    output reg         O_addr_naked     // 地址不匹配 NAK 事件
);
    // 记录数组放在模块内（V2001 端口不允许数组），TB 层次化读取
    reg [7:0] O_bytes [0:MAX_BYTES-1];
    reg       pull;                    // BFM 拉低使能
    reg [3:0] bit_cnt;                 // 0..7 数据位，8=ACK 位
    reg [7:0] shift;
    reg       addr_match;
    reg       scl_d, sda_d;            // 上一拍值
    wire      scl_rise = I_scl && !scl_d;
    wire      scl_fall = !I_scl && scl_d;
    wire      start_c  = I_scl &&  sda_d && !I_sda;
    wire      stop_c   = I_scl && !sda_d &&  I_sda;
    wire      shift_next_match = {shift[6:0], I_sda} == 8'h6C;

    assign I_sda = pull ? 1'bz : 1'b0; // 开漏
    wire sda_now = I_sda;              // 采样（z 由 TB 上拉为 1）

    always @(posedge I_clk or negedge I_rst_n) begin
        if (!I_rst_n) begin
            pull <= 1'b1; bit_cnt <= 0; shift <= 0; addr_match <= 0;
            scl_d <= 1'b1; sda_d <= 1'b1;
            O_byte_cnt <= 16'd0; O_addr_seen <= 1'b0; O_addr_naked <= 1'b0;
        end else begin
            scl_d <= I_scl;
            sda_d <= sda_now;

            if (start_c) begin
                bit_cnt <= 4'd0;
                addr_match <= 1'b0;
                pull <= 1'b1;
            end else if (stop_c) begin
                addr_match <= 1'b0;
                pull <= 1'b1;
            end else if (scl_rise) begin
                if (bit_cnt == 4'd8) begin
                    // 第 9 个时钟上升沿：地址匹配或数据字节 -> ACK
                    if (!addr_match) begin
                        // synthesis translate_off
                        $display("BFM_ADDR shift=%h bits=%0d t=%0t", shift, bit_cnt, $time);
                        // synthesis translate_on
                        if (shift[7:1] == 7'h36) begin
                            addr_match  <= 1'b1;
                            O_addr_seen <= 1'b1;
                            pull        <= 1'b0;       // ACK
                        end else begin
                            O_addr_naked <= 1'b1;
                            pull         <= 1'b1;      // NAK
                        end
                    end else begin
                        pull <= 1'b0;                  // 数据 ACK
                    end
                end
            end else if (scl_fall) begin
                if (bit_cnt <= 4'd7) begin
                    shift  <= {shift[6:0], sda_now};   // 采样数据位
                    bit_cnt <= bit_cnt + 4'd1;
                end else begin
                    // 第 9 个下降沿：释放；地址后的字节写入记录
                    pull <= 1'b1;
                    if (addr_match && O_byte_cnt < MAX_BYTES) begin
                        O_bytes[O_byte_cnt[15:0]] <= shift;
                        O_byte_cnt <= O_byte_cnt + 16'd1;
                    end
                    bit_cnt <= 4'd0;
                end
            end
        end
    end

endmodule
\end{Verbatim}
\subsection{divider\_gate\_sim.v}
\fileintro{门级除法器（AL\_* 库 7355 行网表）的行为级替身（仅仿真）：16 拍延迟、除零给全 1。TD 综合时用真网表，本文件排除。}
\begin{Verbatim}[numbers=left,numbersep=5pt,frame=single,fontsize=\footnotesize,breaklines=true,breakanywhere=true]

// =============================================================================
// divider_gate_sim.v —— divider_gate 门级网表的行为级替身（仅供仿真）
// TD 综合时排除本文件，使用官方 divider_gate.v（AL_* 库单元）。
// 行为：start 上升沿锁存，16 拍后 done 一拍脉冲，quotient = num/den（floor）。
// =============================================================================
`timescale 1ns/1ps
module divider (
    input  wire        clk,
    input  wire [31:0] denominator,
    input  wire [47:0] numerator,
    input  wire        rst,
    input  wire        start,
    output reg         done,
    output reg  [47:0] quotient,
    output reg  [31:0] remainder
);
    reg [47:0] num_q;
    reg [31:0] den_q;
    reg [5:0]  wait_cnt;
    reg busy;

    always @(posedge clk or posedge rst) begin
        if (rst) begin
            done <= 1'b0; busy <= 1'b0;
            quotient <= 0; remainder <= 0;
            num_q <= 0; den_q <= 0; wait_cnt <= 0;
        end else begin
            done <= 1'b0;
            if (!busy && start && !done) begin
                num_q <= numerator;
                den_q <= denominator;
                busy  <= 1'b1;
                wait_cnt <= 0;
            end else if (busy) begin
                if (wait_cnt == 6'd15) begin
                    if (den_q != 0) begin
                        quotient  <= num_q / den_q;
                        remainder <= num_q % den_q;
                    end else begin
                        quotient <= {48{1'b1}};
                        remainder <= num_q;
                    end
                    done <= 1'b1;
                    busy <= 1'b0;
                end else begin
                    wait_cnt <= wait_cnt + 6'd1;
                end
            end
        end
    end
endmodule
\end{Verbatim}
\subsection{ph1p\_logic\_eram\_stub.v}
\fileintro{TD ERAM 原语行为替身 v3（仅仿真）：参数名与官方 wrapper 对齐（ADDR\_WIDTH\_A/B，曾因名字不匹配致端口宽度失效、整帧 X）；地址 Z/X 位清洗归 0 + 未知位单次告警；使能语义：cea/ceb 未接=通、wea/web 仅显式 1 才写；mem 上电清零。}
\begin{Verbatim}[numbers=left,numbersep=5pt,frame=single,fontsize=\footnotesize,breaklines=true,breakanywhere=true]
// =============================================================================
// ph1p_logic_eram_stub.v —— TD ERAM 原语的行为级替身（仅供仿真，TD 综合时用真原语）
// -----------------------------------------------------------------------------
// 端口/参数清单来自官方 blk_mem 生成 RTL 的实例化（addra/addrb/dia/dib/cea/ceb/
// clka/clkb/wea/web/bea/beb/rsta/rstb/doa/dob + ecc 端口）。
// 行为：真双口 RAM，A 口写、B 口读写；bea/beb=0 表示全字节写（忽略）。
//
// v3 修正（2026-10-03，视频通路整链仿真"全 X"事故）：
//   [1] 参数名与官方 wrapper 对齐：wrapper 传的是 ADDR_WIDTH_A/ADDR_WIDTH_B
//       （旧版端口宽度误用 ADDRA_WIDTH/ADDRB_WIDTH，wrapper 传入的宽度
//       没作用到端口上——端口 14bit 对 11bit 连接，高位悬空 Z）。
//   [2] 地址位清洗：端口宽度仍可能大于实际连接位数（formal 高位 Z），
//       mem 索引前把 Z/X 位归 0（真实硬件地址不可能未知）。
//   [3] 使能/写使能语义对齐硬件：cea/ceb 未接（Z）= 直通使能；
//       wea/web 仅显式 1 才写（wrapper 有把 web 悬空的只读用法，
//       若把 Z 当写使能会把 Z 数据灌进 mem——v2 中间版曾犯此错）。
//   [4] mem 上电清零：对齐硬件初值 0，防止 X 沿 AWB 增益累加路径
//       污染整帧（"X 污染锁存"）。
// =============================================================================
`timescale 1ns/1ps
module PH1P_LOGIC_ERAM #(
    parameter integer DATA_WIDTH_A = 20,
    parameter integer DATA_WIDTH_B = DATA_WIDTH_A,
    parameter integer ADDR_WIDTH_A = 14,
    parameter integer ADDR_WIDTH_B = ADDR_WIDTH_A,
    parameter MODE               = "DP",
    parameter IMPLEMENT          = "20KFAST",
    parameter REGMODE_A          = "NORESET",
    parameter REGMODE_B          = "NORESET",
    parameter RESETMODE_A        = "SYNC",
    parameter RESETMODE_B        = "SYNC",
    parameter OREGSET_A          = "RESET",
    parameter OREGSET_B          = "RESET",
    parameter SYNC_RESET_RELEASE_A = "DISABLE",
    parameter SYNC_RESET_RELEASE_B = "DISABLE",
    parameter ASYNC_RESET_RELEASE_A = "ASYNC",
    parameter ASYNC_RESET_RELEASE_B = "ASYNC",
    parameter SSROVERCE          = "DISABLE",
    parameter ECC_ENCODE         = "DISABLE",
    parameter ECC_DECODE         = "DISABLE",
    parameter GSR                = "ENABLE",
    parameter RESETMODE          = "SYNC",
    parameter CLKMODE            = "ASYNC",
    parameter WRITEMODE_A        = "NORMAL",
    parameter WRITEMODE_B        = "NORMAL",
    parameter DATA_DEPTH_A       = 512,
    parameter DATA_DEPTH_B       = 512,
    parameter INIT_FILE          = "NONE",
    parameter FILL_ALL           = "NONE"
)(
    input  wire                     clka,
    input  wire                     clkb,
    input  wire [ADDR_WIDTH_A-1:0]  addra,
    input  wire [ADDR_WIDTH_B-1:0]  addrb,
    input  wire [DATA_WIDTH_A-1:0]  dia,
    input  wire [DATA_WIDTH_B-1:0]  dib,
    input  wire                     cea,
    input  wire                     ceb,
    input  wire                     ocea,
    input  wire                     oceb,
    input  wire                     wea,
    input  wire                     web,
    input  wire                     bea,
    input  wire                     beb,
    input  wire                     rsta,
    input  wire                     rstb,
    output reg  [DATA_WIDTH_A-1:0]  doa,
    output reg  [DATA_WIDTH_B-1:0]  dob,
    output wire                     ecc_sbiterr,
    output wire                     ecc_dbiterr,
    input  wire                     ecc_sbiterrinj,
    input  wire                     ecc_dbiterrinj
);
    localparam DEPTH_A = (DATA_DEPTH_A > 16384) ? 16384 : DATA_DEPTH_A;
    localparam DEPTH_B = (DATA_DEPTH_B > 16384) ? 16384 : DATA_DEPTH_B;
    reg [DATA_WIDTH_A-1:0] mem [0:DEPTH_A-1];
    // 硬件对齐：TD ERAM 上电读出为 0（v3[4]，仅供仿真）
    integer _i;
    initial for (_i = 0; _i < DEPTH_A; _i = _i + 1) mem[_i] = {DATA_WIDTH_A{1'b0}};
    assign ecc_sbiterr = 0; assign ecc_dbiterr = 0;
    // 外脑评审: 未知位(X/Z)单次告警——替身对悬空地址/使能保持宽容语义，
    // 但必须让"连接错误"在仿真阶段可见（只报一次，避免稳定悬空刷屏）。
    reg warned_unknown;
    initial warned_unknown = 1'b0;
    always @(posedge clka or posedge clkb) begin
        if (!warned_unknown &&
            (^{addra, addrb, cea, ceb, wea, web, rsta, rstb} === 1'bx)) begin
            warned_unknown <= 1'b1;
            $display("ERAM_STUB_WARN: unknown (X/Z) on addr/enable pins @ %m (t=%0t)", $time);
        end
    end


    // 地址位清洗（v3[2]）：formal 高位可能因连接位宽不足悬空为 Z，
    // mem 索引带 Z/X 会整体返回 X——把未知位归 0（真实硬件不存在未知地址）。
    reg [ADDR_WIDTH_A-1:0] addra_c;
    reg [ADDR_WIDTH_B-1:0] addrb_c;
    integer _k;
    always @(*) begin
        for (_k = 0; _k < ADDR_WIDTH_A; _k = _k + 1)
            addra_c[_k] = (addra[_k] === 1'b1) ? 1'b1 : 1'b0;
        for (_k = 0; _k < ADDR_WIDTH_B; _k = _k + 1)
            addrb_c[_k] = (addrb[_k] === 1'b1) ? 1'b1 : 1'b0;
    end

    // 使能/写使能语义（v3[3]）：
    //   cea/ceb 未接(Z)=直通使能；rsta/rstb 仅显式 1 复位；
    //   wea/web 仅显式 1 才写（悬空 Z 的 web 不得写——只读端口用法）。
    wire cea_e = (cea !== 1'b0);
    wire ceb_e = (ceb !== 1'b0);

    always @(posedge clka) begin
        if (cea_e && (rsta === 1'b1)) begin
            doa <= 0;
        end else if (cea_e) begin
            if ((wea === 1'b1) && (addra_c < DEPTH_A)) mem[addra_c] <= dia;
            doa <= mem[addra_c];
        end
    end
    always @(posedge clkb) begin
        if (ceb_e && (rstb === 1'b1)) begin
            dob <= 0;
        end else if (ceb_e) begin
            if ((web === 1'b1) && (addrb_c < DEPTH_B)) mem[addrb_c] <= dib;
            dob <= mem[addrb_c];
        end
    end
endmodule

module PH1_LOGIC_ERAM #(
    parameter integer DATA_WIDTH_A = 20,
    parameter integer DATA_WIDTH_B = DATA_WIDTH_A,
    parameter integer ADDR_WIDTH_A = 14,
    parameter integer ADDR_WIDTH_B = ADDR_WIDTH_A,
    parameter MODE               = "DP",
    parameter IMPLEMENT          = "20KFAST",
    parameter REGMODE_A          = "NORESET",
    parameter REGMODE_B          = "NORESET",
    parameter RESETMODE_A        = "SYNC",
    parameter RESETMODE_B        = "SYNC",
    parameter OREGSET_A          = "RESET",
    parameter OREGSET_B          = "RESET",
    parameter SYNC_RESET_RELEASE_A = "DISABLE",
    parameter SYNC_RESET_RELEASE_B = "DISABLE",
    parameter ASYNC_RESET_RELEASE_A = "ASYNC",
    parameter ASYNC_RESET_RELEASE_B = "ASYNC",
    parameter SSROVERCE          = "DISABLE",
    parameter ECC_ENCODE         = "DISABLE",
    parameter ECC_DECODE         = "DISABLE",
    parameter GSR                = "ENABLE",
    parameter RESETMODE          = "SYNC",
    parameter CLKMODE            = "ASYNC",
    parameter WRITEMODE_A        = "NORMAL",
    parameter WRITEMODE_B        = "NORMAL",
    parameter DATA_DEPTH_A       = 512,
    parameter DATA_DEPTH_B       = 512,
    parameter INIT_FILE          = "NONE",
    parameter FILL_ALL           = "NONE"
)(
    input  wire                     clka,
    input  wire                     clkb,
    input  wire [ADDR_WIDTH_A-1:0]  addra,
    input  wire [ADDR_WIDTH_B-1:0]  addrb,
    input  wire [DATA_WIDTH_A-1:0]  dia,
    input  wire [DATA_WIDTH_B-1:0]  dib,
    input  wire                     cea,
    input  wire                     ceb,
    input  wire                     ocea,
    input  wire                     oceb,
    input  wire                     wea,
    input  wire                     web,
    input  wire                     bea,
    input  wire                     beb,
    input  wire                     rsta,
    input  wire                     rstb,
    output reg  [DATA_WIDTH_A-1:0]  doa,
    output reg  [DATA_WIDTH_B-1:0]  dob,
    output wire                     ecc_sbiterr,
    output wire                     ecc_dbiterr,
    input  wire                     ecc_sbiterrinj,
    input  wire                     ecc_dbiterrinj
);
    localparam DEPTH_A = (DATA_DEPTH_A > 16384) ? 16384 : DATA_DEPTH_A;
    localparam DEPTH_B = (DATA_DEPTH_B > 16384) ? 16384 : DATA_DEPTH_B;
    reg [DATA_WIDTH_A-1:0] mem [0:DEPTH_A-1];
    integer _j;
    initial for (_j = 0; _j < DEPTH_A; _j = _j + 1) mem[_j] = {DATA_WIDTH_A{1'b0}};
    assign ecc_sbiterr = 0; assign ecc_dbiterr = 0;
    // 外脑评审: 未知位(X/Z)单次告警——替身对悬空地址/使能保持宽容语义，
    // 但必须让"连接错误"在仿真阶段可见（只报一次，避免稳定悬空刷屏）。
    reg warned_unknown;
    initial warned_unknown = 1'b0;
    always @(posedge clka or posedge clkb) begin
        if (!warned_unknown &&
            (^{addra, addrb, cea, ceb, wea, web, rsta, rstb} === 1'bx)) begin
            warned_unknown <= 1'b1;
            $display("ERAM_STUB_WARN: unknown (X/Z) on addr/enable pins @ %m (t=%0t)", $time);
        end
    end


    reg [ADDR_WIDTH_A-1:0] addra_c;
    reg [ADDR_WIDTH_B-1:0] addrb_c;
    integer _m;
    always @(*) begin
        for (_m = 0; _m < ADDR_WIDTH_A; _m = _m + 1)
            addra_c[_m] = (addra[_m] === 1'b1) ? 1'b1 : 1'b0;
        for (_m = 0; _m < ADDR_WIDTH_B; _m = _m + 1)
            addrb_c[_m] = (addrb[_m] === 1'b1) ? 1'b1 : 1'b0;
    end

    wire cea_e = (cea !== 1'b0);
    wire ceb_e = (ceb !== 1'b0);

    always @(posedge clka) begin
        if (cea_e && (rsta === 1'b1)) begin
            doa <= 0;
        end else if (cea_e) begin
            if ((wea === 1'b1) && (addra_c < DEPTH_A)) mem[addra_c] <= dia;
            doa <= mem[addra_c];
        end
    end
    always @(posedge clkb) begin
        if (ceb_e && (rstb === 1'b1)) begin
            dob <= 0;
        end else if (ceb_e) begin
            if ((web === 1'b1) && (addrb_c < DEPTH_B)) mem[addrb_c] <= dib;
            dob <= mem[addrb_c];
        end
    end
endmodule
\end{Verbatim}
\subsection{tb\_uart\_telemetry.v}
\fileintro{遥测验证：注入 3 帧（valid=1/0/1、坐标各异），uart\_rx\_mon 逐字节解码，校验定界/校验和/坐标；PASS 判据：全帧正确且含 >=2 帧 LOCK。}
\begin{Verbatim}[numbers=left,numbersep=5pt,frame=single,fontsize=\footnotesize,breaklines=true,breakanywhere=true]
// =============================================================================
// tb_uart_telemetry.v —— UART 遥测验证
// 注入 3 帧（valid=1/0/1，坐标各异），解码串行字节流，校验：
//   帧定界 A5/5A、valid 位、cx/cy 全部坐标字节、校验和。
//   （2026-10-03 外脑评审补：原检查器只校 cx 不校 cy——覆盖盲点）
// =============================================================================
`timescale 1ns/1ps
module tb_uart_telemetry;
    localparam CLK_DIV = 646;
    reg clk = 0, rst_n = 0;
    always #7 clk = ~clk;

    reg         frame_tick = 0;
    reg         tv = 0;
    reg [10:0]  cx = 0, cy = 0;
    reg [2:0]   preset = 3'd2;
    wire        txd;

    uart_telemetry #(.CLK_DIV(CLK_DIV)) dut (
        .I_clk(clk), .I_rst_n(rst_n),
        .I_frame_tick(frame_tick), .I_target_valid(tv),
        .I_cx(cx), .I_cy(cy), .I_preset(preset), .O_txd(txd));

    // 激励：三帧，帧间隔 3ms（远大于一包 780µs）
    task send_frame(input v, input integer x, input integer y);
        begin
            tv = v; cx = x; cy = y;
            frame_tick = 1;
            #14;
            frame_tick = 0;
            #3000000;
        end
    endtask

    wire [7:0] rbyte;
    wire       rvalid;
    uart_rx_mon #(.CLK_DIV(CLK_DIV)) u_mon (
        .I_clk(clk), .I_rst_n(rst_n), .I_rxd(txd),
        .O_byte(rbyte), .O_valid(rvalid));

    // 组包校验
    integer state = 0, good = 0, frames = 0, mism = 0;
    reg [7:0] f [0:8];
    always @(posedge clk) begin
        if (rst_n && rvalid) begin
            $display("RXB t=%0t byte=%02x state=%0d", $time, rbyte, state);
            case (state)
                0: if (rbyte == 8'hA5) state = 1;
                8: begin
                    if (rbyte == 8'h5A) begin
                        frames = frames + 1;
                        // 校验和 + 内容
                        if (f[7] == (8'hA5 + f[1] + f[2] + f[3] + f[4] + f[5] + f[6])) begin
                            if ((f[1] == 8'h01) && (f[2] == 8'h52) && (f[3] == 8'h02) &&   // cx=594=0x252
                                (f[4] == 8'h4D) && (f[5] == 8'h01) &&                   // cy=333=0x14D（外脑评审补）
                                (f[6] == 8'h02))
                                good = good + 1;
                        end else mism = mism + 1;
                    end else mism = mism + 1;
                    state = 0;
                end
                default: begin
                    f[state] = rbyte;
                    state = state + 1;
                end
            endcase
        end
    end

    // 帧内容期望：帧0 LOST（00,0,0）不校坐标；帧1 LOCK (0254=594?) 用帧1的 (594, 333)
    // 简化：只对 LOCK 帧校验 X 低位字节与组号（帧1 注入 cx=594=0x252、cy=333）
    initial begin
        repeat (5) @(posedge clk);
        rst_n = 1;
        send_frame(0, 0, 0);
        send_frame(1, 594, 333);
        send_frame(1, 594, 333);
        #4000000;
        if ((frames >= 2) && (good >= 2) && (mism == 0))   // 首帧可能被监视器冷启动吞掉
            $display("TB_RESULT: PASS frames=%0d good=%0d", frames, good);
        else
            $display("TB_RESULT: FAIL frames=%0d good=%0d mism=%0d", frames, good, mism);
        $finish;
    end
    initial begin
        #40000000;
        $display("TB_RESULT: FAIL timeout frames=%0d good=%0d", frames, good);
        $finish;
    end
endmodule
\end{Verbatim}
\subsection{uart\_rx\_mon.v}
\fileintro{UART 接收监视器（仿真）：2FF 同步 + 半位对齐采样（位中心）、假起始重同步；自由运转状态机（无 @ 事件同步）。}
\begin{Verbatim}[numbers=left,numbersep=5pt,frame=single,fontsize=\footnotesize,breaklines=true,breakanywhere=true]
// =============================================================================
// uart_rx_mon.v —— UART 接收监视器（仅供仿真）：8N1，2 分辨率过采样
// =============================================================================
`timescale 1ns/1ps
module uart_rx_mon #(
    parameter integer CLK_DIV = 646
)(
    input  wire       I_clk,
    input  wire       I_rst_n,
    input  wire       I_rxd,
    output reg  [7:0] O_byte,
    output reg        O_valid
);
    reg [1:0] sync;
    always @(posedge I_clk or negedge I_rst_n) begin
        if (!I_rst_n) sync <= 2'b11;
        else          sync <= {sync[0], I_rxd};
    end
    wire rxd = sync[1];

    reg [3:0] bit_cnt;
    reg [7:0] shifter;
    reg [9:0] div;
    reg       busy;
    reg       done_d;

    always @(posedge I_clk or negedge I_rst_n) begin
        if (!I_rst_n) begin
            busy <= 0; bit_cnt <= 0; div <= 0; shifter <= 0;
            O_byte <= 0; O_valid <= 0; done_d <= 0;
        end else begin
            O_valid <= 0;
            done_d  <= O_valid;
            if (!busy) begin
                if (!rxd) begin            // 起始位
                    busy <= 1;
                    div  <= CLK_DIV/2;     // 半拍后采样对齐位中心
                    bit_cnt <= 0;
                end
            end else begin
                div <= div + 1;
                if (div == CLK_DIV - 1) begin
                    div <= 0;
                    if (bit_cnt == 0) begin
                        if (rxd) begin      // 假起始
                            busy <= 0;
                        end else begin
                            bit_cnt <= 1;
                        end
                    end else if (bit_cnt <= 8) begin
                        shifter <= {rxd, shifter[7:1]};
                        bit_cnt <= bit_cnt + 1;
                    end else begin
                        busy <= 0;
                        if (rxd) begin      // 停止位正确
                            O_byte  <= shifter;
                            O_valid <= 1;
                        end
                    end
                end
            end
        end
    end
endmodule
\end{Verbatim}
\subsection{tb\_system\_integration.v}
\fileintro{整链对接仿真（TD 阶段项）：虚拟 RAW10 传感器 -> 官方 ISP -> 重排 -> vision\_top -> OSD 转储 + 四判据（几何/检测/云台方向/UART）。最后一次全链运行受厂商 ISP 仿真环境限制，列 TD 阶段验证。}
\begin{Verbatim}[numbers=left,numbersep=5pt,frame=single,fontsize=\footnotesize,breaklines=true,breakanywhere=true]
// =============================================================================
// tb_system_integration.v —— 嵌入式x视觉 整链对接仿真（核心交付）
// 虚拟 RAW10 传感器（Python 生成 BGGR 场景）-> isp_top（官方 ISP：去马赛克/AWB）
// -> 128->96 重排 -> vision_top（检测/OSD/云台）-> 帧转储 + 断言
//
// 判据（功能级，非像素级——AWB 会自适应改变颜色）：
//   A. ISP 输出几何：每帧 tuser 恰 1 次、tlast 恰 720 次
//   B. 检测锁定：valid=1 且 |cx−539|<=40、|cy−339|<=40（红块中心 539/339）
//   C. 云台方向：duty_x < 1500（块在左半屏）
//   D. UART 遥测：收到 A5..5A 帧且 valid=1、坐标与检测一致
// =============================================================================
`timescale 1ns/1ps
module tb_system_integration;
    reg clk = 0, pwm_clk = 0, rst_n = 0, pwm_rst_n = 0;
    always #7 clk = ~clk;          // ~71.4MHz
    always #10 pwm_clk = ~pwm_clk; // 50MHz

    // ---------------- RAW10 激励 ----------------
    reg [39:0] beats [0:921599];    // 4 帧余量
    reg [1:0]  metas [0:921599];    // bit1=sof bit0=eol
    integer nb;
    initial begin
        $readmemh("raw10_beats.hex", beats);
        $readmemh("raw10_meta.txt", metas);
        nb = 230400 * 4;            // 4 帧
    end

    wire        isp_tready;
    reg  [39:0] raw_data;
    reg         raw_valid, raw_sof, raw_eol;
    integer     bi = 0;

    // 馈送：每拍一 beat（ isp 内部为整拍流水），tsof/teol 随拍
    always @(posedge clk or negedge rst_n) begin
        if (!rst_n) begin
            raw_valid <= 1'b0; raw_sof <= 1'b0; raw_eol <= 1'b0; raw_data <= 0;
        end else if (bi < nb) begin
            raw_valid <= 1'b1;
            raw_data  <= beats[bi];
            raw_sof   <= metas[bi][1];
            raw_eol   <= metas[bi][0];
            bi = bi + 1;
        end else begin
            raw_valid <= 1'b0; raw_sof <= 1'b0; raw_eol <= 1'b0;
        end
    end

    // ---------------- 被测系统集成顶层 ----------------
    wire pwm_x, pwm_y, uart_txd, osd_valid, osd_sof, osd_eol;
    wire [23:0] osd_data;
    wire [12:0] duty_x, duty_y;
    wire [15:0] lost_cnt;

    system_top u_sys (
        .I_clk(clk), .I_rst_n(rst_n),
        .I_pwm_clk(pwm_clk), .I_pwm_rst_n(pwm_rst_n),
        .I_raw10_data(raw_data), .I_raw10_valid(raw_valid),
        .I_raw10_sof(raw_sof), .I_raw10_eol(raw_eol),
        .I_btn_n(2'b11),
        .O_pwm_x(pwm_x), .O_pwm_y(pwm_y),
        .O_duty_x_us(duty_x), .O_duty_y_us(duty_y),
        .O_lost_cnt(lost_cnt),
        .O_uart_txd(uart_txd),
        .O_osd_data(osd_data), .O_osd_valid(osd_valid),
        .O_osd_sof(osd_sof), .O_osd_eol(osd_eol)
    );

    // ---------------- 帧转储（OSD 流） ----------------
    frame_dump24 #(.BASE_NAME("sys_osd"), .MAX_FRAMES(4)) u_dump (
        .I_clk(clk), .I_rst_n(rst_n),
        .I_valid(osd_valid), .I_eol(osd_eol), .I_data(osd_data),
        .O_frame_done());

    // ---------------- ISP 几何检查 ----------------
    integer tuser_cnt = 0, tlast_cnt = 0, frames_done = 0, geo_err = 0;
    always @(posedge clk) begin
        if (rst_n && u_sys.isp_tvalid) begin
            if (u_sys.isp_tuser) begin
                if (frames_done > 0 && (tuser_cnt != 1 || tlast_cnt != 720)) begin
                    $display("GEO_ERR frame=%0d tuser=%0d tlast=%0d", frames_done, tuser_cnt, tlast_cnt);
                    geo_err = geo_err + 1;
                end
                tuser_cnt = 1; tlast_cnt = 0; frames_done = frames_done + 1;
            end
            if (u_sys.isp_tlast) tlast_cnt = tlast_cnt + 1;
        end
    end

    // ---------------- 检测/云台检查 ----------------
    integer det_ok = 0, det_seen = 0;
    wire [10:0] cx  = u_sys.u_vision.t_cx;
    wire [10:0] cy  = u_sys.u_vision.t_cy;
    wire        tv  = u_sys.u_vision.t_valid;
    always @(posedge clk) begin
        if (rst_n && u_sys.u_vision.t_upd && tv && frames_done >= 3) begin
            det_seen = det_seen + 1;
            if ((cx >= 499) && (cx <= 579) && (cy >= 299) && (cy <= 379))
                det_ok = det_ok + 1;
            else
                $display("DET_MISS frame=%0d cx=%0d cy=%0d", frames_done, cx, cy);
        end
    end

    // ---------------- UART 解码监视 ----------------
    wire [7:0] urx_byte; wire urx_valid;
    uart_rx_mon #(.CLK_DIV(646)) u_mon (
        .I_clk(clk), .I_rst_n(rst_n), .I_rxd(uart_txd),
        .O_byte(urx_byte), .O_valid(urx_valid));
    integer u_state = 0, u_good = 0, u_frames = 0;
    reg [7:0] u_b1, u_b2, u_b3, u_b4, u_b5, u_b6, u_b7;
    always @(posedge clk) begin
        if (rst_n && urx_valid) begin
            case (u_state)
                0: if (urx_byte == 8'hA5) u_state = 1;
                1: begin u_b1 = urx_byte; u_state = 2; end
                2: begin u_b2 = urx_byte; u_state = 3; end
                3: begin u_b3 = urx_byte; u_state = 4; end
                4: begin u_b4 = urx_byte; u_state = 5; end
                5: begin u_b5 = urx_byte; u_state = 6; end
                6: begin u_b6 = urx_byte; u_state = 7; end
                7: begin u_b7 = urx_byte; u_state = 8; end
                8: if (urx_byte == 8'h5A) begin
                    u_frames = u_frames + 1;
                    // 校验和 + valid=1 + cx 在框内
                    if ((u_b7 == (8'hA5 + u_b1 + u_b2 + u_b3 + u_b4 + u_b5 + u_b6))
                        && (u_b1 == 8'h01) && (u_b2 >= 8'd24) && (u_b2 <= 8'd67))
                        u_good = u_good + 1;
                    u_state = 0;
                end else u_state = 0;
            endcase
        end
    end

    // ---------------- 终判 ----------------
    integer guard;
    initial begin
        repeat (5) @(posedge clk);
        rst_n = 1;
        repeat (4) @(posedge pwm_clk);
        pwm_rst_n = 1;
        guard = 1;
    end
    always @(posedge clk) begin
        if (rst_n && (frames_done >= 8) && guard && (det_ok >= 2) && (u_good >= 2)) begin
            $display("TB_RESULT: PASS frames=%0d det_ok=%0d uart_good=%0d duty_x=%0d",
                     frames_done, det_ok, u_good, duty_x);
            $finish;
        end
        if (rst_n && (frames_done >= 8) && (det_seen >= 4) && (det_ok == 0) && guard) begin
            $display("TB_RESULT: FAIL detect locked on wrong target/none (det_seen=%0d)", det_seen);
            $finish;
        end
        if (rst_n && (u_dump.frame_idx >= 3) && guard && (u_frames < 2)) begin
            $display("TB_RESULT: FAIL uart no good frames=%0d", u_frames);
            $finish;
        end
    end
    initial begin
        #450000000;
        $display("TB_RESULT: FAIL timeout frames=%0d det_ok=%0d uart_good=%0d geo_err=%0d",
                 frames_done, det_ok, u_good, geo_err);
        $finish;
    end
endmodule
\end{Verbatim}
\subsection{tb\_i2c\_config.v}
\fileintro{摄像头配置链对录：官方 uicfgcs520+uii2c 对 I2C BFM，逐条比对 116 条寄存器写事务。已知遗留：BFM ACK 位相（TD 阶段项）。}
\begin{Verbatim}[numbers=left,numbersep=5pt,frame=single,fontsize=\footnotesize,breaklines=true,breakanywhere=true]
// =============================================================================
// tb_i2c_config.v —— SC500CS 摄像头配置链验证
// uicfgcs500_720p + uii2c + 寄存器 ROM（官方全套）对着 I2C BFM 跑，
// 逐条比对 BFM 收到的写事务与 uics500reg_720p60 ROM：
//   通过判据：cfg_done 拉高 且 前 116 条写事务 (addr,reg_hi,reg_lo,data)
//   与 ROM 完全一致（AE 追加写只计数不校验）。
// =============================================================================
`timescale 1ns/1ps
module tb_i2c_config;
    reg clk_24 = 0;                       // 24MHz MCLK
    always #20.833 clk_24 = ~clk_24;
    reg rst_n = 0;

    wire cam_scl;
    wire cam_sda;
    reg  [15:0] ae = 16'd0, ag = 16'd0;
    wire ae_req = 1'b0;
    wire cfg_done, ae_cfg_done;

    // 上拉（1kΩ 行为模型）：开漏总线的 1 电平由上拉提供
    pullup(cam_sda);

    // 被测：官方配置链
    uicfgcs500_720p #(.CLK_DIV(239)) u_cfg (
        .I_clk(clk_24), .I_rst_n(rst_n),
        .I_ae_req(ae_req), .I_ae(ae), .I_ag(ag),
        .O_cam_scl(cam_scl), .IO_cam_sda(cam_sda),
        .O_cfg_done(cfg_done), .O_ae_cfg_done(ae_cfg_done));

    // BFM（记录数组在 BFM 内部，层次化读取 u_bfm.O_bytes）
    wire [15:0] byte_cnt;
    wire addr_seen, addr_naked;
    i2c_bfm_sc500 #(.MAX_BYTES(2048)) u_bfm (
        .I_clk(clk_24), .I_rst_n(rst_n), .I_scl(cam_scl), .I_sda(cam_sda),
        .O_byte_cnt(byte_cnt),
        .O_addr_seen(addr_seen), .O_addr_naked(addr_naked));

    // 期望 ROM（与被测同源）
    wire [31:0] reg_data;
    wire [8:0]  reg_size;
    reg [8:0]   idx;
    uics500reg_720p60 u_rom (.REG_INDEX(idx), .REG_DATA(reg_data), .REG_SIZE(reg_size));

    integer mism = 0, checked = 0;
    integer i;
    reg [7:0] e_hi, e_lo, e_da;

    // 周期性进度显示
    reg [15:0] last_bc = 0;
    always @(posedge clk_24) begin
        if (rst_n && (byte_cnt % 40 == 0) && byte_cnt != last_bc) begin
            last_bc <= byte_cnt;
            $display("PROG t=%0t tb_wire=%0d bfm_int=%0d match=%b seen=%b", $time, byte_cnt, u_bfm.O_byte_cnt, u_bfm.addr_match, addr_seen);
        end
    end

    // 等待配置完成
    initial begin
        repeat (5) @(posedge clk_24);
        rst_n = 1;
        wait (cfg_done === 1'b1);
        repeat (100) @(posedge clk_24);
        // 终态一次性扫描：前 116 条写事务与 ROM 逐条比对
        checked = 0;
        for (i = 0; i < 116; i = i + 1) begin
            idx = i[8:0];
            e_hi = reg_data[23:16]; e_lo = reg_data[15:8]; e_da = reg_data[7:0];
            if (u_bfm.O_bytes[i*4+0] !== 8'h6C ||
                u_bfm.O_bytes[i*4+1] !== e_hi ||
                u_bfm.O_bytes[i*4+2] !== e_lo ||
                u_bfm.O_bytes[i*4+3] !== e_da) begin
                if (mism < 4)
                    $display("MISM wr#%0d got %h %h %h %h | exp 6C %h %h %h",
                             i, u_bfm.O_bytes[i*4+0], u_bfm.O_bytes[i*4+1],
                             u_bfm.O_bytes[i*4+2], u_bfm.O_bytes[i*4+3], e_hi, e_lo, e_da);
                mism = mism + 1;
            end
            checked = checked + 1;
        end
        if ((mism == 0) && (checked == 116) && (addr_seen === 1'b1))
            $display("TB_RESULT: PASS (writes=%0d)", checked);
        else
            $display("TB_RESULT: FAIL mism=%0d checked=%0d addr_seen=%b naked=%b",
                     mism, checked, addr_seen, addr_naked);
        $finish;
    end

    initial begin
        #60000000;   // 60ms 上限（116 条 @ ~50kHz I2C 约 1.2ms，AE 追加后仍远小于此）
        $display("TB_RESULT: FAIL timeout checked=%0d mism=%0d cfg_done=%b", checked, mism, cfg_done);
        $finish;
    end
endmodule
\end{Verbatim}
\section{板级点亮仿真设施（08，SC520CS 通路）}
\subsection{tb\_video\_passthrough.v}
\fileintro{板级整链仿真 TB：虚拟 RAW10（1/8 占空馈送，AXI 正确握手每拍一消费）-> 官方 ISP -> 弹性串化 -> 逐帧转储。实测 PASS：frames\_dumped=2（每帧 921600 像素）、tuser=3、ser\_ovf=0。踩坑记：向量 int32 溢出致 readmemh 终止、替身参数名、AXI 重复消费。}
\begin{Verbatim}[numbers=left,numbersep=5pt,frame=single,fontsize=\footnotesize,breaklines=true,breakanywhere=true]
`timescale 1ns/1ps
// =============================================================================
// tb_video_passthrough.v —— 视频传输路线整链仿真（摄像头画面到达验证）
// -----------------------------------------------------------------------------
// 链路：虚拟 RAW10 传感器(文件回放) -> isp_top(官方明码: demosaic/AWB/重排)
//       -> 96bit 抽头(AWB 输出) -> ser96to24(1ppc) -> frame_dump24(逐帧转储)
//
// 判据与输出：
//   1. 转储 passthrough_frame_N.hex（每帧 1280x720 RGB888, 行主序）
//      —— 这就是"屏幕应该显示的画面"（仿真里的 HDMI 替身）
//   2. 几何报告：AWB 输出的 tuser/tlast 计数（诚实打印；
//      已知发现：demosaic 的 tuser 链在纯 RTL 仿真可能不激活，
//      上板若"有数据无帧同步"先查这里 —— 见使用指导 PDF）
//   3. PASS = 至少 2 帧完整转储
// =============================================================================
module tb_video_passthrough;
    reg clk = 0, rst_n = 0;
    always #7 clk = ~clk;                 // ~71.4MHz（axi4s_video_aclk 仿真值）

    // ---------------- 虚拟 RAW10 传感器（文件回放） ----------------
    reg [39:0] beats [0:921599];          // 40bit = 4 像素 x RAW10
    reg [1:0]  metas [0:921599];          // bit1=sof(帧首) bit0=eol(行尾)
    integer    nb;
    integer    bi = 0;

    initial begin
        $readmemh("raw10_beats.hex", beats);
        $readmemh("raw10_meta.txt",  metas);
        nb = 230400 * 3;                  // 3 帧馈送: 前 2 帧为验收目标, 第 3 帧兼作流水冲刷
    end

    wire        isp_ready;                // isp 对上游的就绪指示
    reg         raw_valid = 1'b0;
    reg  [39:0] raw_data  = 40'd0;
    reg         raw_sof = 1'b0, raw_eol = 1'b0;

    // 馈送占空 1/8：真实 MIPI 传感器按节拍出图，且下游串化器每 4 拍才
    // 消费 1 个 beat——满速馈送会淹没 1ppc 消费端（溢出丢像素）。
    // 1/8 占空馈送 + 正确 AXI 握手（每拍恰好被消费一次）：
    // - 满 1/8：demosaic 输出 4px/拍 vs 转储侧 1px/拍存在速率差，
    //   满速会让弹性 FIFO 满而搁浅尾帧；1/8 占空留足排空余量。
    // - 历史：valid 跨拍保持会使同一 beat 被重复消费 8 次，eol/tuser
    //   计数 x8 打乱 Bayer 行相位——务必一拍一消费。
    reg [2:0] pres;
    always @(posedge clk or negedge rst_n) begin
        if (!rst_n) pres <= 3'd0;
        else        pres <= pres + 3'd1;
    end
    wire feed_slot = (pres == 3'd0);
    always @(posedge clk or negedge rst_n) begin
        if (!rst_n) begin
            raw_valid <= 1'b0;
        end else if (raw_valid) begin
            if (isp_ready) begin
                raw_valid <= 1'b0;                 // 本拍已被消费
            end
        end else if (feed_slot && (bi < nb)) begin
            raw_data  <= beats[bi];
            raw_sof   <= metas[bi][1];
            raw_eol   <= metas[bi][0];
            raw_valid <= 1'b1;
            bi = bi + 1;
        end
    end

    // ---------------- 官方 ISP 明码链 ----------------
    wire [127:0] isp_tdata;
    wire         isp_tvalid, isp_tuser, isp_tlast;

    isp_top u_isp (
        .axi4s_video_aclk (clk),
        .I_rst_n          (rst_n),
        .I_tlast          (raw_eol),
        .I_tuser          (raw_sof),
        .I_tdata          (raw_data),
        .I_tvalid         (raw_valid),
        .I_tdest          (10'd0),
        .O_tready         (isp_ready),     // 下游就绪 -> ISP 上游流控
        .O_tdata          (isp_tdata),
        .O_tlast          (isp_tlast),
        .O_tuser          (isp_tuser),
        .O_tvalid         (isp_tvalid),
        .I_tready         ()               // 输出端口, 仿真不关心
    );

    // ---------------- 96bit 抽头（AWB 输出）-> 弹性串化 -> 逐帧转储 ----------------
    // 说明：官方链路 ISP 输出送 DDR 帧缓存（深缓冲吸收行级突发）；
    // 仿真没有 DDR，用 px_serializer_fifo（8192 beat 弹性 FIFO + 真反压）
    // 充当"浅 DDR"，保证零丢拍。ser96to24 是无反压窃流设计（13 深 FIFO），
    // 吸不住 demosaic 的行级突发，不能用在整链仿真里。
    wire        ser_valid, ser_sof, ser_eol;
    wire [23:0] ser_data;
    wire [15:0] ser_ovf;

    px_serializer_fifo u_ser (
        .I_clk   (clk), .I_rst_n(rst_n),
        .I_valid (u_isp.awb_O_tvalid),         // 128->96 重排后的 4 像素/拍
        .I_sof   (u_isp.awb_O_tuser),
        .I_eol   (u_isp.awb_O_tlast),
        .I_data  (u_isp.awb_O_tdata),
        .O_ready (isp_ready),                  // 真反压回 ISP 上游
        .O_valid (ser_valid), .O_sof(ser_sof),
        .O_eol   (ser_eol),   .O_data(ser_data),
        .O_dbg_overflows(ser_ovf)
    );

    wire dump_done;
    frame_dump24 #(
        .BASE_NAME ("passthrough_frame"),
        .MAX_FRAMES(3)
    ) u_dump (
        .I_clk(clk), .I_rst_n(rst_n),
        .I_valid(ser_valid), .I_eol(ser_eol), .I_data(ser_data),
        .O_frame_done(dump_done)
    );

    // ---------------- 几何观察（诚实打印，不作 PASS 依据） ----------------
    integer tuser_total = 0, tlast_total = 0;
    always @(posedge clk) begin
        if (rst_n && u_isp.awb_O_tvalid) begin
            if (u_isp.awb_O_tuser) tuser_total = tuser_total + 1;
            if (u_isp.awb_O_tlast) tlast_total = tlast_total + 1;
        end
    end

    // ---------------- 终判 ----------------
    integer dumped = 0;
    always @(posedge clk) begin
        if (rst_n && dump_done) dumped = dumped + 1;
    end

    initial begin
        repeat (5) @(posedge clk);
        rst_n = 1;
    end

    initial begin
        #500000000;   // 看门狗：1/8 馈送约 77ms + 行窗节流, 500ms 足够 3 帧
        $display("TB_RESULT: %s frames_dumped=%0d tuser_total=%0d tlast_total=%0d ser_ovf=%0d",
                 (dumped >= 2) ? "PASS" : "FAIL", dumped, tuser_total, tlast_total, ser_ovf);
        if (tuser_total == 0)
            $display("NOTE: tuser 链在本纯 RTL 仿真环境不激活（已知发现，见使用指导 PDF 第 5 节）");
        $finish;
    end

    // 提前收工：3 帧全部转储完
    always @(posedge clk) begin
        if (rst_n && dumped >= 3) begin
            $display("TB_RESULT: PASS frames_dumped=%0d tuser_total=%0d tlast_total=%0d ser_ovf=%0d",
                     dumped, tuser_total, tlast_total, ser_ovf);
            if (tuser_total == 0)
                $display("NOTE: tuser 链在本纯 RTL 仿真环境不激活（已知发现，见使用指导 PDF 第 5 节）");
            $finish;
        end
    end

endmodule
\end{Verbatim}
\subsection{tb\_isp\_burst.v}
\fileintro{ISP 实流突发测量 TB（2026-10-03 架构级复审新增）：虚拟 RAW10 满速 / 1:9 / 1:6 三档馈送 -> 官方 ISP -> data\_96bit\_to\_128bit(4:3) -> data128\_96 -> ser96to24（与 system\_top 同接法）。实测：AWB 行级突发 maxrun=320（1:6 达 1925）；128 重排 maxrun=3 且 tuser\&\&valid=0；ser96to24 输出 188 像素/零 SOF/丢弃饱和——坐实"128 像素缓冲吸不住行级突发 + 帧同步不可建立"。}
\begin{Verbatim}[numbers=left,numbersep=5pt,frame=single,fontsize=\footnotesize,breaklines=true,breakanywhere=true]
// =============================================================================
// tb_isp_burst.v —— ISP 满速馈送下的 96bit 突发形态测量 + ser96to24 实链丢拍判定
// -----------------------------------------------------------------------------
// 背景（外部架构复审 Gate）：视觉链的 ser96to24（无反压、128 像素缓冲）此前只在
// tpg_720p 的 "1 beat/4clk" 良性节拍下仿真过；真实 ISP 的 AWB 输出在满速馈送下
// 是否行级突发、128 像素缓冲（=32 beat）是否足够，从未实测过。
//
// 结构（2026-10-03 修复后的 system_top 接法）：
//   虚拟 RAW10（满速 / 1:N 分频 / 行突发+消隐 三种模式）
//     -> isp_top
//     -> AWB 96bit 抽头（4px RGB888，tuser/tlast 均在有效拍）
//     -> ser96to24 v4（2048px 寄存器读弹性缓冲，无反压）
//   同时监测 128bit 重排路径（历史问题路径，仅供对照）。
//
// 判据（TB_RESULT）：
//   ser 输出像素 = 3*921600 且 sof=3 且 eol=3*720 且 drop=0 -> PASS
//   （v3 时代本 TB 同一激励下输出 188/2764800 像素、零 SOF——修复的验收对照）
//
// 输出（ISPBURST_REPORT）：
//   * awb 96bit / v128 valid 的最大连续长度、超 42 拍（=128 像素缓冲临界）段数
//   * ser96to24 输出像素数、SOF/EOL 计数、O_drop_cnt（饱和）
//   判定：s_pix == 3*921600 且 drop==0 -> 该模式下达标；否则输出帧像素缺失。
// =============================================================================
`timescale 1ns/1ps
module tb_isp_burst;
    // FEED_DIV: 0=满速（每拍一个 beat）；N=每 N 拍一个 beat（模拟传感器平均占空）
    parameter FEED_DIV = 0;
    // FEED_MODE: 0=均匀分频；1=行突发（320 beat 连发 + 消隐空闲，模拟真实传感器行节奏）
    parameter FEED_MODE = 0;
    parameter IDLE_CYC  = 1330;    // 行突发模式：突发间消隐拍数
    reg clk = 0, rst_n = 0;
    always #7 clk = ~clk;
    initial begin rst_n = 0; #200 rst_n = 1; end

    // ---- 虚拟 RAW10 传感器（满速馈送） ----
    reg [39:0] beats [0:921599];
    reg [1:0]  metas [0:921599];
    integer    nb;
    integer    bi = 0;
    initial begin
        $readmemh("raw10_beats.hex", beats);
        $readmemh("raw10_meta.txt",  metas);
        nb = 230400 * 3;
    end

    wire        isp_up_ready;              // ISP 对上游的就绪输出（I_tready）
    reg         raw_valid = 1'b0;
    reg  [39:0] raw_data  = 40'd0;
    reg         raw_sof = 1'b0, raw_eol = 1'b0;
    integer     feed_beats = 0, stall_cyc = 0;
    integer     div_cnt = 0;
    integer     line_cnt = 0;              // 行突发模式：行内已发 beat 数
    integer     idle_cnt = 0;              // 行突发模式：消隐计时
    reg         pend = 1'b0;               // 当前 beat 已呈现、等待被消费

    // 分频/行突发馈送：每个 beat 恰好呈现一次（valid 保持到 ready），杜绝重复消费
    wire feed_slot = (FEED_MODE == 1) ? (idle_cnt == 0)
                                      : ((FEED_DIV == 0) || (div_cnt == 0));
    always @(posedge clk or negedge rst_n) begin
        if (!rst_n) begin
            raw_valid <= 1'b0;
            pend      <= 1'b0;
            div_cnt   <= 0;
            line_cnt  <= 0;
            idle_cnt  <= 0;
        end else begin
            div_cnt <= (div_cnt >= FEED_DIV) ? 0 : div_cnt + 1;
            if (FEED_MODE == 1) begin
                // 行突发：先连发 320 beat，再空闲 IDLE_CYC 拍
                if (idle_cnt > 0)
                    idle_cnt <= idle_cnt - 1;
                else if (line_cnt >= 320) begin
                    line_cnt <= 0;
                    idle_cnt <= IDLE_CYC;
                end
            end
            if (pend) begin
                if (isp_up_ready) begin
                    pend      <= 1'b0;
                    raw_valid <= 1'b0;
                    feed_beats = feed_beats + 1;
                    if (FEED_MODE == 1) line_cnt = line_cnt + 1;
                end else begin
                    stall_cyc = stall_cyc + 1;
                end
            end else if ((bi < nb) && feed_slot) begin
                raw_data  <= beats[bi];
                raw_sof   <= metas[bi][1];
                raw_eol   <= metas[bi][0];
                bi = bi + 1;
                pend      <= 1'b1;
                raw_valid <= 1'b1;
            end
        end
    end

    // ---- 官方 ISP（O_tready=1：视觉分支恒不反压，最坏情况；与 system_top 语义一致） ----
    wire [127:0] isp_tdata;
    wire         isp_tvalid, isp_tuser, isp_tlast;
    wire [95:0]  tap_tdata;
    wire         tap_tvalid, tap_tuser, tap_tlast;

    isp_top u_isp (
        .axi4s_video_aclk (clk), .I_rst_n (rst_n),
        .I_tlast (raw_eol), .I_tuser (raw_sof), .I_tdata (raw_data),
        .I_tvalid (raw_valid), .I_tdest (10'd0),
        .O_tready (1'b1),
        .O_tdata (isp_tdata), .O_tlast (isp_tlast), .O_tuser (isp_tuser),
        .O_tvalid (isp_tvalid), .I_tready (isp_up_ready),
        .tap_tdata (tap_tdata), .tap_tvalid (tap_tvalid),
        .tap_tuser (tap_tuser), .tap_tlast (tap_tlast)
    );

    // ---- AWB 96bit 抽头（经 isp_top 新增端口引出，验证与 system_top 相同的端口路径） ----
    wire         a_user96, a_valid96, a_last96;
    wire [95:0]  a_data96;
    assign a_user96  = tap_tuser;
    assign a_valid96 = tap_tvalid;
    assign a_last96  = tap_tlast;
    assign a_data96  = tap_tdata;

    // ---- 128bit 重排（视觉链真实输入节奏） ----
    wire         a_user128, a_valid128;
    wire [127:0] a_data128;
    assign a_user128  = u_isp.awb_O_tuser_128;
    assign a_valid128 = u_isp.awb_O_tvalid_128;
    assign a_data128  = u_isp.awb_O_tdata_128;

    // ---- ser96to24（v4 修复后接法：直接抽 AWB 96bit 输出；无反压） ----
    // 修复依据见迭代记录 #N+6/#N+7：原 128 重排路径 tuser 结构性错位 + 3/4 占空节拍，
    // 且 v3 串化器 128px 缓冲吸不住行级突发。本实例即修复后的验收对象。
    wire        s_valid, s_sof, s_eol;
    wire [23:0] s_data;
    wire [15:0] s_drop;
    ser96to24 u_ser (
        .I_clk(clk), .I_rst_n(rst_n),
        .I_tdata(a_data96), .I_tvalid(a_valid96), .I_tuser(a_user96),
        .I_tlast(a_last96),
        .O_tready(), .O_valid(s_valid), .O_sof(s_sof), .O_eol(s_eol),
        .O_data(s_data), .O_drop_cnt(s_drop));

    // ---- 统计 ----
    // FIFO 最大占用（复审要求的安全余量指标）：直接观测 DUT 内部缓冲计数
    integer occ_max = 0;
    always @(posedge clk)
        if (u_ser.cnt > occ_max) occ_max <= u_ser.cnt;

    integer a_run = 0, a_max = 0, a_gt42 = 0, a_beats = 0, a_frames = 0;
    integer v_run = 0, v_max = 0, v_gt42 = 0, v_beats = 0, v_frames = 0;
    integer s_pix = 0, s_sof_cnt = 0, s_eol_cnt = 0;
    integer cyc = 0, first_cyc = -1, last_cyc = -1;
    // 突发长度直方图（96bit AWB）：1-8 / 9-42 / 43-320 / 321-640 / 641-1280 / >1280
    integer h1 = 0, h2 = 0, h3 = 0, h4 = 0, h5 = 0, h6 = 0;
    // tlast 在突发中的位置（1=突发起始拍，~320=突发末拍）
    integer bpos = 0, tl_pos_min = 999999, tl_pos_max = -1, tl_cnt = 0;

    always @(posedge clk) begin
        cyc <= cyc + 1;
        // 96bit valid 连续段 + 直方图 + tlast 位置
        if (a_valid96) begin
            a_run <= a_run + 1;
            a_beats <= a_beats + 1;
            bpos  <= bpos + 1;
            if (a_run + 1 > a_max) a_max <= a_run + 1;
            if (first_cyc < 0) first_cyc <= cyc;
            last_cyc <= cyc;
            if (a_last96) begin
                tl_cnt <= tl_cnt + 1;
                if (bpos + 1 < tl_pos_min) tl_pos_min <= bpos + 1;
                if (bpos + 1 > tl_pos_max) tl_pos_max <= bpos + 1;
            end
        end else begin
            if (a_run > 42) a_gt42 <= a_gt42 + 1;
            if (a_run >= 1  && a_run <= 8)    h1 <= h1 + 1;
            if (a_run >= 9  && a_run <= 42)   h2 <= h2 + 1;
            if (a_run >= 43 && a_run <= 320)  h3 <= h3 + 1;
            if (a_run >= 321 && a_run <= 640) h4 <= h4 + 1;
            if (a_run >= 641 && a_run <= 1280)h5 <= h5 + 1;
            if (a_run > 1280)                 h6 <= h6 + 1;
            a_run <= 0;
            bpos  <= 0;
        end
        if (a_valid96 && a_user96) a_frames <= a_frames + 1;
        // 128bit valid 连续段（ser96to24 实际看到的节拍）
        if (a_valid128) begin
            v_run <= v_run + 1;
            v_beats <= v_beats + 1;
            if (v_run + 1 > v_max) v_max <= v_run + 1;
        end else begin
            if (v_run > 42) v_gt42 <= v_gt42 + 1;
            v_run <= 0;
        end
        if (a_valid128 && a_user128) v_frames <= v_frames + 1;
        // ser 输出
        if (s_valid) s_pix <= s_pix + 1;
        if (s_sof) s_sof_cnt <= s_sof_cnt + 1;
        if (s_eol) s_eol_cnt <= s_eol_cnt + 1;
    end

    // ---- 尾帧冲刷 + 报告 ----
    initial begin
        wait (bi >= nb);
        @(posedge clk);
        repeat (50000) @(posedge clk);
        $display("ISPBURST_REPORT:");
        $display("  awb96 : beats=%0d frames=%0d maxrun=%0d runs_gt42=%0d",
                 a_beats, a_frames, a_max, a_gt42);
        $display("  v128  : beats=%0d frames=%0d maxrun=%0d runs_gt42=%0d",
                 v_beats, v_frames, v_max, v_gt42);
        $display("  ser   : pixels=%0d sof=%0d eol=%0d drop=%0d (expect_pix=%0d)",
                 s_pix, s_sof_cnt, s_eol_cnt, s_drop, 921600*3);
        $display("  occ   : ser FIFO max occupancy = %0d / 2048 px (%0d%%)",
                 occ_max, (occ_max*100)/2048);
        $display("  duty  : a_beats/(last-first+1) = %0d/%0d", a_beats,
                 (last_cyc - first_cyc + 1));
        $display("  feed  : beats=%0d stall_cyc=%0d (ISP 上游反压拍数)", feed_beats, stall_cyc);
        $display("  hist  : 1-8=%0d 9-42=%0d 43-320=%0d 321-640=%0d 641-1280=%0d >1280=%0d",
                 h1, h2, h3, h4, h5, h6);
        $display("  tlast : cnt=%0d pos_in_burst: min=%0d max=%0d",
                 tl_cnt, tl_pos_min, tl_pos_max);
        if ((s_pix == 921600*3) && (s_sof_cnt == 3) && (s_eol_cnt == 720*3) && (s_drop == 0))
            $display("TB_RESULT: PASS ser pixels=%0d sof=%0d eol=%0d drop=0", s_pix, s_sof_cnt, s_eol_cnt);
        else
            $display("TB_RESULT: FAIL ser pixels=%0d (expect %0d) sof=%0d eol=%0d drop=%0d",
                     s_pix, 921600*3, s_sof_cnt, s_eol_cnt, s_drop);
        $finish;
    end

    // 看门狗
    initial begin
        #200_000_000;
        $display("ISPBURST WATCHDOG TIMEOUT bi=%0d", bi);
        $finish;
    end
endmodule
\end{Verbatim}
\subsection{px\_serializer\_fifo.v}
\fileintro{仿真专用弹性串化器：8192 beat FIFO + 真反压（替代无反压的 ser96to24 做整链转储）；充当"浅 DDR"吸收 ISP 行级突发。}
\begin{Verbatim}[numbers=left,numbersep=5pt,frame=single,fontsize=\footnotesize,breaklines=true,breakanywhere=true]
`timescale 1ns/1ps
// =============================================================================
// px_serializer_fifo.v —— 仿真专用弹性串化器（96bit 4像素/拍 -> 1像素/拍）
// -----------------------------------------------------------------------------
// 与 ser96to24（无反压窃流设计）不同：本模块带 8192 beat 弹性 FIFO，
// 并把"空位不足"如实反压回上游（O_ready -> isp 上游流控）。
// 官方链路里 ISP 输出去往 DDR 帧缓存（深缓冲吸收行级突发）；
// 仿真里没有 DDR，用本模块充当"浅 DDR"，保证零丢拍。
// 每个存储字 = {eol, sof, data[95:0]}；sof 标在首像素、eol 标在第 4 像素。
// 仅供仿真，禁止综合。
// =============================================================================
module px_serializer_fifo #(
    parameter DEPTH = 8192                    // beat 数（2 的幂）
)(
    input  wire        I_clk,
    input  wire        I_rst_n,
    input  wire        I_valid,
    input  wire        I_sof,
    input  wire        I_eol,
    input  wire [95:0] I_data,
    output wire        O_ready,               // 对上游的反压（空位充足才为 1）
    output reg         O_valid,
    output reg         O_sof,
    output reg         O_eol,
    output reg  [23:0] O_data,
    output reg  [15:0] O_dbg_overflows        // 调试：因满被拒的 beat 计数
);
    reg [97:0] mem [0:DEPTH-1];
    reg [12:0] waddr, raddr;
    reg [13:0] count;                         // FIFO 内 beat 数
    reg [1:0]  px_sel;                        // 当前 beat 内像素指针

    assign O_ready = (count < DEPTH - 4);

    wire wr        = I_valid && O_ready;
    wire rd        = (count != 14'd0);
    wire beat_done = rd && (px_sel == 2'd3);  // 本拍第 4 像素已输出 = beat 消费完

    wire [97:0] rword = mem[raddr];
    wire [23:0] rpx   = (px_sel == 2'd0) ? rword[23:0]   :
                        (px_sel == 2'd1) ? rword[47:24]  :
                        (px_sel == 2'd2) ? rword[71:48]  :
                                           rword[95:72];

    always @(posedge I_clk or negedge I_rst_n) begin
        if (!I_rst_n) begin
            waddr <= 13'd0; raddr <= 13'd0; count <= 14'd0; px_sel <= 2'd0;
            O_valid <= 1'b0; O_sof <= 1'b0; O_eol <= 1'b0; O_data <= 24'd0;
            O_dbg_overflows <= 16'd0;
        end else begin
            O_valid <= 1'b0;

            if (wr) begin
                mem[waddr] <= {I_eol, I_sof, I_data};
                waddr <= waddr + 13'd1;
            end
            if (beat_done)
                raddr <= raddr + 13'd1;

            // count: 写 +1，beat 消费完 -1（同拍两者相抵）
            if (wr && !beat_done)        count <= count + 14'd1;
            else if (!wr && beat_done)   count <= count - 14'd1;

            if (rd) begin
                O_valid <= 1'b1;
                O_data  <= rpx;
                O_sof   <= rword[96] && (px_sel == 2'd0);   // 首像素携带 sof
                O_eol   <= rword[97] && (px_sel == 2'd3);   // 末像素携带 eol
                px_sel <= (px_sel == 2'd3) ? 2'd0 : px_sel + 2'd1;
            end else begin
                O_sof <= 1'b0;
                O_eol <= 1'b0;
            end

            if (I_valid && !O_ready)
                O_dbg_overflows <= O_dbg_overflows + 16'd1;
        end
    end

endmodule
\end{Verbatim}
\subsection{filelist\_sim.f}
\fileintro{板级仿真编译清单（全部仿真专用文件，禁止加入 TD 工程）。}
\begin{Verbatim}[numbers=left,numbersep=5pt,frame=single,fontsize=\footnotesize,breaklines=true,breakanywhere=true]
# filelist_sim.f —— 视频通路整链仿真编译清单（ModelSim 用）
# 约定：D:/fpga-ptz-tracker 为仓库根；本清单供 msim_passthrough.do 使用
# 注意：全部为仿真专用文件，禁止加入 TD 工程

# ---- 视觉方向（串化器 + 帧转储） ----
D:/fpga-ptz-tracker/03_vision/step1_video_link/03_project/rtl/ser96to24.v
D:/fpga-ptz-tracker/03_vision/step1_video_link/03_project/rtl/frame_dump24.v

# ---- 官方 ISP 明码副本（04_embedded/sim）----
D:/fpga-ptz-tracker/04_embedded/03_project/sim/isp_top.v
D:/fpga-ptz-tracker/04_embedded/03_project/sim/data128_96.v
D:/fpga-ptz-tracker/04_embedded/03_project/sim/data_96bit_to_128bit.v
D:/fpga-ptz-tracker/04_embedded/03_project/sim/demosaic_4x_2_0/demosaic.v
D:/fpga-ptz-tracker/04_embedded/03_project/sim/demosaic_4x_2_0/zhenghe.v
D:/fpga-ptz-tracker/04_embedded/03_project/sim/demosaic_4x_2_0/bilinear_interpolation.v
D:/fpga-ptz-tracker/04_embedded/03_project/sim/demosaic_4x_2_0/line_buffer_demosaic.v
D:/fpga-ptz-tracker/04_embedded/03_project/sim/demosaic_4x_2_0/mipi_to_raw_converter.v
D:/fpga-ptz-tracker/04_embedded/03_project/sim/demosaic_4x_2_0/raw_matrix_3x3_buffer.v
D:/fpga-ptz-tracker/04_embedded/03_project/sim/demosaic_4x_2_0/blk_mem_gen_demosaic/blk_mem_gen_demosaic.v
D:/fpga-ptz-tracker/04_embedded/03_project/sim/demosaic_4x_2_0/blk_mem_gen_demosaic/RTL/ram_c36c2ca3b3a7.v
D:/fpga-ptz-tracker/04_embedded/03_project/sim/demosaic_4x_2_0/blk_mem_gen_zhenghe/blk_mem_gen_zhenghe.v
D:/fpga-ptz-tracker/04_embedded/03_project/sim/demosaic_4x_2_0/blk_mem_gen_zhenghe/RTL/ram_d32970204f32.v
D:/fpga-ptz-tracker/04_embedded/03_project/sim/awb.v
D:/fpga-ptz-tracker/04_embedded/03_project/sim/awb/signal_delay.v
D:/fpga-ptz-tracker/04_embedded/03_project/sim/awb/ram/blk_mem_gen_awb_delay_signal.v

# ---- 行为级替身（TD 工程禁止加入）----
D:/fpga-ptz-tracker/04_embedded/03_project/rtl/divider_gate_sim.v
D:/fpga-ptz-tracker/04_embedded/03_project/rtl/ph1p_logic_eram_stub.v

# ---- 本 TB ----
D:/fpga-ptz-tracker/08_board_bringup/01_video_passthrough/sim/px_serializer_fifo.v
D:/fpga-ptz-tracker/08_board_bringup/01_video_passthrough/sim/tb_video_passthrough.v
\end{Verbatim}
\subsection{run\_sim.bat}
\fileintro{板级仿真一键脚本（重建 work 库防 vopt 缓存旧设计）。}
\begin{Verbatim}[numbers=left,numbersep=5pt,frame=single,fontsize=\footnotesize,breaklines=true,breakanywhere=true]
@echo off
rem run_sim.bat —— 视频通路整链仿真一键运行（ModelSim SE 命令行）
rem 前提：ModelSim 安装于 D:\ModelSim（如不同请改下面的 VSIM_DIR）
setlocal
set VSIM_DIR=D:\ModelSim\win64
set PATH=%VSIM_DIR%;%PATH%

cd /d %~dp0
if exist work rmdir /s /q work
vlib work
vlog -f filelist_sim.f
if errorlevel 1 (
    echo ==== 编译失败，检查 filelist_sim.f ====
    exit /b 1
)
vsim -c tb_video_passthrough -do "run -all; quit -f"
echo ==== 仿真结束，转储文件 passthrough_frame_N.hex 见本目录 ====
endlocal
\end{Verbatim}
\subsection{msim\_passthrough.do}
\fileintro{ModelSim do 脚本等价版本。}
\begin{Verbatim}[numbers=left,numbersep=5pt,frame=single,fontsize=\footnotesize,breaklines=true,breakanywhere=true]
# msim_passthrough.do —— 视频通路整链仿真一键脚本
# 用法：在 sim 目录打开 ModelSim 控制台，执行 do msim_passthrough.do
# 或直接运行 run_sim.bat（免 GUI，命令行出结果）

# 重建 work 库(避免 vopt 缓存导致跑旧设计)
if {[file exists work]} { vdel -all }
vlib work
vmap work work

# 编译（-f 清单见 filelist_sim.f）
vlog -f filelist_sim.f

# 仿真（命令行模式，跑完自动退出）
vsim -c tb_video_passthrough -do "run -all; quit -f"
\end{Verbatim}
\section{Python 金标准与工具链}
\subsection{golden\_model.py}
\fileintro{算术金标准（唯一真值来源）：HSV/阈值/平滑/质心/云台/帧差/整链，所有 RTL 按其逐位复刻；含 MIN\_PIXELS/勘探参数表。}
\begin{Verbatim}[numbers=left,numbersep=5pt,frame=single,fontsize=\footnotesize,breaklines=true,breakanywhere=true]
# -*- coding: utf-8 -*-
"""
golden_model.py —— 视觉方向全链路位精确（bit-true）金标准模型
=================================================================
本文件是整个视觉软件包的"算术真值来源"：
  RTL（ModelSim 仿真）与上位机（Python）必须逐位一致。
任何对算术的修改必须先改这里，再同步 RTL，并在迭代报告中记录。

定点算术约定（与 RTL 严格一一对应，见 02_技术文档）：
  1) RGB->HSV：
     max = max(R,G,B), min = min(R,G,B), delta = max-min
     V = max            (8bit)
     S = delta          (8bit, 未归一化饱和度代理，省一个除法器)
     inv[d] = floor(3840/d)                (12bit 查找表, d=1..255)
     3扇区优先级 R > G > B（同 max 时取先者）：
       R 最大: h = clamp(((G-B)*inv[d]) >>> 6, 0, 60), H = 0  + h
       G 最大: h = clamp(((B-R)*inv[d]) >>> 6, 0, 60), H = 120+ h
       B 最大: h = clamp(((R-G)*inv[d]) >>> 6, 0, 60), H = 240+ h
     delta==0 时 H=0（消色差）。
     >>> 为算术右移（floor），Verilog 中对 signed 用 >>>。
     已知限制：H 只能落在 {0..60,120..180,240..300} 三个扇区，
     品红/青/黄绿等中间色调会被并入相邻扇区（详见方案对比报告）。
  2) 阈值分割（掩码 mask=1 表示目标像素）：
     mask = H在[h_lo..h_hi]（支持环绕） 且 S>=s_lo 且 v_lo<=V<=v_hi
  3) 3x3 二值平滑（多数表决）：
     sum(3x3窗内mask) >= 5 -> 1
     边界策略：行/列复制（replicate）；
     y==0: 上两行均取当前行; y==1: line1(顶)取 row0（RTL 中
     顶行缓存在 y==1 时视为无效，用 row0 顶替）——与 RTL 逐拍对应。
  4) 质心/包围框（对平滑后的 mask）：
     count=summask, sum_x=sumx·mask, sum_y=sumy·mask
     count >= MIN_PIXELS 才有效: cx = floor(sum_x/count), cy = floor(sum_y/count)
     bbox = 平滑后 mask 的 min/max x/y（无目标时全 0 且 valid=0）
  5) 云台控制（gimbal_ctrl，Q8.8 定点乘）：
     err = invert ? (center-target) : (target-center)
     kp_term = (err * KP) >>> 8（有符号算术右移）
     |err|<=DZ 时 kp_term=0；限幅 [SERVO_MIN_US, SERVO_MAX_US]；
     单帧限速 SLEW_US；丢失目标保持原脉宽并累加 lost_cnt。
"""

import numpy as np

# ---------------- 全局几何参数（与 RTL 参数默认值一致） ----------------
IMG_W = 1280
IMG_H = 720
CENTER_X = IMG_W // 2      # 640
CENTER_Y = IMG_H // 2      # 360
MIN_PIXELS = 30            # 质心有效性像素门限

# HSV 预置阈值组（8 组，与 hsv_threshold.v 中 PRESETS 完全一致）
# 每组: (h_lo, h_hi, s_lo, v_lo, v_hi)
HSV_PRESETS = [
    (335, 25,  80, 60, 255),   # 0: 红色（环绕窗口）
    (110, 145, 60, 40, 255),   # 1: 绿色
    (200, 250, 60, 40, 255),   # 2: 蓝色
    (20,  70,  80, 80, 255),   # 3: 橙黄
    (170, 195, 60, 60, 255),   # 4: 青
    (300, 330, 60, 40, 255),   # 5: 品红
    (0,   359, 10, 220, 255),  # 6: 白色（低饱和高亮度）
    (0,   359, 0,  0,   60),   # 7: 黑色（低亮度）
]


def build_inv_lut():
    """inv[d] = floor(3840/d)，与 rtl/hsv_inv_lut.v（自动生成）逐项一致。"""
    lut = np.zeros(256, dtype=np.int32)
    for d in range(1, 256):
        lut[d] = 3840 // d
    return lut


_INV = build_inv_lut()


def rgb_to_hsv(r, g, b):
    """向量化 RGB888 -> (H[0..359], S=delta, V=max)，逐位复刻 RTL 算术。"""
    r = np.asarray(r, dtype=np.int32)
    g = np.asarray(g, dtype=np.int32)
    b = np.asarray(b, dtype=np.int32)
    mx = np.maximum(np.maximum(r, g), b)
    mn = np.minimum(np.minimum(r, g), b)
    d = mx - mn
    v = mx
    s = d
    inv = np.where(d > 0, _INV[np.clip(d, 0, 255)], 0)

    # 3 扇区，优先级 R > G > B
    is_r = (r >= g) & (r >= b)
    is_g = (~is_r) & (g >= r) & (g >= b)
    is_b = ~(is_r | is_g)

    h_r = np.clip(((g - b) * inv) >> 6, 0, 60)      # >> 为 floor 算术移位
    h_g = np.clip(((b - r) * inv) >> 6, 0, 60)
    h_b = np.clip(((r - g) * inv) >> 6, 0, 60)

    h = np.where(is_r, h_r, np.where(is_g, 120 + h_g, 240 + h_b))
    h = np.where(d == 0, 0, h)
    return h.astype(np.int32), s.astype(np.int32), v.astype(np.int32)


def hsv_threshold(h, s, v, preset_idx):
    """HSV -> 二值掩码（支持 H 窗口环绕）。"""
    h_lo, h_hi, s_lo, v_lo, v_hi = HSV_PRESETS[preset_idx]
    if h_lo <= h_hi:
        in_h = (h >= h_lo) & (h <= h_hi)
    else:
        in_h = (h >= h_lo) | (h <= h_hi)
    return (in_h & (s >= s_lo) & (v >= v_lo) & (v <= v_hi)).astype(np.uint8)


def smooth3x3_majority(mask, w=IMG_W, hgt=IMG_H):
    """3x3 多数表决平滑，边界复制 + 前两行特殊处理，与 RTL 逐拍语义一致。"""
    m = mask.astype(np.int32)
    out = np.zeros_like(m)
    for y in range(hgt):
        row_top = m[y - 2] if y >= 2 else None
        row_mid = m[y - 1] if y >= 1 else None
        cur = m[y]
        if y == 0:
            top = mid = cur
        elif y == 1:
            top = mid = row_mid          # RTL: line1 无效，用 row0 顶替
        else:
            top, mid = row_top, row_mid
        ext = np.empty(w + 2, dtype=np.int32)
        for row, name in ((top, 't'), (mid, 'm'), (cur, 'b')):
            e = np.empty(w + 2, dtype=np.int32)
            e[1:w + 1] = row
            e[0] = row[0]
            e[w + 1] = row[w - 1]
            if name == 't':
                te = e
            elif name == 'm':
                me = e
            else:
                be = e
        s = (te[:-2] + te[1:-1] + te[2:] +
             me[:-2] + me[1:-1] + me[2:] +
             be[:-2] + be[1:-1] + be[2:])
        out[y] = (s >= 5).astype(np.uint8)
    return out


def centroid_bbox(mask_smooth, w=IMG_W, hgt=IMG_H, min_pixels=MIN_PIXELS):
    """平滑掩码 -> 质心 + 包围框 + 计数（floor 除法，除零保护）。"""
    ys, xs = np.nonzero(mask_smooth)
    cnt = int(xs.size)
    res = dict(count=cnt, valid=0, cx=0, cy=0,
               min_x=0, max_x=0, min_y=0, max_y=0)
    if cnt >= min_pixels:
        sx = int(xs.sum())
        sy = int(ys.sum())
        res.update(valid=1, cx=sx // cnt, cy=sy // cnt,
                   min_x=int(xs.min()), max_x=int(xs.max()),
                   min_y=int(ys.min()), max_y=int(ys.max()))
    return res


def gimbal_ctrl_step(err, duty_prev, kp_q8, kd_q8=0, err_prev=0,
                     deadzone=25, slew=200, servo_min=700, servo_max=2300,
                     invert=False):
    """gimbal_ctrl.v 单轴单帧步进的位精确复刻。err 为有符号像素误差。"""
    if invert:
        err = -err
    if abs(err) <= deadzone:
        kp_term = 0
    else:
        kp_term = (err * kp_q8) >> 8          # Python >> 对负数即 floor
    kd_term = ((err - err_prev) * kd_q8) >> 8 if kd_q8 else 0
    step = kp_term + kd_term
    if step > slew:
        step = slew
    elif step < -slew:
        step = -slew
    duty = duty_prev + step
    if duty > servo_max:
        duty = servo_max
    elif duty < servo_min:
        duty = servo_min
    return duty


# ---------------- 帧差运动检测（motion_detect.v 的位精确复刻，v2） ----------------

def frame_diff_motion(frames, mot_th=12):
    """帧差运动检测：对应 motion_detect.v（2026-09-20 修正版）。

    位精确语义（与 RTL 严格一一对应）：
      - 第 0 帧只填充前帧缓存，无输出脉冲（prev_ready 门控）；
      - 每帧 sof 拍（(0,0) 像素）只写缓存、不参与比对；
      - 第 N>=1 帧：diff = cur ^ prev（除 (0,0) 全帧比对），
        输出 cnt 与 diff 像素包围框，valid = (cnt >= mot_th)；
        cnt==0 时 bbox 全 0。
    frames: [HxW 0/1 uint8] 时间顺序。
    返回 list：第 0 帧为 None，其余为 dict(cnt, valid, min_x, max_x, min_y, max_y)。
    """
    out = []
    prev = None
    for m in frames:
        if prev is None:
            out.append(None)
            prev = m.copy()
            continue
        d = (m.astype(np.int8) != prev.astype(np.int8))
        d[0, 0] = False                       # sof 拍存而不比
        ys, xs = np.nonzero(d)
        cnt = int(xs.size)
        if cnt:
            r = dict(cnt=cnt, valid=int(cnt >= mot_th),
                     min_x=int(xs.min()), max_x=int(xs.max()),
                     min_y=int(ys.min()), max_y=int(ys.max()))
        else:
            r = dict(cnt=0, valid=0, min_x=0, max_x=0, min_y=0, max_y=0)
        out.append(r)
        prev = m.copy()
    return out


# ---------------- 整链封装（图像 -> 掩码/质心/框） ----------------

def full_detect(rgb_img, preset_idx=0):
    """rgb_img: HxWx3 uint8。返回(掩码, 平滑掩码, 质心结果)。"""
    r = rgb_img[:, :, 0].ravel().astype(np.int32)
    g = rgb_img[:, :, 1].ravel().astype(np.int32)
    b = rgb_img[:, :, 2].ravel().astype(np.int32)
    h, s, v = rgb_to_hsv(r, g, b)
    mask = hsv_threshold(h, s, v, preset_idx).reshape(IMG_H, IMG_W)
    sm = smooth3x3_majority(mask)
    return mask, sm, centroid_bbox(sm)


if __name__ == '__main__':
    # 自检：三基色 + 消色差 + 位精确 LUT
    for rgb, exp_h in [((224, 32, 32), 0), ((32, 224, 32), 120),
                       ((32, 32, 224), 240), ((128, 128, 128), 0),
                       ((255, 255, 255), 0), ((0, 0, 0), 0)]:
        hh, ss, vv = rgb_to_hsv(*rgb)
        print(f"RGB{rgb} -> H={int(hh)} S={int(ss)} V={int(vv)} (期望H~={exp_h})")
    print("inv[1..8] =", _INV[1:9].tolist())
    print("golden_model 自检完成")
\end{Verbatim}
\subsection{gen\_vectors.py}
\fileintro{向量生成器：16 TB 的激励/期望 + 3 个生成 RTL（倒数表/字库/跑马灯表）+ OSD 金标准渲染器（位精确复刻 osd\_track 绘制规则）。}
\begin{Verbatim}[numbers=left,numbersep=5pt,frame=single,fontsize=\footnotesize,breaklines=true,breakanywhere=true]
# -*- coding: utf-8 -*-
"""
gen_vectors.py —— 生成全部测试向量（ModelSim $readmemh 格式）与期望输出，
并自动生成两份 RTL：hsv_inv_lut.v（HSV 倒数查找表）、font_rom_8x16.v（OSD 字库）。
所有期望值由 golden_model.py 位精确算出。用法：
    python gen_vectors.py all       # 生成全部
    python gen_vectors.py list      # 列出各目标
"""
import os
import sys
import numpy as np

sys.path.insert(0, os.path.dirname(os.path.abspath(__file__)))
from golden_model import (IMG_W, IMG_H, CENTER_X, CENTER_Y, HSV_PRESETS,
                          rgb_to_hsv, hsv_threshold, smooth3x3_majority,
                          centroid_bbox, gimbal_ctrl_step, full_detect)

ROOT = os.path.dirname(os.path.dirname(os.path.abspath(__file__)))
RNG = np.random.default_rng(20260918)  # 固定种子，保证可复现


def w(path, lines):
    os.makedirs(os.path.dirname(path), exist_ok=True)
    with open(path, 'w') as f:
        f.write('\n'.join(lines) + '\n')
    print(f"  写 {path} ({len(lines)} 行)")


def sp(step):
    return os.path.join(ROOT, '03_vision', step, '03_project')


# ---------------- 生成 RTL：HSV 倒数 LUT ----------------
def gen_inv_lut():
    lines = ["// 自动生成：tools/gen_vectors.py  inv[d]=floor(3840/d) d=1..255",
             "// 供 rgb2hsv.v 求色调用，综合为 Block ROM",
             "module hsv_inv_lut (", "    input  wire [7:0]  I_d,", "    output reg  [11:0] O_inv", ");",
             "always @(*) begin", "    case (I_d)"]
    for d in range(1, 256):
        lines.append(f"        8'd{d}: O_inv = 12'd{3840 // d};")
    lines += ["        default: O_inv = 12'd0;", "    endcase", "end", "endmodule"]
    w(os.path.join(sp('step2_color_centroid'), 'rtl', 'hsv_inv_lut.v'), lines)


# ---------------- 生成 RTL：8x16 字库 ----------------
FONT_CHARS = " 0123456789ABCDEFGHIJKLMNOPQRSTUVWXYZ:.-+/"
FONT_TEXTS = ["LOCK", "LOST", "TRK", "X", "Y"]


_FONT_CACHE = {}


def gen_font_glyph(ch):
    """单个字符的 8x16 点阵（带缓存），与 font_rom_8x16.v 逐位一致。"""
    if not _FONT_CACHE:
        import cv2
        for c in FONT_CHARS:
            img = np.zeros((16, 8), dtype=np.uint8)
            scale = 0.55 if c != ' ' else 0.1
            t = cv2.putText(img, c, (0, 14), cv2.FONT_HERSHEY_SIMPLEX,
                            scale, 255, 1, cv2.LINE_AA)
            _FONT_CACHE[c] = (t >= 110).astype(np.uint8)
    return _FONT_CACHE.get(ch, np.zeros((16, 8), dtype=np.uint8))


def gen_font():
    """cv2 Hershey 轮廓字渲染到 8x16 点阵，阈值化为位图 -> case 表 ROM。"""
    glyph = {ch: gen_font_glyph(ch) for ch in FONT_CHARS}
    lines = ["// 自动生成：tools/gen_vectors.py (cv2 Hershey->8x16 点阵)",
             "// 地址 {char[7:0], row[3:0]} -> 该行 8 位点阵（MSB=最左像素）",
             "module font_rom_8x16 (", "    input  wire [11:0] I_addr,",
             "    output reg  [7:0]  O_bits", ");", "always @(*) begin",
             "    case (I_addr)"]
    for ch in FONT_CHARS:
        bits = glyph[ch]
        for r in range(16):
            val = 0
            for c in range(8):
                val = (val << 1) | int(bits[r, c])
            lines.append(f"        12'h{(ord(ch) << 4 | r):03X}: O_bits = 8'h{val:02X};  // '{ch}' row{r}")
    lines += ["        default: O_bits = 8'h00;", "    endcase", "end", "endmodule"]
    w(os.path.join(sp('step4_osd_integration'), 'rtl', 'font_rom_8x16.v'), lines)
    return glyph


# ---------------- 场景生成 ----------------
def scene_red_square(x0=560, y0=300, size=90, noise=0):
    img = np.full((IMG_H, IMG_W, 3), 128, dtype=np.uint8)      # 灰底
    img[y0:y0 + size, x0:x0 + size] = (224, 32, 32)            # 红块 RGB
    if noise:
        ys = RNG.integers(0, IMG_H, noise)
        xs = RNG.integers(0, IMG_W, noise)
        cols = RNG.integers(0, 256, (noise, 3))
        img[ys, xs] = cols
    return img


def scene_two_objects():
    img = np.full((IMG_H, IMG_W, 3), 100, dtype=np.uint8)
    img[200:320, 300:420] = (224, 32, 32)     # 红
    img[430:560, 800:940] = (200, 32, 200)    # 品红（3扇区算法会混入 H=0 扇区，靠 s_lo 抑制：S=168 仍过 -> 检验质心偏移告警场景）
    return img


def scene_small_target():
    img = np.full((IMG_H, IMG_W, 3), 128, dtype=np.uint8)
    img[100:115, 100:115] = (224, 32, 32)     # 15x15=225 >= MIN_PIXELS
    return img


def scene_tiny_target():
    img = np.full((IMG_H, IMG_W, 3), 128, dtype=np.uint8)
    img[300:303, 300:303] = (224, 32, 32)     # 9 px < MIN_PIXELS -> invalid
    return img


def scene_to_beats(img):
    """图像 -> 96bit(4像素/拍) AXI-S beat hex 行 + tuser/tlast 标记行。"""
    flat = img.reshape(-1, 3)
    lines, flags = [], []
    n = flat.shape[0]
    for i in range(0, n, 4):
        px = flat[i:i + 4]
        while len(px) < 4:
            px = np.vstack([px, px[-1]])
        val = 0
        for k, (r, g, b) in enumerate(px):
            val |= (int(r) << (88 - 24 * k)) | (int(g) << (80 - 24 * k)) | (int(b) << (72 - 24 * k))
        lines.append(f"{val:024X}")
        beat_idx = (i // 4) % (IMG_W // 4)
        line_idx = (i // 4) // (IMG_W // 4)
        tuser = 1 if i == 0 else 0
        tlast = 1 if beat_idx == (IMG_W // 4 - 1) else 0
        flags.append(f"{(tuser << 1) | tlast:X}")   # 单字符（2bit 向量限 1 位数字）
    return lines, flags



TPG_BARS = [(64,64,64),(128,128,128),(192,192,192),(224,224,224),
            (32,32,224),(32,224,32),(16,96,16),(16,16,96)]

def tpg_frame(box_x=500, box_y=300, box_en=True, box_w=80, box_h=80):
    """位精确复刻 tpg_720p.v 的输出图案。"""
    img = np.zeros((IMG_H, IMG_W, 3), dtype=np.int64)
    for x in range(IMG_W):
        c = TPG_BARS[min(x // 160, 7)]
        img[:, x, :] = c
    if box_en:
        img[box_y:box_y+box_h, box_x:box_x+box_w] = (224, 32, 32)
    return img


def gen_tb_tpg():
    d = os.path.join(sp('step1_video_link'), 'sim')
    img = tpg_frame()
    w(os.path.join(d, 'tpg_exp.hex'),
      [f"{p[0]:02X}{p[1]:02X}{p[2]:02X}" for p in img.reshape(-1, 3)])
    return img


# ---------------- 各 TB 向量 ----------------

def gen_tb_rgb2hsv():
    d = os.path.join(sp('step2_color_centroid'), 'sim')
    n = 10000
    rgb = RNG.integers(0, 256, (n, 3))
    corners = np.array([[0, 0, 0], [255, 255, 255], [255, 0, 0], [0, 255, 0],
                        [0, 0, 255], [255, 255, 0], [0, 255, 255], [255, 0, 255],
                        [128, 128, 128], [224, 32, 32]], dtype=np.int64)
    rgb = np.vstack([corners, rgb])
    h, s, v = rgb_to_hsv(rgb[:, 0], rgb[:, 1], rgb[:, 2])
    w(os.path.join(d, 'rgb2hsv_in.hex'), [f"{r:02X}{g:02X}{b:02X}" for r, g, b in rgb])
    w(os.path.join(d, 'rgb2hsv_exp.hex'), [f"{hh:03X}{ss:02X}{vv:02X}" for hh, ss, vv in zip(h, s, v)])


def gen_tb_threshold():
    d = os.path.join(sp('step2_color_centroid'), 'sim')
    n = 4000
    h = RNG.integers(0, 360, n)
    s = RNG.integers(0, 256, n)
    v = RNG.integers(0, 256, n)
    pres = np.arange(n) % 8
    m = np.array([hsv_threshold(h[i], s[i], v[i], pres[i]) for i in range(n)])
    w(os.path.join(d, 'threshold_in.hex'),
      [f"{hh:03X}{ss:02X}{vv:02X}{p:01X}" for hh, ss, vv, p in zip(h, s, v, pres)])
    w(os.path.join(d, 'threshold_exp.hex'), [str(int(x)) for x in m])


def gen_tb_smooth():
    d = os.path.join(sp('step2_color_centroid'), 'sim')
    W, H, FRAMES = 64, 48, 4
    masks = (RNG.random((FRAMES, H, W)) < 0.35).astype(np.uint8)
    for f in range(FRAMES):
        y0, x0 = 10 + f * 3, 8 + f * 5
        masks[f, y0:y0 + 12, x0:x0 + 14] = 1
    lines_in, lines_exp = [], []
    for f in range(FRAMES):
        sm = smooth3x3_majority(masks[f], W, H)
        lines_in += [f"{int(x)}" for x in masks[f].ravel()]
        lines_exp += [f"{int(x)}" for x in sm.ravel()]
    w(os.path.join(d, 'smooth_in.hex'), lines_in)
    w(os.path.join(d, 'smooth_exp.hex'), lines_exp)
    with open(os.path.join(d, 'smooth_geo.txt'), 'w') as fh:
        fh.write(f"{W} {H} {FRAMES}\n")


def gen_tb_centroid():
    """迷你几何（128x72）质心单元测试：实心块 / 噪声+小块 / 微块(无效)。"""
    from golden_model import MIN_PIXELS as MP
    d = os.path.join(sp('step3_placeholder'), 'sim')  # 占位，下面覆盖
    d = os.path.join(sp('step2_color_centroid'), 'sim')
    W, H = 128, 72
    masks = []
    m = np.zeros((H, W), dtype=np.uint8)
    m[20:60, 30:120] = 1
    masks.append(m)                                   # 帧0: 90x40 实心块
    m2 = np.zeros((H, W), dtype=np.uint8)
    ys = RNG.integers(0, H, 150)
    xs = RNG.integers(0, W, 150)
    m2[ys, xs] = 1
    m2[10:20, 90:120] = 1
    masks.append(m2)                                  # 帧1: 噪声+30x10 块
    m3 = np.zeros((H, W), dtype=np.uint8)
    m3[30:33, 30:33] = 1                              # 帧2: 9px < 门限
    masks.append(m3)
    lines_in, lines_exp = [], []
    for m in masks:
        r = centroid_bbox(m, W, H, MP)
        lines_in += [f"{int(x)}" for x in m.ravel()]
        lines_exp.append(f"{r['valid']:X}{r['count']:05X}{r['cx']:03X}{r['cy']:03X}"
                         f"{r['min_x']:03X}{r['max_x']:03X}{r['min_y']:03X}{r['max_y']:03X}")
    w(os.path.join(d, 'centroid_in.hex'), lines_in)
    w(os.path.join(d, 'centroid_exp.txt'), lines_exp)
    with open(os.path.join(d, 'centroid_geo.txt'), 'w') as fh:
        fh.write(f"{W} {H}" + chr(10))


def gen_tb_chain():
    """detect 链 TB：真实场景（灰底+红块+噪声+品红干扰物），3 帧。"""
    d = os.path.join(sp('step2_color_centroid'), 'sim')
    scenes = [scene_red_square(560, 300, 90, noise=1500),
              scene_two_objects(),
              scene_small_target()]
    beats, flags, exp = [], [], []
    for img in scenes:
        b, f = scene_to_beats(img)
        beats += b
        flags += f
        _, sm, r = full_detect(img)
        exp.append(f"{r['valid']:X}{r['count']:05X}{r['cx']:03X}{r['cy']:03X}"
                   f"{r['min_x']:03X}{r['max_x']:03X}{r['min_y']:03X}{r['max_y']:03X}")
    w(os.path.join(d, 'chain_beats.hex'), beats)
    w(os.path.join(d, 'chain_flags.hex'), flags)
    w(os.path.join(d, 'chain_exp.txt'), exp)


def gen_tb_gimbal():
    d = os.path.join(sp('step3_closed_loop'), 'sim')
    rows, exp = [], []
    # 轴参数（与 gimbal_ctrl.v 默认参数一致）
    KPX, DZX, SLEW, SMIN, SMAX, CX = 20, 25, 100, 700, 2300, CENTER_X
    KPY, DZY, CY = 18, 20, CENTER_Y
    dx = dy = 1500
    dpx = dpy = 0
    lost = 1
    cases = [(1, 840, 360), (1, 840, 360), (1, 900, 300), (1, 700, 200),
             (1, 641, 361), (0, 0, 0), (0, 0, 0), (1, 100, 600),
             (1, 640, 360), (1, 660, 350)]
    for valid, cx, cy in cases:
        rows.append(f"{valid:X}{cx:03X}{cy:03X}")
        if valid:
            lost = 0
            ex, ey = cx - CX, cy - CY
            ndx = gimbal_ctrl_step(ex, dx, KPX, 0, 0, DZX, SLEW, SMIN, SMAX)
            ndy = gimbal_ctrl_step(ey, dy, KPY, 0, 0, DZY, SLEW, SMIN, SMAX)
            dx, dy = ndx, ndy
            dpx = 1 if abs(ex) <= DZX else 0
            dpy = 1 if abs(ey) <= DZY else 0
        else:
            lost += 1
            dpx = dpy = 0
        exp.append(f"{dx:03X}{dy:03X}{lost:04X}{dpx:X}{dpy:X}")
    w(os.path.join(d, 'gimbal_stim.txt'), rows)
    w(os.path.join(d, 'gimbal_exp.txt'), exp)


def render_osd(img, state, disp_cnt):
    """位精确复刻 osd_track.v v3（双缓冲版）的叠加规则。
    state: None = 空状态（复位影子，无文本无框，跑马灯照滚）；
           dict = {valid,cx,cy,min_x,max_x,min_y,max_y, lat,preset,algo,
                   mot: None 或 dict(cnt,valid,min_x,max_x,min_y,max_y)}
    disp_cnt: 帧序计数器（RTL 在 sof 拍 +1，帧 0 绘制期间值为 1）。
    绘制优先级 文字L1>L2>跑马灯>十字>目标框>运动框>底图。"""
    out = img.copy().astype(np.int64)
    H, W = img.shape[:2]
    TXT_X0, TXT_Y0, TXT2_Y0, TXT3_Y0 = 16, 16, 40, 64

    def glyph_bits(ch):
        return gen_font_glyph(ch)

    def draw_text(out, chars, x0_band, y0_band):
        for i, ch in enumerate(chars):
            if ch == ' ':
                continue
            bits = glyph_bits(ch)
            gx, gy = x0_band + i * 8, y0_band
            for ry in range(16):
                row_bits = bits[ry]
                for rx in range(8):
                    if row_bits[rx]:
                        yy, xx = gy + ry, gx + rx
                        if yy < H and xx < W:
                            out[yy, xx] = (255, 255, 255)

    # ---- 文本内容（逐字对应 RTL txt1/txt2 赋值） ----
    if state is not None:
        v = state['valid']
        algo = state.get('algo', 0)
        pre = state.get('preset', 0)
        cx, cy = state['cx'], state['cy']
        t1 = ('EDG' if algo else 'TRK') + str(pre) + ' '
        t1 += 'LOCK' if v else 'LOST'
        t1 += ' X' + (f"{cx:04d}" if v else '----') + 'Y' + (f"{cy:04d}" if v else '----')
        lat = min(state.get('lat', 0), 99)
        br, ct = state.get('bright', 0), state.get('contrast', 64)
        b_abs = (~br + 1) & 0x7F if br & 0x40 else br
        mcnt = min(state.get('mcnt', 0), 9999)
        t2 = (f"D{lat:02d} B{'-' if br & 0x40 else '+'}{b_abs:02d} "
              f"C{ct >> 3:02d} M{mcnt:04d} ")
        draw_text(out, t1.ljust(20), TXT_X0, TXT_Y0)
        draw_text(out, t2.ljust(20), TXT_X0, TXT2_Y0)

        if v:
            x0, x1, y0, y1 = state['min_x'], state['max_x'], state['min_y'], state['max_y']
            rows = sorted(set([y0, min(y0 + 1, y1), max(y1 - 1, y0), y1]))
            cols = sorted(set([x0, min(x0 + 1, x1), max(x1 - 1, x0), x1]))
            for y in rows:
                if 0 <= y < H:
                    out[y, max(x0, 0):min(x1 + 1, W)] = (255, 0, 0)
            for x in cols:
                if 0 <= x < W:
                    out[max(y0, 0):min(y1 + 1, H), x] = (255, 0, 0)
            for d in list(range(-12, -2)) + list(range(3, 13)):
                if 0 <= cy < H and 0 <= cx + d < W:
                    out[cy, cx + d] = (0, 255, 0)
                if 0 <= cx < W and 0 <= cy + d < H:
                    out[cy + d, cx] = (0, 255, 0)
        mot = state.get('mot')
        if mot and mot.get('valid'):
            x0, x1, y0, y1 = mot['min_x'], mot['max_x'], mot['min_y'], mot['max_y']
            rows = sorted(set([y0, min(y0 + 1, y1), max(y1 - 1, y0), y1]))
            cols = sorted(set([x0, min(x0 + 1, x1), max(x1 - 1, x0), x1]))
            for y in rows:
                if 0 <= y < H:
                    out[y, max(x0, 0):min(x1 + 1, W)] = (255, 128, 0)
            for x in cols:
                if 0 <= x < W:
                    out[max(y0, 0):min(y1 + 1, H), x] = (255, 128, 0)

    # ---- 跑马灯（disp_cnt 驱动，与目标包无关） ----
    mar_off = (disp_cnt & 0xFF) << 1
    mar_msg = 'ANLOGIC PH1P35 FPGA VISION PTZ TRACKER 2026                    '
    glyph_cache = {}
    for y in range(TXT3_Y0, min(TXT3_Y0 + 16, H)):
        row = y - TXT3_Y0
        for x in range(TXT_X0, W):
            pos = (x - TXT_X0 + mar_off) & 0xFFF
            # 外脑评审修复同步: 原实现 (pos>>3)&63 %63 在 pos>=512 后字符
            # 索引截断错位（与旧 RTL 同错）；正确公式 (pos>>3) % 63
            ci = (pos >> 3) % 63
            ch = mar_msg[ci]
            if ch == ' ':
                continue
            if ch not in glyph_cache:
                glyph_cache[ch] = glyph_bits(ch)
            if glyph_cache[ch][row][pos & 7]:
                out[y, x] = (0, 255, 255)
    return out


def gen_tb_osd():
    """OSD TB v2：3 帧（f0 空状态 / f1 换装后 / f2 连续帧），包在 f0 中帧注入。
    金标准 3 份（mar_off 随 disp_cnt 2/4/6 递增），验证双缓冲消撕裂 + 整帧一致。"""
    d = os.path.join(sp('step4_osd_integration'), 'sim')
    img = scene_red_square(560, 300, 90, noise=800)
    mask, sm, r = full_detect(img)
    state = {'valid': r['valid'], 'cx': r['cx'], 'cy': r['cy'],
             'min_x': r['min_x'], 'max_x': r['max_x'],
             'min_y': r['min_y'], 'max_y': r['max_y'],
             # lat=1：注入帧 fid=0、首显示帧 disp_cnt=1 -> 延迟 1 帧；
             # 真实链路稳态同样为 1（第 N 帧结果在第 N+1 帧显示）
             'lat': 1, 'preset': 0, 'algo': 0, 'mcnt': 64,
             # 运动框覆盖（外脑评审补：原 mot=None，运动框从未被像素级测试）：
             # 运动框 (200,200)-(260,240) 与目标框 (~560..649, 301..390) 明显不同，
             # 若 RTL 误用目标框坐标画运动框，逐像素对比必然失败
             'mot': {'valid': 1, 'cnt': 64,
                     'min_x': 200, 'max_x': 260,
                     'min_y': 200, 'max_y': 240},
             'bright': 0, 'contrast': 64}
    w(os.path.join(d, 'osd_base.hex'),
      [f"{p[0]:02X}{p[1]:02X}{p[2]:02X}" for p in img.reshape(-1, 3)])
    for f, dc in ((0, 1), (1, 2), (2, 3)):
        st = None if f == 0 else state
        out = render_osd(img, st, dc)
        w(os.path.join(d, f'osd_exp_f{f}.hex'),
          [f"{p[0]:02X}{p[1]:02X}{p[2]:02X}" for p in out.reshape(-1, 3)])
    with open(os.path.join(d, 'osd_coord.txt'), 'w') as fh:
        fh.write(f"{r['valid']:X}{r['cx']:03X}{r['cy']:03X}{r['min_x']:03X}"
                 f"{r['max_x']:03X}{r['min_y']:03X}{r['max_y']:03X}\n")
    return r


def gen_tb_motion():
    """帧差运动检测 TB（64x48 x 6 帧）：静止->出现->移动->消失->亚门限->末像素边界。
    期望值由 golden_model.frame_diff_motion 位精确生成。"""
    from golden_model import frame_diff_motion
    d = os.path.join(sp('step5_expansion'), 'sim')
    W, H, TH = 64, 48, 12
    m = np.zeros((H, W), dtype=np.uint8)
    frames = [m.copy()]                                       # f0 静止背景（只填缓存）
    m[10:18, 10:18] = 1; frames.append(m.copy())              # f1 出现 8x8 块
    m = np.zeros((H, W), dtype=np.uint8); m[20:28, 30:38] = 1
    frames.append(m.copy())                                   # f2 移动（旧位消失+新位出现）
    m = np.zeros((H, W), dtype=np.uint8); frames.append(m.copy())  # f3 消失
    m = np.zeros((H, W), dtype=np.uint8)
    for yy, xx in ((5, 5), (40, 50), (30, 20), (10, 60), (44, 3)):
        m[yy, xx] = 1
    frames.append(m.copy())                                   # f4 5px 亚门限（valid=0）
    # f5 末像素边界（外脑评审新增）：仅 (W-1,H-1) 一个差分像素。
    # 用于回归"帧末锁存漏末像素"缺陷（v3.1 已修：锁存延迟一拍）。
    m = np.zeros((H, W), dtype=np.uint8); m[H-1, W-1] = 1
    frames.append(m.copy())
    lines_in = []
    exp = []
    res = frame_diff_motion(frames, TH)
    for f, mm in enumerate(frames):
        lines_in += [f"{int(v)}" for v in mm.ravel()]
        r = res[f]
        if r is None:                     # 首帧无输出 -> 期望全 0
            exp.append("000000000000000")
            continue
        packed = (r['cnt'] << 43) | (r['valid'] << 42) | (r['min_x'] << 31) \
                 | (r['max_x'] << 20) | (r['min_y'] << 10) | r['max_y']
        exp.append(f"{packed:015X}")
    w(os.path.join(d, 'motion_in.hex'), lines_in)
    w(os.path.join(d, 'motion_exp.txt'), exp)
    with open(os.path.join(d, 'motion_geo.txt'), 'w') as fh:
        fh.write(f"{W} {H} {TH} {len(frames)}" + chr(10))


def gen_all():
    print("== 生成 HSV 倒数 LUT RTL ==")
    gen_inv_lut()
    print("== 生成 8x16 字库 RTL ==")
    gen_font()
    print("== TB 向量 ==")
    gen_tb_rgb2hsv()
    gen_tb_threshold()
    gen_tb_smooth()
    gen_tb_centroid()
    gen_tb_chain()
    gen_tb_gimbal()
    gen_tb_osd()
    gen_tb_tpg()
    gen_tb_motion()
    print("全部向量生成完成")


if __name__ == '__main__':
    gen_all()
\end{Verbatim}
\subsection{gen\_raw10\_sensor.py}
\fileintro{虚拟 RAW10 传感器（BGGR 灰底+红块+0.2\% 确定性噪声，4 帧）：打包先转 Python int 再移位（int32 溢出坑）；R 位点相位为厂商 demosaic 实测定稿。}
\begin{Verbatim}[numbers=left,numbersep=5pt,frame=single,fontsize=\footnotesize,breaklines=true,breakanywhere=true]
# -*- coding: utf-8 -*-
"""
gen_raw10_sensor.py —— 虚拟 SC500CS RAW10 传感器（嵌入式集成仿真用）
生成 BGGR Bayer 场景（灰底 + 红块），打包为 isp_top 输入格式：
  40bit/拍 = 4 像素 x RAW10，P0=[39:30] P1=[29:20] P2=[19:10] P3=[9:0]
  每行 320 拍，tuser=帧首拍，tlast=行尾拍（与官方 raw10_unpacket_2lane 输出对流对接）
输出：04_embedded/03_project/sim/{raw10_beats.hex, raw10_meta.txt}
"""
import os
import numpy as np

W, H = 1280, 720
BOX = (500, 300, 80, 80)          # x0,y0,w,h —— 与 tb_vision_top 的 TPG 块相近
ROOT = os.path.dirname(os.path.dirname(os.path.abspath(__file__)))
OUT = os.path.join(ROOT, '04_embedded', '03_project', 'sim')

def gen():
    os.makedirs(OUT, exist_ok=True)
    # 10bit 灰底；BGGR：偶行 B G B G...，奇行 G R G R...
    raw = np.full((H, W), 512, dtype=np.int32)
    x0, y0, w, h = BOX
    # 红块：R 位点=900，G/B 位点=64（R 位点 = 奇行奇列）
    ys, xs = np.mgrid[y0:y0+h, x0:x0+w]
    # 红块: 对角棋盘布 R=900(与厂商 demosaic 的 Bayer 相位约定对齐,
    # 实测定稿——标准 BGGR 单相位会在插值后出现隔行绿纹)
    # R 位点=偶行偶列：与厂商 demosaic 的**实测**相位对齐（2026-10-03 定稿）。
    # 注：标准命名下 BGGR 的 R 应在奇行奇列；厂商链路对行/列相位有一格偏移，
    # 以实测定稿为准——红块经该链路渲染为红色即为验证（见板级使用指导第 5 节）。
    r_sites = ((ys & 1) == 0) & ((xs & 1) == 0)
    region_r = np.where(r_sites, 900, 64)
    raw[y0:y0+h, x0:x0+w] = region_r
    # 少量确定性噪声（LFSR 风格，可复现）
    rng = np.random.default_rng(500)
    noise_mask = rng.random((H, W)) < 0.002
    raw[noise_mask] = rng.integers(100, 900, int(noise_mask.sum()))

    # 打包：每 4 像素一拍
    flat = raw.ravel()
    beats, metas = [], []
    n = W * H
    beats_per_line = W // 4
    for i in range(0, n, 4):
        # 必须先转 Python int 再移位：numpy int32 左移 30 位会溢出
        # （背景拍高 10bit 被截断、噪声拍变负数，readmemh 遇 '-' 终止加载全 X）
        p0, p1, p2, p3 = int(flat[i]), int(flat[i+1]), int(flat[i+2]), int(flat[i+3])
        val = (p0 << 30) | (p1 << 20) | (p2 << 10) | p3
        beats.append(f"{val:010X}")
        beat_in_line = (i // 4) % beats_per_line
        sof = 1 if i == 0 else 0
        eol = 1 if beat_in_line == beats_per_line - 1 else 0
        metas.append(f"{(sof << 1) | eol:X}")   # 单字符（2bit 向量限 1 位数字）
    FRAMES = 4    # 集成仿真需连续 4 帧
    with open(os.path.join(OUT, 'raw10_beats.hex'), 'w') as f:
        f.write('\n'.join(beats * FRAMES) + '\n')
    with open(os.path.join(OUT, 'raw10_meta.txt'), 'w') as f:
        f.write('\n'.join(metas * FRAMES) + '\n')
    # 期望检测框（R 位点上取值后经 demosaic/AWB 的框位置不变）
    with open(os.path.join(OUT, 'raw10_expect.txt'), 'w') as f:
        f.write(f"{x0} {y0} {x0+w-1} {y0+h-1} {x0+w//2} {y0+h//2}\n")
    print(f"RAW10 虚拟传感器：{len(beats)} 拍，块 {BOX}，输出 {OUT}")

if __name__ == '__main__':
    gen()
\end{Verbatim}
\subsection{run\_all.py}
\fileintro{一键回归执行器：18 个 TB 的编译/运行/逐位比对/汇总。}
\begin{Verbatim}[numbers=left,numbersep=5pt,frame=single,fontsize=\footnotesize,breaklines=true,breakanywhere=true]
# -*- coding: utf-8 -*-
"""
run_all.py —— 视觉软件包整体测试调度器
流程：编译全部 RTL+TB（vlog）-> 逐 TB 仿真（vsim -c）-> 解析 TB_RESULT
      -> 帧级 hex 比对（TPG/OSD/整链）-> 汇总打印 + 写迭代测试报告。
用法：
    python run_all.py            # 全量
    python run_all.py tb_ser96to24 [tb_xxx ...]   # 只跑指定 TB
"""
import os
import re
import sys
import glob
import hashlib
import subprocess
import datetime

sys.path.insert(0, os.path.dirname(os.path.abspath(__file__)))
from gen_vectors import tpg_frame, render_osd, sp
from golden_model import full_detect

ROOT = os.path.dirname(os.path.dirname(os.path.abspath(__file__)))
MSIM = r'D:\ModelSim\win64'
VLOG = os.path.join(MSIM, 'vlog.exe')
VSIM = os.path.join(MSIM, 'vsim.exe')
VLIB = os.path.join(MSIM, 'vlib.exe')

import os as _os
P1 = sp('step1_video_link')
P2 = sp('step2_color_centroid')
P3 = sp('step3_closed_loop')
P4 = sp('step4_osd_integration')

def r(step, *parts):
    return os.path.join(ROOT, '03_vision', step, '03_project', *parts)

STEP1_RTL = [r('step1_video_link', 'rtl', f) for f in
             ['ser96to24.v', 'tpg_720p.v', 'sim_video_src.v', 'frame_dump24.v',
              'pulse_sync.v', 'cdc_4phase.v']]
STEP2_RTL = [r('step2_color_centroid', 'rtl', f) for f in
             ['hsv_inv_lut.v', 'rgb2hsv.v', 'hsv_threshold.v',
              'bin_smooth3x3.v', 'centroid_accum.v', 'detect_top.v']]
STEP3_RTL = [r('step3_closed_loop', 'rtl', f) for f in ['gimbal_ctrl.v', 'pwm_servo.v']]
STEP5_RTL = [r('step5_expansion', 'rtl', f) for f in
             ['edge_detect_96.v', 'video_adj.v', 'motion_detect.v', 'marquee_rom_64.v']]
STEP4_RTL = [r('step4_osd_integration', 'rtl', f) for f in
             ['font_rom_8x16.v', 'osd_track.v', 'vision_top.v']]
EMB_RTL = [os.path.join(ROOT, '04_embedded', '03_project', 'rtl', f) for f in
           ['led_buzzer.v', 'uart_telemetry.v']]

# name -> (sim_cwd_dir, 编译文件列表, 输出工件检查)
TBS = {
    'tb_ser96to24':  (r('step1_video_link', 'sim'), STEP1_RTL + [r('step1_video_link', 'tb', 'tb_ser96to24.v')], []),
    'tb_tpg':        (r('step1_video_link', 'sim'), STEP1_RTL + [r('step1_video_link', 'tb', 'tb_tpg.v')],
                      [('frame', 'tpg_act_0.hex', 'tpg_exp.hex')]),
    'tb_pulse_sync': (r('step1_video_link', 'sim'), STEP1_RTL + [r('step1_video_link', 'tb', 'tb_pulse_sync.v')], []),
    'tb_cdc4phase':  (r('step1_video_link', 'sim'), STEP1_RTL + [r('step1_video_link', 'tb', 'tb_cdc4phase.v')], []),
    'tb_cdc_reset_skew': (r('step1_video_link', 'sim'),
                          [os.path.join(ROOT, '03_vision', 'step1_video_link', '03_project', 'rtl', 'cdc_4phase.v'),
                           r('step1_video_link', 'tb', 'tb_cdc_reset_skew.v')], []),
    'tb_rgb2hsv':    (r('step2_color_centroid', 'sim'), STEP1_RTL[:1] + STEP2_RTL[:2] + [r('step2_color_centroid', 'tb', 'tb_rgb2hsv.v')], []),
    'tb_hsv_threshold': (r('step2_color_centroid', 'sim'), STEP2_RTL[:4] + [r('step2_color_centroid', 'tb', 'tb_hsv_threshold.v')], []),
    'tb_bin_smooth': (r('step2_color_centroid', 'sim'), STEP2_RTL[3:4] + [r('step2_color_centroid', 'tb', 'tb_bin_smooth.v')], []),
    'tb_centroid':   (r('step2_color_centroid', 'sim'), STEP2_RTL[4:5] + [r('step2_color_centroid', 'tb', 'tb_centroid.v')], []),
    'tb_detect_chain': (r('step2_color_centroid', 'sim'), STEP1_RTL[:1] + STEP1_RTL[2:3] + STEP2_RTL + [r('step2_color_centroid', 'tb', 'tb_detect_chain.v')], []),
    'tb_gimbal':     (r('step3_closed_loop', 'sim'), STEP3_RTL[:1] + [r('step3_closed_loop', 'tb', 'tb_gimbal.v')], []),
    'tb_pwm':        (r('step3_closed_loop', 'sim'), STEP3_RTL[1:2] + [r('step3_closed_loop', 'tb', 'tb_pwm.v')], []),
    'tb_osd_track':  (r('step4_osd_integration', 'sim'), STEP1_RTL[5:6] + STEP4_RTL[:2] + STEP5_RTL[3:4] + [r('step4_osd_integration', 'tb', 'tb_osd_track.v')],
                      [('frame', 'osd_act_f0.hex', 'osd_exp_f0.hex'),
                       ('frame', 'osd_act_f1.hex', 'osd_exp_f1.hex'),
                       ('frame', 'osd_act_f2.hex', 'osd_exp_f2.hex')]),
    'tb_vision_top': (r('step4_osd_integration', 'sim'), STEP1_RTL + STEP2_RTL + STEP3_RTL + STEP4_RTL + STEP5_RTL + EMB_RTL + [r('step4_osd_integration', 'tb', 'tb_vision_top.v')],
                      []),  # v2 OSD 金标准渲染器待更新，功能断言由 TB 内部完成
    'tb_closed_loop': (r('step4_osd_integration', 'sim'), STEP1_RTL + STEP2_RTL + STEP3_RTL + STEP4_RTL + STEP5_RTL + EMB_RTL + [r('step4_osd_integration', 'tb', 'tb_closed_loop.v')], []),
    'tb_algo_switch': (r('step4_osd_integration', 'sim'), STEP1_RTL + STEP2_RTL + STEP3_RTL + STEP4_RTL + STEP5_RTL + EMB_RTL + [r('step4_osd_integration', 'tb', 'tb_algo_switch.v')], []),
    # ---------------- 嵌入式方向（04_embedded） ----------------
    'tb_uart_telemetry': (
        os.path.join(ROOT, '04_embedded', '03_project', 'sim'),
        STEP1_RTL[:1] + STEP1_RTL[4:6] + STEP2_RTL[1:6] + STEP2_RTL[0:1] + STEP3_RTL + STEP4_RTL + [
            os.path.join(ROOT, '04_embedded', '03_project', 'rtl', 'uart_telemetry.v'),
            os.path.join(ROOT, '04_embedded', '03_project', 'tb', 'uart_rx_mon.v'),
            os.path.join(ROOT, '04_embedded', '03_project', 'tb', 'tb_uart_telemetry.v')], []),
    # ---------------- v2 拓展（step5） ----------------
    'tb_motion_detect': (r('step5_expansion', 'sim'),
                         STEP5_RTL[:4] + [r('step5_expansion', 'tb', 'tb_motion_detect.v')], []),
    # TD 阶段项（默认回归外，见 04_embedded 报告）：
    #   tb_i2c_config      —— 需与真机 I2C 波形对照后收敛位相
    #   tb_system_integration —— 官方 demosaic 的 tuser 链在纯 RTL 仿真不激活（TD/真机项）
}


def sha_lines(path):
    h = hashlib.sha256()
    with open(path, 'r') as f:
        for line in f:
            h.update(line.strip().upper().encode())
    return h.hexdigest()


def check_frame(cwd, act, exp, log):
    a, e = os.path.join(cwd, act), os.path.join(cwd, exp)
    if not os.path.exists(a):
        return False, f'{act} 不存在'
    if sha_lines(a) == sha_lines(e):
        return True, f'{act} 与金标准逐像素一致'
    # 找第一个差异行
    with open(a) as fa, open(e) as fe:
        for i, (la, le) in enumerate(zip(fa, fe)):
            if la.strip().upper() != le.strip().upper():
                return False, f'{act} 首个差异在第 {i} 像素: got={la.strip()} exp={le.strip()}'
        return False, f'{act} 行数不一致'


def check_vision(cwd, act, _):
    # v2：OSD 三行+跑马灯的金标准渲染器尚未更新，暂时只检查转储完整性
    a = os.path.join(cwd, act)
    if not os.path.exists(a):
        return False, f'{act} 不存在'
    n = sum(1 for _ in open(a))
    if n != 921600:
        return False, f'帧行数 {n} != 921600'
    with open(os.path.join(cwd, 'vision_duty.txt')) as f:
        dx, dy, lost, status = map(int, f.read().split())
    if not (700 <= dx < 1500):
        return False, f'云台 X 方向错误: duty_x={dx}'
    if dy != 1500:
        return False, f'云台 Y 死区错误: duty_y={dy}'
    if lost != 0:
        return False, f'目标不应丢失: lost={lost}'
    return True, f'帧完整 {n} px; duty=({dx},{dy}) 方向正确; lost=0'
    """整链：OSD 帧 + 云台方向 + 质心正确性。"""
    a = os.path.join(cwd, act)
    if not os.path.exists(a):
        return False, f'{act} 不存在'
    # 期望：TPG 图案 -> 金标准检测 -> 叠加渲染
    img = tpg_frame()
    _, _, res = full_detect(img.astype('uint8'))
    exp_img = render_osd(img, res)
    exp_lines = [f"{p[0]:02X}{p[1]:02X}{p[2]:02X}" for p in exp_img.reshape(-1, 3)]
    with open(a) as f:
        act_lines = [l.strip().upper() for l in f]
    if len(act_lines) != len(exp_lines):
        return False, f'帧行数 {len(act_lines)} != {len(exp_lines)}'
    diff = next((i for i, (x, y) in enumerate(zip(act_lines, exp_lines)) if x != y), None)
    if diff is not None:
        return False, f'帧差异@{diff}: got={act_lines[diff]} exp={exp_lines[diff]}'
    # 云台方向
    with open(os.path.join(cwd, 'vision_duty.txt')) as f:
        dx, dy, lost, status = map(int, f.read().split())
    if not (700 <= dx < 1500):
        return False, f'云台 X 方向错误: duty_x={dx}（目标在中心左侧，应 <1500 且 >=700）'
    if dy != 1500:
        return False, f'云台 Y 死区错误: duty_y={dy}（cy=340, err=-20 恰在 DZ_Y=20 内，应保持 1500）'
    if lost != 0:
        return False, f'目标不应丢失: lost={lost}'
    return True, (f'OSD 帧 921600 像素全一致; duty=({dx},{dy})µs 方向正确; '
                  f'centroid=({res["cx"]},{res["cy"]}) count={res["count"]}')


def run_tb(name):
    cwd, files, checks = TBS[name]
    os.chdir(cwd)
    if not os.path.isdir('work'):
        subprocess.run([VLIB, 'work'], capture_output=True)
    out = {'name': name, 'compile': '', 'sim': '', 'pass': False, 'detail': ''}
    cp = subprocess.run([VLOG, '-quiet'] + files, capture_output=True, text=True,
                        errors='replace')
    out['compile'] = (cp.stdout + cp.stderr).strip()
    if cp.returncode != 0:
        out['detail'] = '编译失败:\n' + out['compile'][-1500:]
        return out
    cs = subprocess.run([VSIM, '-c', '-quiet', name, '-do', 'run -all; quit -f'],
                        capture_output=True, text=True, errors='replace', timeout=3600)
    log = cs.stdout + cs.stderr
    out['sim'] = log
    ok = re.search(r'TB_RESULT: PASS', log) is not None
    msgs = [l for l in log.splitlines()
            if re.search(r'(MISM|_ERR|FATAL|LOCKED|MEAS|FRAME\d|STEP\d|DUTY_REPORT)', l)]
    out['detail'] = '\n'.join(msgs[-12:])
    # 工件检查
    for kind, act, exp in checks:
        if kind == 'frame':
            ok2, msg = check_frame(cwd, act, exp, log)
        else:
            ok2, msg = check_vision(cwd, act, exp)
        out['detail'] += f'\n[{kind}] {msg}'
        ok = ok and ok2
    out['pass'] = ok
    return out


def main():
    only = sys.argv[1:]
    names = only if only else list(TBS.keys())
    print(f'==== 视觉软件包整体测试  {datetime.datetime.now():%Y-%m-%d %H:%M:%S} ====')
    results = []
    for n in names:
        print(f'--- {n} ...', flush=True)
        try:
            res = run_tb(n)
        except subprocess.TimeoutExpired:
            res = {'name': n, 'pass': False, 'compile': '', 'sim': '',
                   'detail': '仿真超时(1h)'}
        tag = 'PASS' if res['pass'] else 'FAIL'
        print(f'    ==> {tag}')
        if res['detail']:
            for line in res['detail'].splitlines():
                print('        ' + line)
        results.append(res)
        os.chdir(ROOT)
    npass = sum(1 for x in results if x['pass'])
    print(f'\n==== 汇总: {npass}/{len(results)} PASS ====')
    for x in results:
        if not x['pass']:
            print(f"  FAIL: {x['name']}")
            if x.get('compile') and 'Error' in x['compile']:
                print('    (编译错误见迭代日志)')
    # 写迭代测试日志
    logdir = os.path.join(ROOT, '03_vision', 'step4_osd_integration', '04_迭代文件')
    os.makedirs(logdir, exist_ok=True)
    stamp = datetime.datetime.now().strftime('%Y%m%d_%H%M%S')
    with open(os.path.join(logdir, f'测试运行_{stamp}.log'), 'w', encoding='utf-8') as f:
        for x in results:
            f.write(f"===== {x['name']} : {'PASS' if x['pass'] else 'FAIL'} =====\n")
            if x.get('compile'):
                f.write('-- compile --\n' + x['compile'][-3000:] + '\n')
            f.write('-- sim tail --\n')
            f.write('\n'.join(x['sim'].splitlines()[-40:]) + '\n\n')
    print(f'日志已写入 04_迭代文件/测试运行_{stamp}.log')
    return 0 if npass == len(results) else 1


if __name__ == '__main__':
    sys.exit(main())
\end{Verbatim}
\subsection{rtl\_model\_check.py}
\fileintro{RTL 周期精确转录模型（调试辅助）：失配时先在 Python 里定位再改 RTL。}
\begin{Verbatim}[numbers=left,numbersep=5pt,frame=single,fontsize=\footnotesize,breaklines=true,breakanywhere=true]
# -*- coding: utf-8 -*-
"""bin_smooth3x3.v v6 的严格周期精确转录，与金标准对拍。"""
import numpy as np
import sys, os
sys.path.insert(0, os.path.dirname(os.path.abspath(__file__)))
from golden_model import smooth3x3_majority

W, H = 64, 48

def rtl_model(flat_in, total):
    x = 0; y = 0
    a_msk = 0
    started = 0
    t0 = t1 = m0 = m1 = b0 = b1 = 0
    line0 = [0]*W; line1 = [0]*W
    outs = []
    for g in range(total):
        i_mask = flat_in[g]
        i_sof = 1 if g % (W*H) == 0 else 0
        i_eol = 1 if (g % W == W-1) else 0
        # 站B（寄存相位 = 令牌 g-1）
        if y == 0:
            t_sel = m_sel = b_sel = i_mask
        elif y == 1:
            t_sel = m_sel = line0[0]; b_sel = i_mask
        else:
            t_sel = line1[0]; m_sel = line0[0]; b_sel = i_mask
        win_main = t1 + t0 + t_sel + m1 + m0 + m_sel + b1 + b0 + b_sel
        win_last = t0 + 2*t_sel + m0 + 2*m_sel + b1 + 2*b0
        main_en = 1 if x >= 1 else 0
        last_en = 1 if (x == 0 and started) else 0
        if 12 <= g <= 15:
            print(f"g={g} x={x} y={y} t1={t1} t0={t0} tsel={t_sel} "
                  f"wm={win_main} main={main_en} om={om} im={i_mask}")
        om = 1 if (main_en and win_main >= 5) else (1 if (last_en and win_last >= 5) else 0)
        outs.append(om)
        # 寄存更新
        t1, t0 = t0, t_sel
        m1, m0 = m0, m_sel
        b1, b0 = b0, b_sel
        if x == 0:
            t1, m1, b1 = t_sel, m_sel, b_sel
        d0_old = line0[0]
        line0 = [a_msk] + line0[:W-1]
        line1 = [d0_old] + line1[:W-1]
        a_msk = i_mask
        if i_sof:
            x = 1; y = 0; started = 1
        elif i_eol:
            x = 0; y = 0 if y == H-1 else y + 1
        else:
            x += 1
    return outs

if __name__ == '__main__':
    d = '../03_vision/step2_color_centroid/03_project/sim'
    fin = np.loadtxt(d+'/smooth_in.hex', dtype=int).tolist()
    fex = np.loadtxt(d+'/smooth_exp.hex', dtype=int)
    total = W*H*4
    outs = rtl_model(fin[:total], total)
    nmis = 0
    for i, om in enumerate(outs):
        if om != fex[i]:
            if nmis < 6:
                y, x = divmod(i, W)
                print(f"MISM out{i} (r{y} c{x}) model={om} exp={fex[i]}")
            nmis += 1
    print(f"总输出 {len(outs)}（期望 {total}），失配 {nmis}")
\end{Verbatim}
\subsection{dump2png.py}
\fileintro{转储帧（hex）-> PNG（含 X 像素统计），板级仿真"画面"肉眼验证用。}
\begin{Verbatim}[numbers=left,numbersep=5pt,frame=single,fontsize=\footnotesize,breaklines=true,breakanywhere=true]
# -*- coding: utf-8 -*-
"""
dump2png.py —— 把 frame_dump24 转储的 hex 帧转成 PNG（肉眼看仿真画面）

用法：
  python tools/dump2png.py <仿真目录> <前缀> <帧号...> [--out <输出目录>]
示例：
  python tools/dump2png.py 08_board_bringup/01_video_passthrough/sim passthrough_frame 0 1 2 3

输入：<仿真目录>/<前缀>_<帧号>.hex（frame_dump24 格式：每行 6 位十六进制 RGB888）
输出：<输出目录>/<前缀>_<帧号>.png（默认输出到 <仿真目录>/../sim_out 或 sim_out）
"""
import os
import sys

import numpy as np
from PIL import Image

W, H = 1280, 720


def load_frame(path):
    px = np.zeros((H, W, 3), dtype=np.uint8)
    i = 0
    n_x = 0
    with open(path, 'r') as f:
        for line in f:
            line = line.strip()
            if len(line) < 6:
                continue
            try:
                v = int(line[:6], 16)
            except ValueError:
                n_x += 1          # 'xxxxxx' = X 传播（未初始化 RAM 读出）
                i += 1
                continue
            y, x = divmod(i, W)
            if y < H:
                px[y, x] = ((v >> 16) & 0xFF, (v >> 8) & 0xFF, v & 0xFF)
            i += 1
    if i != W * H:
        print(f'  提示: {path} 像素数 {i} != {W*H}（帧不完整也会尽量出图）')
    if n_x:
        print(f'  提示: {path} 含 {n_x} 个 X 像素（按黑画，多为未初始化 RAM 的帧首现象）')
    return px


def main():
    argv = sys.argv[1:]
    outdir = 'sim_out'
    if '--out' in argv:
        outdir = argv[argv.index('--out') + 1]
        argv = argv[:argv.index('--out')] + argv[argv.index('--out') + 2:]
    if len(argv) < 1:
        sys.exit(__doc__)
    simdir = argv[0]
    prefix = argv[1] if len(argv) > 1 else 'passthrough_frame'
    frames = [int(a) for a in argv[2:]] or [0]
    os.makedirs(outdir, exist_ok=True)
    for n in frames:
        p = os.path.join(simdir, f'{prefix}_{n}.hex')
        if not os.path.exists(p):
            print(f'跳过（不存在）: {p}')
            continue
        img = load_frame(p)
        out = os.path.join(outdir, f'{prefix}_{n}.png')
        Image.fromarray(img).save(out)
        mean = img.mean(axis=(0, 1))
        print(f'{out}  亮度均值 R={mean[0]:.1f} G={mean[1]:.1f} B={mean[2]:.1f}')


if __name__ == '__main__':
    main()
\end{Verbatim}

\section{官方复用与生成文件说明}

\begin{itemize}
  \item \textbf{官方复用（不改动）}：edge\_detect\_96.v（Lab2 Sobel 边缘核，149 行，整链仿真已验证接线）；
        官方 ISP 明码副本（isp\_top/data128\_96/awb/demosaic 树/raw10 解包等，位于 04\_embedded/03\_project/sim/，仅用于仿真）；
        SC520 适配版 Lab1 完整 TD 工程（08\_board\_bringup，含加密 MIPI/HDMI 核，不可改）。
  \item \textbf{生成文件（禁止手改）}：hsv\_inv\_lut.v、font\_rom\_8x16.v、marquee\_rom\_64.v，
        均由 tools/gen\_vectors.py 从算法/字库源生成（节选见上文）。
  \item \textbf{仿真专用（禁止进 TD 工程）}：\*\_stub.v / \*\_sim.v、全部 tb\_\*.v、px\_serializer\_fifo.v。
\end{itemize}

\section{开放项与上板检查单}

\begin{enumerate}
  \item \textbf{（已修复，2026-10-03 修复轮 \#N+7 / 加固 \#N+9）ISP 行级突发 vs 视觉抽头}：
        实测坐实"AWB 输出为行级突发（每行恰 320 beat，行间消隐；2160 段=3 帧×720 行）"后：
        ① ser96to24 v4/v4.1——容量 128→\textbf{2048 像素} + \textbf{寄存器读}（+1 拍，
        兼容 LUTRAM/ERAM）+ \textbf{帧同步看门狗与 tlast 门控恢复}（tuser 永久缺失不再死锁，
        故障注入 TB 验证）+ O\_occ\_max 观测（经 status[6] 引出 sync\_err）；
        ② 抽头改 AWB 96bit 直抽。验收：\code{pixels=2764800, sof=3, eol=2160, drop=0}
        （修复前 188 像素/0 SOF），\textbf{FIFO 最大占用实测 1601/2048（78\%）}。
        遗留：真机上板用 O\_drop\_cnt 与帧计数双重确认；多行合并突发的最大行数需真机确认
        （3 行=2880px 会超 2048；不盲目升 4096——双实例 16 块 ERAM 超出整机预算）。
  \item \textbf{（已修复，\#N+7）128bit 重排路径 tuser 结构性错位}：改 AWB 直抽后该路径不再
        参与视觉链（isp\_top 新增 tap\_* 只增端口）。
  \item \textbf{（已修复，\#N+7）system\_top 握手}：\code{I\_tvalid} 直用 RAW 有效（与官方
        design\_top\_wrapper 同款）；\code{O\_tready} 显式 1'b1。注：官方 wrapper 亦悬空该网
        （有效行为=无背压），本次显式化消除 Z 传播。
  \item \textbf{ERAM 写模式 + 预算}：motion\_detect v3.1 "写延迟一拍"结构解耦成立；
        TD gate 报告确认 prev\_mem → ERAM（57 块）。v4 串化器 4 块/实例；跑马灯 ROM 1 块。
        \textbf{预算}：视觉链 \textbf{62/108（57.4\%）} + 官方 DDR/HDMI 链 39 块 ⇒ 真机整合
        ≈94\%——若吃紧：边缘支路 ser 降 DEPTH 512（1 块）或走共享串化器重构；接近冻结期
        不建议大改；最终以 P\&R 布局报告为准。
  \item \textbf{LUTRAM 映射（已解除）}：bin\_smooth3x3 行缓存推断为 Distributed RAM
        （gate 报告 256 dram lut）；ser96to24 v4 因注册读设计对 LUTRAM/ERAM 双映射皆正确。
  \item \textbf{gimbal\_ctrl 帧拍组合路径（已解除）}：pwm\_clk(50MHz) WNS \textbf{+8.747ns}。
        像素域曾两度指认 OSD 跑马灯 ROM 路径：先 case→数组（\#N+7）——但 \#N+9 复综合发现
        该改动\textbf{并未真正消除 mux 树}（TD 仍映射为 ~10 级 LUT3 链，WNS 一度回退到
        −1.025ns，"+0.541"是综合器估算运气）；最终\textbf{结构性流水化}（地址+ROM 读前瞻
        4 像素 + 4 级移位对齐）后 \textbf{WNS +2.441ns、0 违规端点}。教训：换写法≠映射变浅，
        必须看 gate 报告的逻辑级数。gate 级净延迟 ~68\% 为估算，P\&R 后复核。
  \item \textbf{I2C BFM ACK 位相}（tb\_i2c\_config）：列 TD 阶段项，需真机逻辑分析仪对照；
        寄存器序列（116 条）本身已由 ROM 逐条确认。
  \item \textbf{厂商 ISP 仿真环境限制}：无背压满速馈送下 AWB 行数完整性在行突发模式下
        实测零丢行（早期 668/720 观察与馈送模型有关）；上板仍以 O\_drop\_cnt 实测为准。
  \item \textbf{CDC 复位同源（已固化）}：cdc\_4phase 头部契约 + \code{tb\_cdc\_reset\_skew}
        （4 轮错开释放 PASS）；上板确认两复位同源，禁止运行时单侧复位。
  \item \textbf{算法切换帧专项 TB（已完成）}：输出侧 SOF 提交 + 提交沿组合选择 + 超时回退，
        经 \code{tb\_algo\_switch}（第 18 项）覆盖，全绿。
\end{enumerate}

\vspace{6pt}
{\small\color{black!60} —— 文档由 tools/make\_review\_pdf.py 自动生成；
代码以仓库实际文件为唯一真值，本文档中的说明如与代码冲突，以代码为准。}

\end{document}